% Corrigés des exercices du cours « Méthodes numériques pour physiciens et
% ingénieurs » (P. Woafo), rédigés à partir du cours, de Burden & Faires,
% Numerical Analysis (10e éd.), et de sources publiques citées.
\documentclass[11pt,a4paper,oneside]{report}

\usepackage[T1]{fontenc}
\usepackage[utf8]{inputenc}
\usepackage[french]{babel}
\usepackage{lmodern}
\usepackage{microtype}
\usepackage[a4paper,margin=2.5cm]{geometry}
\usepackage{amsmath,amssymb,mathtools}
\usepackage{array,booktabs}
\usepackage{enumitem}
\usepackage{titlesec}
\usepackage{needspace}
\usepackage{algpseudocode}
\usepackage{listings}
\usepackage{tikz}
\usetikzlibrary{arrows.meta,patterns}
\usepackage{xurl}
\usepackage[hidelinks]{hyperref}

% ---------------------------------------------------------------
% Numérotation et titres
% ---------------------------------------------------------------
\renewcommand{\thechapter}{\Roman{chapter}}
\setcounter{secnumdepth}{-1}
\setcounter{tocdepth}{1}
\titleformat{\chapter}[display]
  {\normalfont\bfseries\centering}{}{0pt}{\LARGE}
\titlespacing*{\chapter}{0pt}{-20pt}{24pt}
\titleformat{\section}{\normalfont\large\bfseries\needspace{8\baselineskip}}{}{0pt}{}
\titleformat{\subsection}{\normalfont\normalsize\bfseries}{}{0pt}{}

% ---------------------------------------------------------------
% Structure de chaque corrigé : pourquoi, puis comment, puis vérification
% ---------------------------------------------------------------
\newcommand{\Enonce}[1]{\par\noindent\textit{Énoncé.}~\textit{#1}\par}
\newcommand{\Pourquoi}{\par\medskip\needspace{4\baselineskip}\noindent\textbf{Pourquoi ?}\quad}
\newcommand{\Etapes}{\par\medskip\needspace{4\baselineskip}\noindent\textbf{Comment ? (étapes)}\par}
\newcommand{\Etape}[1]{\par\smallskip\noindent\textit{Étape #1.}\quad}
\newcommand{\Verif}{\par\medskip\needspace{3\baselineskip}\noindent\textbf{Vérification numérique.}\quad}
\newcommand{\Refs}{\par\medskip\noindent\textbf{Références.}\quad}
\newcommand{\Attention}{\par\medskip\noindent\textbf{Remarque sur l'énoncé.}\quad}
\newcommand{\Question}[1]{\par\bigskip\needspace{8\baselineskip}\noindent\textbf{#1}\quad}

% ---------------------------------------------------------------
% Algorithmes en français (sans cadre)
% ---------------------------------------------------------------
\algrenewcommand\algorithmicrequire{\textbf{Données :}}
\algrenewcommand\algorithmicensure{\textbf{Résultat :}}
\algrenewcommand\algorithmicfor{\textbf{Pour}}
\algrenewcommand\algorithmicdo{\textbf{faire}}
\algrenewcommand\algorithmicwhile{\textbf{Tant que}}
\algrenewcommand\algorithmicif{\textbf{Si}}
\algrenewcommand\algorithmicthen{\textbf{alors}}
\algrenewcommand\algorithmicelse{\textbf{Sinon}}
\algrenewcommand\algorithmicrepeat{\textbf{Répéter}}
\algrenewcommand\algorithmicuntil{\textbf{Jusqu'à}}
\algrenewcommand\algorithmicreturn{\textbf{Retourner}}
\algrenewcommand\algorithmicend{\textbf{Fin}}
\algrenewtext{EndFor}{\textbf{Fin pour}}
\algrenewtext{EndWhile}{\textbf{Fin tant que}}
\algrenewtext{EndIf}{\textbf{Fin si}}
\algnewcommand\algorithmicprint{\textbf{Écrire}}
\algnewcommand\Print{\State\algorithmicprint{} }
\algnewcommand\algorithmicread{\textbf{Lire}}
\algnewcommand\Read{\State\algorithmicread{} }
\algnewcommand{\Stop}{\State\textbf{Arrêt}}
\algrenewcommand\algorithmiccomment[1]{\hfill\textit{\{#1\}}}
\newenvironment{algo}[1][]
  {\par\smallskip\needspace{5\baselineskip}\noindent\textbf{Algorithme#1}\par\nopagebreak
   \begin{algorithmic}[0]}
  {\end{algorithmic}\par\smallskip}

% ---------------------------------------------------------------
% Programmes
% ---------------------------------------------------------------
\lstset{
  basicstyle=\ttfamily\small,
  columns=fullflexible,
  keepspaces=true,
  frame=lines,
  framesep=4pt,
  xleftmargin=1em,
  breaklines=true,
  showstringspaces=false,
  inputencoding=utf8,
  extendedchars=true,
  literate={é}{{\'e}}1 {è}{{\`e}}1 {à}{{\`a}}1 {ê}{{\^e}}1 {ç}{{\c c}}1
           {É}{{\'E}}1 {ô}{{\^o}}1 {î}{{\^i}}1 {ù}{{\`u}}1
}
\newcommand{\programme}[1]{%
  \par\smallskip\noindent\textit{Programme Python} (fichier \fichier{programmes/#1}) :
  \lstinputlisting[language=Python]{programmes/#1}}

% ---------------------------------------------------------------
% Macros mathématiques
% ---------------------------------------------------------------
\newcommand{\dd}{\mathrm{d}}
\newcommand{\R}{\mathbb{R}}
\newcommand{\abs}[1]{\left|#1\right|}
\newcommand{\BF}[1]{[BF, #1]}
\newcommand{\fichier}[1]{\nolinkurl{#1}}

\setlength{\parskip}{0.35em}
\setlength{\parindent}{0pt}
\allowdisplaybreaks

\hypersetup{
  pdftitle={Corrigés des exercices -- Méthodes numériques pour physiciens et ingénieurs},
  pdfsubject={Méthodes numériques},
  pdflang={fr},
  hypertexnames=false
}

\begin{document}

\begin{titlepage}
  \centering
  \vspace*{3cm}
  {\Huge\bfseries Corrigés des exercices\par}
  \vspace{1cm}
  {\Large Méthodes numériques pour physiciens et ingénieurs\par}
  \vspace{0.3cm}
  {\large cours de P. Woafo (Université de Yaoundé I)\par}
  \vfill
  {\normalsize Corrigés rédigés à partir du cours et de\par
   R. L. Burden, J. D. Faires, A. M. Burden, \textit{Numerical Analysis},\par
   10\textsuperscript{e} édition, Cengage Learning, 2016.\par}
  \vspace{2cm}
\end{titlepage}

\chapter*{Avant-propos}
\addcontentsline{toc}{chapter}{Avant-propos}

Ce document rassemble les corrigés de tous les exercices du cours
\og Méthodes numériques pour physiciens et ingénieurs \fg{} : les exercices
placés dans les chapitres I à VI, puis les dix-neuf exercices de la fin du
document.

\section{Comment chaque corrigé est construit}

Chaque corrigé suit le même plan, dans cet ordre :
\begin{enumerate}[label=\arabic*.]
  \item \textbf{Pourquoi ?} On explique d'abord l'idée : quel est le
        problème, pourquoi la méthode choisie convient, quelle difficulté
        elle résout et à quelles conditions elle marche.
  \item \textbf{Comment ? (étapes)} On déroule ensuite la solution étape par
        étape : mise en équations, calculs intermédiaires, puis algorithme et,
        quand l'énoncé le demande, programme.
  \item \textbf{Vérification numérique.} Chaque résultat chiffré et chaque
        algorithme ont été réellement calculés ou exécutés. Les valeurs
        données sont celles que ces calculs ont produites.
  \item \textbf{Références.} On indique où trouver la méthode dans le livre de
        référence ou dans une source publique.
\end{enumerate}
Lorsqu'un énoncé est incomplet, ambigu ou semble contenir une erreur, on le
dit dans un paragraphe \og Remarque sur l'énoncé \fg{}, et on précise
l'interprétation retenue. Rien n'est supposé sans être signalé.

\section{Références utilisées}

La référence principale est le livre
R. L. Burden, J. D. Faires et A. M. Burden, \textit{Numerical Analysis},
10\textsuperscript{e} édition, Cengage Learning, 2016, noté \textbf{[BF]}.
Les renvois du type \BF{alg.~6.1, p.~368} donnent le numéro de l'algorithme ou
du théorème et la page imprimée du livre. Les autres sources (articles,
encyclopédies en ligne) sont listées dans la bibliographie à la fin du
document.

\section{Programmes et calculs}

Les programmes demandés par les énoncés se trouvent dans le dossier
\fichier{latex-corriges/programmes/} (langages C, Fortran 90 et Pascal). Ils ont
été compilés avec \texttt{gcc}, \texttt{gfortran} et Free Pascal
(\texttt{fpc}) ; le script \fichier{tester.sh} les compile et les exécute
tous avec les données utilisées dans ce document (résultats dans
\fichier{calculs/programmes_resultats.txt}). Les calculs numériques et les
tests d'algorithmes sont dans \fichier{latex-corriges/calculs/} (Python avec
NumPy et SciPy) : \fichier{calculs.py}, \fichier{complements.py} et
\fichier{tests_algorithmes.py}, avec leurs sorties dans les fichiers
\fichier{*_resultats.txt} et \fichier{resultats.txt}.

Dans les algorithmes, la flèche $\leftarrow$ désigne l'affectation, et les
notations sont celles du cours.


\tableofcontents

\chapter{Chapitre I : systèmes d'équations algébriques linéaires}

% ------------------------------------------------------------------
\section{Exercice I.1 : résolution du système triangulaire}

\Enonce{Écrire l'algorithme de résolution du système algébrique triangulaire.}

\Pourquoi Après la triangularisation de Gauss, la dernière équation ne
contient plus qu'une seule inconnue, $x_n$. L'avant-dernière contient $x_{n-1}$
et $x_n$, et ainsi de suite : l'équation $k$ ne fait intervenir que
$x_k, x_{k+1}, \dots, x_n$. Si l'on résout les équations \emph{de bas en
haut}, chaque équation n'a donc plus qu'une seule inconnue nouvelle, $x_k$, car
$x_{k+1}, \dots, x_n$ ont déjà été calculés. C'est la \og remontée \fg{}
(substitution arrière). Elle exige $a_{kk} \neq 0$ pour tout $k$, ce que la
triangularisation avec pivot garantit lorsque la matrice est régulière.

\Etapes
\Etape{1} La dernière équation $a_{nn}x_n = b_n$ donne directement
$x_n = b_n/a_{nn}$.

\Etape{2} Pour $k = n-1$, puis $n-2$, \dots, jusqu'à $1$ (ordre décroissant,
c'est indispensable), l'équation $k$ s'écrit
\[
  a_{kk}x_k + \sum_{l=k+1}^{n} a_{kl}x_l = b_k ,
\]
où toutes les $x_l$ de la somme sont déjà connues. On calcule d'abord la somme
$s = b_k - \sum_{l=k+1}^{n} a_{kl}x_l$, puis $x_k = s/a_{kk}$. C'est
exactement la forme condensée du cours.

\Etape{3} Coût : l'étape $k$ demande $n-k$ multiplications et une division, soit
au total $1 + \sum_{k=1}^{n-1}(n-k+1) = (n^2+n)/2$ multiplications ou divisions
\BF{p.~371}. C'est négligeable devant le coût de la triangularisation (de
l'ordre de $n^3/3$).

\begin{algo}[ : remontée]
  \Require $n$, la matrice triangulaire supérieure $(a_{kl})$, le second membre $(b_k)$
  \State $x_n \leftarrow b_n / a_{nn}$
  \For{$k$ allant de $n-1$ à $1$ (pas $-1$)}
    \State $s \leftarrow b_k$
    \For{$l$ allant de $k+1$ à $n$}
      \State $s \leftarrow s - a_{kl}\,x_l$
    \EndFor
    \State $x_k \leftarrow s / a_{kk}$
  \EndFor
  \Print $x_1, \dots, x_n$
\end{algo}

\Verif Cet algorithme est la fonction \texttt{remontee} du programme
\fichier{gauss.py} (exercice 10 de la fin du document). Sur un système
$6\times 6$ aléatoire dont on connaît la solution, l'erreur maximale obtenue est
$2{,}4\times 10^{-16}$.

\Refs \BF{alg.~6.1, étapes 8 et 9, p.~368 ; décompte des opérations
p.~370--371}.

% ------------------------------------------------------------------
\section{Exercice I.2 : échange de lignes (pivot maximal)}

\Enonce{Écrire l'algorithme qui permet d'interchanger les lignes des systèmes
algébriques dû à la règle de pivot maximal.}

\Pourquoi À l'étape $k$ de la triangularisation, on divise par $a_{kk}$
(multiplicateur $a_{lk}/a_{kk}$). Deux problèmes peuvent survenir :
\begin{itemize}
  \item si $a_{kk} = 0$, la méthode s'arrête (division par zéro) alors que le
        système peut très bien avoir une solution ;
  \item si $\abs{a_{kk}}$ est petit devant les autres coefficients de la
        colonne, les multiplicateurs sont grands et amplifient les erreurs
        d'arrondi. Burden et Faires montrent sur le système
        $0{,}003000\,x_1 + 59{,}14\,x_2 = 59{,}17$,
        $5{,}291\,x_1 - 6{,}130\,x_2 = 46{,}78$ (solution exacte $x_1 = 10$,
        $x_2 = 1$) qu'en arithmétique à quatre chiffres on obtient
        $x_1 \approx -10{,}00$ sans échange de lignes \BF{ex.~1, p.~376}.
\end{itemize}
La règle du pivot maximal (pivot partiel) choisit comme équation pivotale,
parmi les lignes $k$ à $n$, celle dont le coefficient de $x_k$ est le plus
grand en valeur absolue. Les multiplicateurs vérifient alors
$\abs{a_{lk}/a_{kk}} \le 1$ : les erreurs ne sont plus amplifiées. Échanger deux
équations ne change pas la solution du système.

\Etapes
\Etape{1} À l'étape $k$, chercher l'indice $p \ge k$ tel que
$\abs{a_{pk}} = \max_{k \le i \le n}\abs{a_{ik}}$.

\Etape{2} Si $a_{pk} = 0$, toute la colonne sous la diagonale est nulle : la
matrice est singulière, on arrête.

\Etape{3} Si $p \neq k$, échanger les lignes $k$ et $p$ de la matrice (il suffit
de le faire pour les colonnes $k$ à $n$, les colonnes précédentes étant déjà
nulles) \emph{et} les composantes $b_k$ et $b_p$ du second membre. L'échange
utilise une variable auxiliaire $t$.

\begin{algo}[ : recherche du pivot et échange (étape $k$)]
  \State $p \leftarrow k$
  \For{$i$ allant de $k+1$ à $n$}
    \If{$\abs{a_{ik}} > \abs{a_{pk}}$} \State $p \leftarrow i$ \EndIf
  \EndFor
  \If{$a_{pk} = 0$} \Print \og matrice singulière \fg{} \Stop \EndIf
  \If{$p \neq k$}
    \For{$j$ allant de $k$ à $n$}
      \State $t \leftarrow a_{kj}$ ; $a_{kj} \leftarrow a_{pj}$ ; $a_{pj} \leftarrow t$
    \EndFor
    \State $t \leftarrow b_k$ ; $b_k \leftarrow b_p$ ; $b_p \leftarrow t$
  \EndIf
\end{algo}

\Refs \BF{§~6.2, pivot partiel, alg.~6.2, p.~378--379}.

% ------------------------------------------------------------------
\section{Exercice I.3 : méthode de Gauss complète avec pivot maximal}

\Enonce{Écrire l'algorithme complet de la méthode de Gauss (triangularisation
et résolution du système triangulaire) en prenant soin de faire jouer la règle
du pivot maximal.}

\Pourquoi On assemble les deux morceaux précédents : pour chaque étape $k$, on
choisit d'abord le pivot (exercice I.2), puis on élimine $x_k$ des lignes
$k+1$ à $n$ avec l'algorithme de l'étape $k$ du cours ; une fois le système
triangulaire obtenu, on fait la remontée (exercice I.1). Le pivot doit être
choisi \emph{à chaque étape}, car les coefficients changent d'une étape à
l'autre.

\Etapes
\Etape{1} Pour $k = 1, \dots, n-1$ : choix du pivot et échange éventuel.

\Etape{2} Toujours à l'étape $k$, pour chaque ligne $l > k$ : calcul du
multiplicateur $m = a_{lk}/a_{kk}$ (calculé une seule fois, ce qui évite de
refaire la division pour chaque colonne), puis $b_l \leftarrow b_l - m\,b_k$ et
$a_{lj} \leftarrow a_{lj} - m\,a_{kj}$ pour $j = k+1, \dots, n$.

\Etape{3} Après la boucle, vérifier $a_{nn} \neq 0$, puis remontée.

\begin{algo}[ : Gauss avec pivot maximal]
  \Require $n$, $A = (a_{ij})$, $b = (b_i)$
  \For{$k$ allant de $1$ à $n-1$}
    \State recherche du pivot $p$ et échange des lignes $k$ et $p$ \Comment{exercice I.2}
    \For{$l$ allant de $k+1$ à $n$}
      \State $m \leftarrow a_{lk}/a_{kk}$
      \State $b_l \leftarrow b_l - m\,b_k$
      \For{$j$ allant de $k+1$ à $n$}
        \State $a_{lj} \leftarrow a_{lj} - m\,a_{kj}$
      \EndFor
      \State $a_{lk} \leftarrow 0$
    \EndFor
  \EndFor
  \If{$a_{nn} = 0$} \Print \og matrice singulière \fg{} \Stop \EndIf
  \State remontée \Comment{exercice I.1}
\end{algo}

Le coût total est de $n^3/3 + n^2 - n/3$ multiplications ou divisions
\BF{p.~371} ; la recherche du pivot n'ajoute que des comparaisons.

\Verif Exemple traité à la main (et par le programme \fichier{gauss.py}) :
le système
\[
  2x_2 + x_3 = 5, \qquad x_1 + x_2 + x_3 = 6, \qquad 2x_1 + x_2 + 3x_3 = 13
\]
ne peut pas être traité sans échange car $a_{11} = 0$.
\begin{itemize}
  \item $k = 1$ : le plus grand $\abs{a_{i1}}$ est $2$ (ligne 3). On échange
        les lignes 1 et 3 : $(2\ 1\ 3 \mid 13)$, $(1\ 1\ 1 \mid 6)$,
        $(0\ 2\ 1 \mid 5)$. Multiplicateurs $1/2$ et $0$ : la ligne 2 devient
        $(0\ \tfrac12\ {-\tfrac12} \mid -\tfrac12)$.
  \item $k = 2$ : entre $\tfrac12$ (ligne 2) et $2$ (ligne 3), le pivot est
        $2$ ; on échange les lignes 2 et 3. Multiplicateur
        $\tfrac12/2 = \tfrac14$ : la ligne 3 devient
        $(0\ 0\ {-\tfrac34} \mid -\tfrac74)$.
  \item Remontée : $x_3 = 7/3$, $x_2 = (5 - 7/3)/2 = 4/3$,
        $x_1 = (13 - 4/3 - 3\cdot 7/3)/2 = 7/3$.
\end{itemize}
Le programme affiche bien $x_1 = 2{,}3333333333$, $x_2 = 1{,}3333333333$,
$x_3 = 2{,}3333333333$. Un test sur une matrice $5\times 5$ aléatoire avec
$a_{11} = 0$ donne une erreur de $10^{-15}$ (\fichier{tests_algorithmes.py}).

\Refs \BF{alg.~6.1, p.~368, et alg.~6.2, p.~378--379}.

% ------------------------------------------------------------------
\section{\texorpdfstring{Exercice I.4 : calcul des $l_{ij}$ et des $s_{ij}$}{Exercice I.4 : calcul des lij et des sij}}

\Enonce{Écrire l'algorithme permettant de calculer les $l_{ij}$ et $s_{ij}$.}

\Pourquoi Une fois $A = LS$ connue, résoudre $AX = b$ revient à résoudre deux
systèmes triangulaires ($LY = b$, puis $SX = Y$), soit environ $n^2$
opérations, au lieu de refaire toute l'élimination (environ $n^3/3$). C'est
très avantageux quand on doit résoudre plusieurs systèmes de même matrice $A$
et de seconds membres différents. Il faut organiser le calcul dans un ordre
où chaque élément ne dépend que d'éléments déjà calculés : c'est pourquoi le
cours calcule alternativement une colonne de $L$ puis une ligne de $S$.

\Etapes
\Etape{1} Écrire le produit $LS$ terme à terme. Comme $l_{jk} = 0$ pour $k > j$,
$s_{kp} = 0$ pour $k > p$ et $s_{pp} = 1$ :
\[
  a_{jp} = \sum_{k=1}^{p-1} l_{jk}s_{kp} + l_{jp} \quad (j \ge p), \qquad
  a_{pj} = \sum_{k=1}^{p-1} l_{pk}s_{kj} + l_{pp}s_{pj} \quad (j > p).
\]

\Etape{2} En isolant l'élément inconnu, on obtient les formules générales qui
contiennent toutes celles du cours (pour $p = 1$, les sommes sont vides) :
\[
  l_{jp} = a_{jp} - \sum_{k=1}^{p-1} l_{jk}s_{kp} \quad (j = p, \dots, n), \qquad
  s_{pj} = \frac{1}{l_{pp}}\Bigl(a_{pj} - \sum_{k=1}^{p-1} l_{pk}s_{kj}\Bigr)
  \quad (j = p+1, \dots, n).
\]

\Etape{3} Ordre du calcul : pour $p = 1, 2, \dots, n$, d'abord la colonne $p$ de
$L$ (elle utilise les colonnes $1$ à $p-1$ de $L$ et les lignes $1$ à $p-1$ de
$S$, déjà connues), puis la ligne $p$ de $S$ (elle utilise en plus $l_{pp}$, que
l'on vient de calculer). Il faut $l_{pp} \neq 0$ ; sinon la décomposition sans
échange de lignes est impossible.

\begin{algo}[ : décomposition $A = LS$ (Crout)]
  \Require $n$, $A = (a_{ij})$
  \For{$p$ allant de $1$ à $n$}
    \For{$j$ allant de $p$ à $n$} \Comment{colonne $p$ de $L$}
      \State $l_{jp} \leftarrow a_{jp}$
      \For{$k$ allant de $1$ à $p-1$}
        \State $l_{jp} \leftarrow l_{jp} - l_{jk}\,s_{kp}$
      \EndFor
    \EndFor
    \If{$l_{pp} = 0$} \Print \og décomposition impossible \fg{} \Stop \EndIf
    \State $s_{pp} \leftarrow 1$
    \For{$j$ allant de $p+1$ à $n$} \Comment{ligne $p$ de $S$}
      \State $s_{pj} \leftarrow a_{pj}$
      \For{$k$ allant de $1$ à $p-1$}
        \State $s_{pj} \leftarrow s_{pj} - l_{pk}\,s_{kj}$
      \EndFor
      \State $s_{pj} \leftarrow s_{pj}/l_{pp}$
    \EndFor
  \EndFor
\end{algo}

\Verif Sur une matrice $5\times 5$ aléatoire à diagonale dominante, on
obtient $\max\abs{LS - A} = 9\times 10^{-16}$ (programme
\fichier{tests_algorithmes.py}).

\Refs \BF{alg.~6.4, p.~410} avec le choix $u_{ii} = 1$ : c'est la
décomposition de Crout \BF{p.~410}.

% ------------------------------------------------------------------
\section{\texorpdfstring{Exercice I.5 : inversion d'une matrice et produit $L^{-1}b$}{Exercice I.5 : inversion d'une matrice et produit L\textasciicircum -1 b}}

\Enonce{Trouver dans la littérature l'algorithme d'inversion d'une matrice et de
calcul des éléments du produit $L^{-1}b$.}

\Pourquoi Le système $SX = L^{-1}b$ du cours demande le vecteur $Y = L^{-1}b$.
On ne calcule jamais la matrice $L^{-1}$ elle-même : $Y$ est la solution de
$LY = b$, système triangulaire \emph{inférieur} que l'on résout de haut en bas
(\og descente \fg{}), comme la remontée mais dans l'autre sens. C'est moins
coûteux et plus précis que de former $L^{-1}$. De même, Burden et Faires
rappellent qu'il n'est pas efficace de calculer $A^{-1}$ pour résoudre
$AX = b$ \BF{p.~392} ; l'inverse n'est utile que si on en a besoin pour
elle-même.

\Etapes
\Etape{1} \emph{Produit $Y = L^{-1}b$.} La première équation de $LY = b$ donne
$y_1 = b_1/l_{11}$ ; la ligne $i$ ne contient que $y_1, \dots, y_i$, d'où
\[
  y_i = \frac{1}{l_{ii}}\Bigl(b_i - \sum_{j=1}^{i-1} l_{ij}\,y_j\Bigr),
  \qquad i = 2, \dots, n \quad \BF{p.~411}.
\]

\Etape{2} \emph{Inverse.} Si $B = A^{-1}$, alors $AB = I$ : la colonne $j$ de
$B$, notée $B_j$, est la solution de $AB_j = e_j$, où $e_j$ est la colonne $j$
de la matrice identité \BF{p.~392--393}. On résout donc $n$ systèmes de même
matrice $A$. Avec la décomposition $A = LS$ faite une seule fois, chaque
colonne coûte une descente ($LY = e_j$) et une remontée ($SB_j = Y$).
Burden et Faires procèdent de façon équivalente par élimination de Gauss sur la
matrice augmentée $[A \mid I]$ \BF{p.~393} ; la méthode de Jordan du cours
transforme directement $[A \mid I]$ en $[I \mid A^{-1}]$.

\begin{algo}[ : descente $Y = L^{-1}b$]
  \State $y_1 \leftarrow b_1/l_{11}$
  \For{$i$ allant de $2$ à $n$}
    \State $s \leftarrow b_i$
    \For{$j$ allant de $1$ à $i-1$} \State $s \leftarrow s - l_{ij}\,y_j$ \EndFor
    \State $y_i \leftarrow s/l_{ii}$
  \EndFor
\end{algo}

\begin{algo}[ : inverse de $A$]
  \State décomposer $A = LS$ \Comment{exercice I.4}
  \For{$j$ allant de $1$ à $n$}
    \State $e \leftarrow$ colonne $j$ de l'identité ($e_i = 0$ pour $i \ne j$, $e_j = 1$)
    \State résoudre $LY = e$ par descente
    \State résoudre $SX = Y$ par remontée \Comment{exercice I.1}
    \For{$i$ allant de $1$ à $n$} \State $b_{ij} \leftarrow x_i$ \EndFor
  \EndFor
  \Print la matrice $B = A^{-1}$
\end{algo}

\Verif Sur la même matrice $5\times 5$ : erreur de résolution de $LSX = b$
égale à $6\times 10^{-17}$, et $\max\abs{AA^{-1} - I} = 2\times 10^{-16}$.

\Refs \BF{§~6.3, p.~391--393 (inverse) ; §~6.5, p.~411 (descente et
remontée)}.

% ------------------------------------------------------------------
\section{\texorpdfstring{Exercice I.6 : coefficients $m_{pp}$ et $m_{pj}$ de Cholesky}{Exercice I.6 : coefficients mpp et mpj de Cholesky}}

\Enonce{Écrire l'algorithme de calcul des coefficients $m_{pp}$ et $m_{pj}$.}

\Pourquoi Pour une matrice symétrique, $A = M^{t}M$ avec une seule matrice
triangulaire supérieure $M$ à calculer : on stocke moitié moins de nombres et
on fait environ deux fois moins d'opérations que pour $LS$. Le calcul des
$m_{pp}$ demande des racines carrées : il faut que les quantités sous la racine
soient strictement positives. C'est le cas si et seulement si $A$ est
\emph{définie positive} (et pas seulement symétrique) \BF{§~6.6,
p.~423--424}. L'algorithme doit donc tester ce signe.

\Etapes
\Etape{1} $M$ étant triangulaire supérieure ($m_{kp} = 0$ si $k > p$), le
coefficient $(p, j)$ de $M^{t}M$ pour $p \le j$ vaut
\[
  a_{pj} = \sum_{k=1}^{p} m_{kp}\,m_{kj}
         = \sum_{k=1}^{p-1} m_{kp}\,m_{kj} + m_{pp}\,m_{pj} .
\]

\Etape{2} Pour $j = p$, on obtient
$a_{pp} = \sum_{k<p} m_{kp}^2 + m_{pp}^2$, d'où
$m_{pp} = \sqrt{a_{pp} - \sum_{k=1}^{p-1} m_{kp}^2}$. Pour $j > p$, on obtient
$m_{pj} = \bigl(a_{pj} - \sum_{k=1}^{p-1} m_{kp}m_{kj}\bigr)/m_{pp}$. Ce sont
les formules du cours (avec $m_{11} = \sqrt{a_{11}}$ et
$m_{1j} = a_{1j}/m_{11}$ pour $p = 1$).

\Etape{3} On calcule $M$ ligne par ligne : la ligne $p$ n'utilise que les lignes
$1$ à $p-1$, déjà calculées.

\begin{algo}[ : Cholesky $A = M^{t}M$]
  \For{$p$ allant de $1$ à $n$}
    \State $s \leftarrow a_{pp}$
    \For{$k$ allant de $1$ à $p-1$} \State $s \leftarrow s - m_{kp}^2$ \EndFor
    \If{$s \le 0$} \Print \og $A$ n'est pas définie positive \fg{} \Stop \EndIf
    \State $m_{pp} \leftarrow \sqrt{s}$
    \For{$j$ allant de $p+1$ à $n$}
      \State $s \leftarrow a_{pj}$
      \For{$k$ allant de $1$ à $p-1$} \State $s \leftarrow s - m_{kp}\,m_{kj}$ \EndFor
      \State $m_{pj} \leftarrow s/m_{pp}$
    \EndFor
  \EndFor
\end{algo}

\Verif Pour la matrice de \BF{ex.~4, p.~424},
\[
  A = \begin{pmatrix} 4 & -1 & 1\\ -1 & 4{,}25 & 2{,}75\\ 1 & 2{,}75 & 3{,}5 \end{pmatrix},
  \quad\text{l'algorithme donne}\quad
  M = \begin{pmatrix} 2 & -0{,}5 & 0{,}5\\ 0 & 2 & 1{,}5\\ 0 & 0 & 1 \end{pmatrix},
\]
qui est bien la transposée de la matrice $L$ trouvée par Burden et Faires
\BF{p.~424}, et $M^{t}M = A$ exactement.

\Refs \BF{alg.~6.6, p.~423--424} (forme $A = LL^{t}$, avec $L = M^{t}$).

% ------------------------------------------------------------------
\section{Exercice I.7 : Jacobi et Gauss-Seidel avec condition d'arrêt}

\Enonce{Écrire les algorithmes des méthodes de Jacobi et de Gauss-Seidel avec
imposition de la condition d'arrêt.}

\Pourquoi Une méthode itérative produit une suite infinie
$X^{(0)}, X^{(1)}, \dots$ ; il faut donc décider quand s'arrêter. Le cours
recommande le test relatif
$d = \sum_i\abs{x_i^{(k+1)} - x_i^{(k)}}/\sum_i\abs{x_i^{(k+1)}} < \varepsilon$,
qui ne dépend pas de l'ordre de grandeur des inconnues. Comme la convergence
n'est pas toujours assurée, on limite aussi le nombre d'itérations à $k_{\max}$
(une condition suffisante de convergence des deux méthodes est que $A$ soit à
diagonale strictement dominante \BF{th.~7.21, p.~464}). La différence entre
les deux méthodes impose une différence de programmation :
\begin{itemize}
  \item Jacobi n'utilise que les valeurs de l'itération précédente : il faut
        deux tableaux, l'ancien $x$ et le nouveau $y$ ;
  \item Gauss-Seidel utilise chaque nouvelle valeur dès qu'elle est calculée :
        un seul tableau suffit, mis à jour sur place ; il faut seulement garder
        l'ancienne valeur de $x_i$ pour le test d'arrêt.
\end{itemize}

\Etapes
\Etape{1} Valeurs initiales $x_i = 0$ (ou $1$), comme le propose le cours.

\Etape{2} À chaque itération, calculer pour $i = 1, \dots, n$ la nouvelle valeur
par la formule du cours, en accumulant le numérateur et le dénominateur de $d$.

\Etape{3} Arrêter si $d < \varepsilon$ (convergence) ou si $k = k_{\max}$
(échec, à signaler).

\begin{algo}[ : Jacobi]
  \Require $n$, $A$, $b$, $\varepsilon$, $k_{\max}$
  \For{$i$ allant de $1$ à $n$} \State $x_i \leftarrow 0$ \EndFor
  \State $k \leftarrow 0$
  \Repeat
    \State $k \leftarrow k+1$ ; $\mathit{num} \leftarrow 0$ ; $\mathit{den} \leftarrow 0$
    \For{$i$ allant de $1$ à $n$}
      \State $s \leftarrow b_i$
      \For{$j$ allant de $1$ à $n$, $j \neq i$} \State $s \leftarrow s - a_{ij}\,x_j$ \EndFor
      \State $y_i \leftarrow s/a_{ii}$
      \State $\mathit{num} \leftarrow \mathit{num} + \abs{y_i - x_i}$ ; $\mathit{den} \leftarrow \mathit{den} + \abs{y_i}$
    \EndFor
    \For{$i$ allant de $1$ à $n$} \State $x_i \leftarrow y_i$ \EndFor
  \Until{$\mathit{num}/\mathit{den} < \varepsilon$ \textbf{ou} $k \ge k_{\max}$}
  \If{$\mathit{num}/\mathit{den} \ge \varepsilon$} \Print \og pas de convergence \fg{} \EndIf
  \Print $k$, $x_1, \dots, x_n$
\end{algo}

\begin{algo}[ : Gauss-Seidel]
  \Require $n$, $A$, $b$, $\varepsilon$, $k_{\max}$
  \For{$i$ allant de $1$ à $n$} \State $x_i \leftarrow 0$ \EndFor
  \State $k \leftarrow 0$
  \Repeat
    \State $k \leftarrow k+1$ ; $\mathit{num} \leftarrow 0$ ; $\mathit{den} \leftarrow 0$
    \For{$i$ allant de $1$ à $n$}
      \State $s \leftarrow b_i$
      \For{$j$ allant de $1$ à $n$, $j \neq i$} \State $s \leftarrow s - a_{ij}\,x_j$
        \Comment{$x_j$ est déjà la nouvelle valeur si $j < i$}
      \EndFor
      \State $t \leftarrow s/a_{ii}$
      \State $\mathit{num} \leftarrow \mathit{num} + \abs{t - x_i}$ ; $\mathit{den} \leftarrow \mathit{den} + \abs{t}$
      \State $x_i \leftarrow t$
    \EndFor
  \Until{$\mathit{num}/\mathit{den} < \varepsilon$ \textbf{ou} $k \ge k_{\max}$}
  \Print $k$, $x_1, \dots, x_n$
\end{algo}

\Verif Sur le système de \BF{ex.~1, p.~456} ($10x_1 - x_2 + 2x_3 = 6$,
$-x_1 + 11x_2 - x_3 + 3x_4 = 25$, $2x_1 - x_2 + 10x_3 - x_4 = -11$,
$3x_2 - x_3 + 8x_4 = 15$, solution $(1, 2, -1, 1)$), avec
$\varepsilon = 10^{-3}$ : Jacobi s'arrête après 10 itérations sur
$(1{,}0001 ; 1{,}9998 ; -0{,}9998 ; 0{,}9998)$ et Gauss-Seidel après 5 itérations
sur $(1{,}0001 ; 2{,}0000 ; -1{,}0000 ; 1{,}0000)$. Ce sont les mêmes nombres
d'itérations que dans le livre (10 et 5), qui utilise un test d'arrêt voisin
fondé sur la norme maximale \BF{p.~458--461}.

\Refs \BF{alg.~7.1, p.~459 (Jacobi) ; alg.~7.2, p.~462 (Gauss-Seidel) ;
th.~7.21, p.~464}.

\chapter{Chapitre II : équations algébriques non linéaires}

% ------------------------------------------------------------------
\section{Exercice II.1 : Newton-Raphson avec condition d'arrêt}

\Enonce{Écrire l'algorithme de la méthode de Newton-Raphson en y intégrant la
condition d'arrêt.}

\Pourquoi La suite $x_{n+1} = x_n - f(x_n)/f'(x_n)$ ne s'arrête jamais
d'elle-même : il faut un critère qui garantisse la précision $E$ demandée. Le
cours donne la majoration $\abs{x_n - \bar{x}} \le \abs{f(x_n)}/m$ si
$\abs{f'} \ge m > 0$ sur un intervalle $[a,b]$ contenant $x_n$ et la racine
$\bar{x}$. Elle vient du théorème des accroissements finis :
$f(x_n) = f(x_n) - f(\bar{x}) = f'(\xi)(x_n - \bar{x})$ avec $\xi$ entre $x_n$
et $\bar{x}$, donc $\abs{x_n - \bar{x}} = \abs{f(x_n)}/\abs{f'(\xi)} \le
\abs{f(x_n)}/m$. Le test $\abs{f(x_n)}/m < E$ garantit donc
$\abs{x_n - \bar{x}} < E$. Deux sécurités sont nécessaires en plus : arrêter si
$f'(x_n) = 0$ (la tangente est horizontale, la formule n'a plus de sens) et
limiter le nombre d'itérations, car Newton peut diverger si $x_0$ est mal choisi
(remarque du cours).

\Etapes
\Etape{1} Localiser la racine dans $[a,b]$ (changement de signe de $f$),
choisir $x_0 \in [a,b]$ et calculer $m = \min(\abs{f'(a)}, \abs{f'(b)})$
(valable lorsque $\abs{f'}$ est monotone sur $[a,b]$, comme le suppose le
cours ; sinon prendre le minimum de $\abs{f'}$ sur $[a,b]$).

\Etape{2} Itérer $x \leftarrow x - f(x)/f'(x)$.

\Etape{3} Après chaque itération, tester $\abs{f(x)}/m < E$ ; arrêter aussi si
$k$ atteint $k_{\max}$.

\begin{algo}[ : Newton-Raphson]
  \Require $f$, $f'$, $a$, $b$, $x_0$, $E$, $k_{\max}$
  \State $m \leftarrow \min(\abs{f'(a)}, \abs{f'(b)})$
  \State $x \leftarrow x_0$ ; $k \leftarrow 0$
  \Repeat
    \If{$f'(x) = 0$} \Print \og dérivée nulle : changer $x_0$ \fg{} \Stop \EndIf
    \State $x \leftarrow x - f(x)/f'(x)$
    \State $k \leftarrow k+1$
  \Until{$\abs{f(x)}/m < E$ \textbf{ou} $k \ge k_{\max}$}
  \If{$\abs{f(x)}/m \ge E$} \Print \og pas de convergence : changer $x_0$ \fg{}
  \Else \ \Print $x$, $k$
  \EndIf
\end{algo}

\Verif Pour $f(x) = (x-5)e^{x} + 5$ (exercice 1 de la fin du document), sur
$[4{,}9 ; 5]$ : $f'(x) = (x-4)e^{x}$ est croissante, $m = \abs{f'(4{,}9)} =
120{,}861$ ; depuis $x_0 = 5$ et avec $E = 10^{-12}$, l'algorithme s'arrête
après 4 itérations sur $x = 4{,}965114231744$.

\Refs \BF{alg.~2.3, p.~67}.

% ------------------------------------------------------------------
\section{Exercice II.2 : méthode de bissection}

\Enonce{Écrire l'algorithme de la méthode de bissection.}

\Pourquoi Si $f$ est continue et si $f(a)$ et $f(b)$ sont de signes opposés,
le théorème des valeurs intermédiaires garantit une racine dans $[a,b]$. En
gardant à chaque étape la moitié où le signe change, on conserve toujours une
racine, et la longueur de l'intervalle est divisée par 2. La méthode converge
donc \emph{toujours} (contrairement à Newton), mais lentement. Son avantage :
on connaît à l'avance le nombre d'étapes. Pour que l'intervalle final soit de
longueur inférieure à $\varepsilon$, il faut $(b-a)/2^n < \varepsilon$,
c'est-à-dire $n > \log_2\bigl((b-a)/\varepsilon\bigr)$.

Le test sur le signe se fait par le produit $f(a)f(c)$ : il fonctionne que $f$
soit croissante ou décroissante (le cours suppose $f$ croissante, ce qui
n'est pas nécessaire). On garde $f(a)$ en mémoire pour n'évaluer $f$ qu'une
fois par étape. Burden et Faires conseillent de tester le \emph{signe} plutôt
que le produit si l'on craint un dépassement de capacité \BF{p.~52}.

\Etapes
\Etape{1} Vérifier $f(a)f(b) < 0$.

\Etape{2} Tant que $b - a > \varepsilon$ : calculer $c = (a+b)/2$ et $f(c)$ ;
si $f(c) = 0$, la racine est $c$ ; si $f(a)f(c) < 0$, la racine est dans
$[a,c]$ et on pose $b \leftarrow c$ ; sinon elle est dans $[c,b]$ et on pose
$a \leftarrow c$, $f(a) \leftarrow f(c)$.

\Etape{3} Le résultat est le milieu du dernier intervalle.

\begin{algo}[ : bissection]
  \Require $f$, $a$, $b$, $\varepsilon$
  \State $f_a \leftarrow f(a)$
  \If{$f_a\,f(b) > 0$} \Print \og pas de changement de signe sur $[a,b]$ \fg{} \Stop \EndIf
  \While{$b - a > \varepsilon$}
    \State $c \leftarrow (a+b)/2$ ; $f_c \leftarrow f(c)$
    \If{$f_c = 0$} \Print $c$ \Stop \EndIf
    \If{$f_a\,f_c < 0$} \State $b \leftarrow c$
    \Else \State $a \leftarrow c$ ; $f_a \leftarrow f_c$
    \EndIf
  \EndWhile
  \Print $(a+b)/2$
\end{algo}

\Verif Avec la même fonction que ci-dessus, sur $[4{,}9 ; 5]$ et
$\varepsilon = 10^{-10}$ : $\log_2(0{,}1/10^{-10}) = 29{,}9$, il faut donc 30
étapes, et l'algorithme donne effectivement $x = 4{,}9651142317$ en 30 étapes.

\Refs \BF{alg.~2.1, p.~49}.

% ------------------------------------------------------------------
\section{Exercice II.3 : schéma de Lin pour les racines réelles}

\Enonce{Établir l'algorithme de ce schéma de Lin pour les solutions réelles.}

\Pourquoi Si $x$ est une racine, la dernière quantité de Horner est nulle :
$b_n = a_n + x\,b_{n-1}(x) = 0$, d'où $x = -a_n/b_{n-1}(x)$. La racine est donc
un point fixe de la fonction $g(x) = -a_n/b_{n-1}(x)$, et la méthode de Lin
consiste à itérer $x_{k+1} = g(x_k)$ : c'est une formule d'itération du premier
ordre au sens du chapitre II. Une telle itération converge vers une racine $r$
lorsque $\abs{g'(r)} < 1$ et s'en éloigne lorsque $\abs{g'(r)} > 1$ ; c'est ce
qui explique la remarque du cours (\og elle n'est pas toujours convergente \fg).
Quand l'itération a convergé, $b_n \approx 0$ et
$P_n(x) = (x - r)Q(x) + b_n$ : les nombres $b_0, \dots, b_{n-1}$ sont les
coefficients du quotient $Q$, de degré $n-1$, sur lequel on recommence pour
trouver la racine suivante (déflation).

\Etapes
\Etape{1} Partir d'une valeur $x_0$.

\Etape{2} Calculer par Horner $b_0 = a_0$, $b_j = a_j + x\,b_{j-1}$ pour
$j = 1, \dots, n-1$ (on s'arrête à $b_{n-1}$).

\Etape{3} Nouvelle valeur $x' = -a_n/b_{n-1}$ ; arrêter si
$\abs{x' - x} < \varepsilon$, ou si $k_{\max}$ itérations ont été faites sans
convergence (changer alors $x_0$).

\Etape{4} Déflation : recalculer $b_0, \dots, b_{n-1}$ avec la racine trouvée,
remplacer $(a_0, \dots, a_n)$ par $(b_0, \dots, b_{n-1})$, faire $n \leftarrow
n-1$ et recommencer à l'étape 1 tant que $n \ge 1$.

\begin{algo}[ : méthode de Lin avec déflation]
  \Require $n$, $a_0, \dots, a_n$, $x_0$, $\varepsilon$, $k_{\max}$
  \While{$n \ge 1$}
    \State $x \leftarrow x_0$ ; $k \leftarrow 0$
    \Repeat
      \State $k \leftarrow k+1$ ; $b_0 \leftarrow a_0$
      \For{$j$ allant de $1$ à $n-1$} \State $b_j \leftarrow a_j + x\,b_{j-1}$ \EndFor
      \If{$b_{n-1} = 0$} \Print \og changer $x_0$ \fg{} \Stop \EndIf
      \State $x' \leftarrow -a_n/b_{n-1}$ ; $d \leftarrow \abs{x' - x}$ ; $x \leftarrow x'$
    \Until{$d < \varepsilon$ \textbf{ou} $k \ge k_{\max}$}
    \If{$d \ge \varepsilon$} \Print \og pas de convergence \fg{} \Stop \EndIf
    \Print racine $x$
    \State $b_0 \leftarrow a_0$
    \For{$j$ allant de $1$ à $n-1$} \State $b_j \leftarrow a_j + x\,b_{j-1}$ \EndFor
    \For{$j$ allant de $0$ à $n-1$} \State $a_j \leftarrow b_j$ \EndFor
    \State $n \leftarrow n-1$
  \EndWhile
\end{algo}
(Pour $n = 1$, la boucle de Horner est vide et $x = -a_1/a_0$ est obtenu en une
itération.)

\Verif Pour $P(x) = x^3 - 6x^2 + 11x - 6 = (x-1)(x-2)(x-3)$, avec
$\varepsilon = 10^{-10}$ : depuis $x_0 = 0$ ou $1{,}5$, la méthode converge vers
$1$ (53 et 56 itérations) ; depuis $x_0 = 2{,}5$ ou $3{,}5$, elle converge vers
$3$ (7 itérations) ; elle ne converge jamais vers $2$. La théorie l'explique :
ici $g(x) = 6/(x^2 - 6x + 11)$ et $g'(x) = -6(2x-6)/(x^2-6x+11)^2$, d'où
$g'(1) = 2/3$ (convergence lente), $g'(2) = 4/3 > 1$ (racine répulsive) et
$g'(3) = 0$ (convergence rapide). C'est un exemple concret de la non-convergence
annoncée par le cours.

\Refs Méthode du cours (Lin, 1941). Pour la théorie des itérations de point
fixe et le critère $\abs{g'} < 1$ : \BF{§~2.2, th.~2.4, p.~61}. Horner
et déflation : \BF{alg.~2.7, p.~94 ; déflation p.~95}.

\chapter{Chapitre III : interpolation et moindres carrés}

% ------------------------------------------------------------------
\section{Exercice III.1 : programme d'interpolation de Lagrange}

\Enonce{Traduire cet algorithme dans un des langages de programmation de votre
choix (Pascal, Fortran, C).}

\Pourquoi Le polynôme de Lagrange s'écrit
$f(x_p) = \sum_k y_k\,f_k(x_p)$ avec
$f_k(x_p) = \prod_{i \ne k}(x_p - x_i)/(x_k - x_i)$. L'algorithme du cours
calcule chaque produit $f_k(x_p)$ dans la variable $y_s$ (boucle intérieure),
puis l'ajoute, multiplié par $y_k$, à la somme $y_p$ (boucle extérieure). On
n'a jamais besoin des coefficients du polynôme : on évalue directement sa
valeur en $x_p$.

\Etapes
\Etape{1} Écrire une fonction qui reçoit la liste des $x_k$, la liste des $y_k$
et $x_p$ (en Python, \texttt{len(x)} donne le nombre de points $n+1$).

\Etape{2} Traduire les deux boucles imbriquées ; le test $i \ne k$ évite le
facteur nul $(x_k - x_k)$ au dénominateur.

\Etape{3} En Python, les listes commencent à l'indice 0 :
\texttt{range(len(x))} parcourt les indices $0$ à $n$ au lieu de $1$ à $n+1$.

\programme{lagrange.py}

\Verif Avec les données de l'exercice III.4 ci-dessous, le programme affiche
$y_p = 15{,}793629$ pour $x_p = 102$.

\Refs \BF{§~3.1, p.~106--108} pour la formule de Lagrange ; une alternative
programmable est l'algorithme de Neville \BF{alg.~3.1, p.~120}.

% ------------------------------------------------------------------
\section{\texorpdfstring{Complément : formules de $a_0$ et $a_1$ pour la droite des moindres
carrés}{Complément : formules de a0 et a1 pour la droite des moindres carrés}}

\Enonce{Le cours demande d'établir en exercice les formules de $a_0$ et $a_1$
pour $f(x) = a_0 + a_1x$.}

\Pourquoi Le reste $R(a_0, a_1) = \sum_k\bigl[f_k - (a_0 + a_1x_k)\bigr]^2$ est
une fonction du second degré des deux inconnues, positive : son minimum est
atteint là où ses deux dérivées partielles s'annulent. On obtient deux
équations linéaires (les \og équations normales \fg).

\Etapes
\Etape{1} Dérivées partielles :
\[
  \frac{\partial R}{\partial a_0} = -2\sum_k\bigl(f_k - a_0 - a_1x_k\bigr) = 0,
  \qquad
  \frac{\partial R}{\partial a_1} = -2\sum_k x_k\bigl(f_k - a_0 - a_1x_k\bigr) = 0 .
\]

\Etape{2} En développant :
\[
  m\,a_0 + \Bigl(\sum x_k\Bigr)a_1 = \sum f_k, \qquad
  \Bigl(\sum x_k\Bigr)a_0 + \Bigl(\sum x_k^2\Bigr)a_1 = \sum x_kf_k .
\]

\Etape{3} Par la règle de Cramer, avec $D = m\sum x_k^2 - \bigl(\sum x_k\bigr)^2$
(non nul dès que deux $x_k$ sont distincts) :
\[
  a_0 = \frac{\sum x_k^2\sum f_k - \sum x_k\sum x_kf_k}{D}, \qquad
  a_1 = \frac{m\sum x_kf_k - \sum x_k\sum f_k}{D} ,
\]
ce sont les formules du cours.

\Refs \BF{§~8.1, p.~507--508}.

% ------------------------------------------------------------------
\section{\texorpdfstring{Exercice III.2 : calcul des coefficients $b_k$ et $A_{kj}$}{Exercice III.2 : calcul des coefficients bk et Akj}}

\Enonce{Écrire les algorithmes de calculs des coefficients $b_k$ et $A_{kj}$.}

\Pourquoi Les équations normales $\sum_j A_{kj}a_j = b_k$ viennent de
$\partial R/\partial a_k = 0$ pour
$R = \sum_i\bigl[f_i - \sum_j a_jg_j(x_i)\bigr]^2$. Chaque coefficient est une
somme sur les $m$ mesures : une boucle sur $i$ suffit. Deux remarques rendent
le calcul plus économique :
\begin{itemize}
  \item $A_{kj} = A_{jk}$ (la matrice est symétrique) : on ne calcule que
        $j \ge k$ et on recopie ;
  \item les valeurs $g_k(x_i)$ servent plusieurs fois : on les calcule une seule
        fois et on les range dans un tableau $G_{ik} = g_k(x_i)$.
\end{itemize}
Une fois $A$ et $b$ calculés, le système se résout par la méthode de Gauss
avec pivot (exercice I.3).

\Etapes
\Etape{1} Remplir $G_{ik} = g_k(x_i)$ pour $i = 1, \dots, m$ et
$k = 0, \dots, n$.

\Etape{2} Pour chaque $k$ : $b_k = \sum_i G_{ik}f_i$, puis pour $j \ge k$ :
$A_{kj} = \sum_i G_{ik}G_{ij}$ et $A_{jk} = A_{kj}$.

\Etape{3} Résoudre $A\,a = b$ par Gauss.

\begin{algo}[ : équations normales]
  \Require $m$, $(x_i, f_i)$, $n$, les fonctions $g_0, \dots, g_n$
  \For{$i$ allant de $1$ à $m$}
    \For{$k$ allant de $0$ à $n$} \State $G_{ik} \leftarrow g_k(x_i)$ \EndFor
  \EndFor
  \For{$k$ allant de $0$ à $n$}
    \State $b_k \leftarrow 0$
    \For{$i$ allant de $1$ à $m$} \State $b_k \leftarrow b_k + G_{ik}\,f_i$ \EndFor
    \For{$j$ allant de $k$ à $n$}
      \State $A_{kj} \leftarrow 0$
      \For{$i$ allant de $1$ à $m$} \State $A_{kj} \leftarrow A_{kj} + G_{ik}\,G_{ij}$ \EndFor
      \State $A_{jk} \leftarrow A_{kj}$
    \EndFor
  \EndFor
  \State résoudre $A\,a = b$ par la méthode de Gauss avec pivot \Comment{exercice I.3}
  \Print $a_0, \dots, a_n$
\end{algo}

\Refs \BF{§~8.1, équations normales, p.~508--510}.

% ------------------------------------------------------------------
\section{Exercice III.3 : moindres carrés sur un tableau de mesures}

\Enonce{$x = 1, 2, 3, 4, 5$ ; $y = 1{,}6$ ; $2$ ; $2{,}3$ ; $2{,}4$ ; $2{,}5$.
Déterminer par les moindres carrés a) $y = a_0 + a_1/x$ ;
b) $y = a_0 + a_1/x + a_2e^{x}$.}

\Pourquoi On applique l'exercice précédent avec $g_0 = 1$, $g_1 = 1/x$
(et $g_2 = e^{x}$ en b). Les modèles sont non linéaires en $x$ mais
\emph{linéaires en les coefficients} : les équations normales restent un
système linéaire. Le choix $1/x$ est naturel car les mesures croissent de moins
en moins vite.

\Etapes
\Etape{1} Cas a). Sommes nécessaires ($m = 5$) :
\[
  \sum\tfrac1{x_i} = \tfrac{137}{60} = 2{,}283333, \quad
  \sum\tfrac1{x_i^2} = 1{,}463611, \quad
  \sum y_i = 10{,}8, \quad
  \sum\tfrac{y_i}{x_i} = 4{,}466667 .
\]

\Etape{2} Système normal :
\[
  \begin{pmatrix} 5 & 2{,}283333\\ 2{,}283333 & 1{,}463611 \end{pmatrix}
  \begin{pmatrix} a_0\\ a_1 \end{pmatrix}
  = \begin{pmatrix} 10{,}8\\ 4{,}466667 \end{pmatrix}.
\]
Déterminant $D = 5 \times 1{,}463611 - 2{,}283333^2 = 2{,}104444$ ; Cramer :
\[
  a_0 = \frac{1{,}463611 \times 10{,}8 - 2{,}283333 \times 4{,}466667}{D}
      = \frac{5{,}608111}{2{,}104444} = 2{,}664889,
\]
\[
  a_1 = \frac{5 \times 4{,}466667 - 2{,}283333 \times 10{,}8}{D}
      = \frac{-2{,}326667}{2{,}104444} = -1{,}105597 .
\]
Donc $y \approx 2{,}664889 - 1{,}105597/x$, avec un reste $R = 0{,}017529$
(écarts aux mesures : $0{,}0407$ ; $-0{,}1121$ ; $0{,}0036$ ; $0{,}0115$ ;
$0{,}0562$).

\Etape{3} Cas b). Sommes supplémentaires :
$\sum e^{x_i} = 233{,}204184$, $\sum e^{x_i}/x_i = 56{,}440158$,
$\sum e^{2x_i} = 25472{,}839781$, $\sum y_ie^{x_i} = 567{,}392556$. Le système
$3\times 3$
\[
  \begin{pmatrix}
    5 & 2{,}283333 & 233{,}204184\\
    2{,}283333 & 1{,}463611 & 56{,}440158\\
    233{,}204184 & 56{,}440158 & 25472{,}839781
  \end{pmatrix}
  \begin{pmatrix} a_0\\ a_1\\ a_2 \end{pmatrix}
  = \begin{pmatrix} 10{,}8\\ 4{,}466667\\ 567{,}392556 \end{pmatrix}
\]
résolu par Gauss avec pivot donne
\[
  a_0 = 2{,}567598, \qquad a_1 = -0{,}990980, \qquad a_2 = 0{,}000964,
\]
avec $R = 0{,}009502$.

\Etape{4} Interprétation. Ajouter une fonction ne peut que diminuer $R$ (le
modèle a contient le modèle b avec $a_2 = 0$) : $R$ passe de $0{,}0175$ à
$0{,}0095$. Le coefficient $a_2$ est petit parce que $e^{x}$ est très grand
($e^5 \approx 148$). Les nombres de la matrice sont d'ordres de grandeur très
différents (de $1{,}46$ à $25\,473$) : son conditionnement vaut environ
$1{,}4\times 10^{5}$, ce qui justifie le pivot et la double précision.

\Verif Valeurs obtenues par \fichier{calculs.py} et \fichier{complements.py}.

\Refs \BF{§~8.1, p.~506--513}.

% ------------------------------------------------------------------
\section{Exercice III.4 : interpolation de Lagrange directe et inverse}

\Enonce{Du tableau $x$ : $93{,}0$ ; $96{,}2$ ; $100{,}0$ ; $104{,}2$ ; $108{,}7$,
$y$ : $11{,}38$ ; $12{,}80$ ; $14{,}70$ ; $17{,}07$ ; $19{,}91$, déterminer $y_p$
pour $x_p = 102$ et $x_p$ pour $y_p = 13{,}5$.}

\Pourquoi Avec 5 points, on utilise le polynôme de Lagrange de degré 4. Pour la
question inverse (trouver $x$ connaissant $y$), le cours indique d'échanger les
rôles de $x$ et de $y$ : on interpole $x$ comme fonction de $y$. C'est possible
parce que $y$ est strictement croissante sur le tableau (la fonction réciproque
existe).

\Etapes
\Etape{1} Poids de Lagrange en $x_p = 102$ :
$f_k(102) = 0{,}043410$ ; $-0{,}218191$ ; $0{,}791622$ ; $0{,}413053$ ;
$-0{,}029894$ (leur somme vaut 1, ce qui est un bon contrôle : le polynôme
d'interpolation de la fonction constante $1$ est $1$).

\Etape{2} $y_p = \sum_k y_kf_k(102) = 15{,}793629$, soit $y_p \approx 15{,}79$.

\Etape{3} Interpolation inverse : les nœuds sont maintenant les $y_k$ et les
valeurs les $x_k$. Poids en $13{,}5$ : $-0{,}084007$ ; $0{,}710729$ ;
$0{,}435996$ ; $-0{,}069802$ ; $0{,}007084$, d'où $x_p = 97{,}655750$.

\Verif Le programme \fichier{lagrange.py} donne $y_p = 15{,}793629$ et (en
échangeant les deux listes) $x_p = 97{,}655750$. La réponse du cours
($15{,}79$ et $97{,}65$) est retrouvée ; arrondi au centième, $x_p$ vaut
plutôt $97{,}66$ ($97{,}65$ est la valeur tronquée).

\Refs \BF{§~3.1, p.~106--108} ; interpolation inverse itérée : \BF{p.~122}.

\chapter{Chapitre IV : intégration et dérivation numériques}

% ------------------------------------------------------------------
\section{Exercice IV.1 : dérivation à partir de l'interpolation quadratique}

\Enonce{Établir la formule de dérivation numérique à partir de l'interpolation
quadratique.}

\Pourquoi On ne connaît $f$ qu'en quelques points. On la remplace par le
polynôme de degré 2 qui passe par trois points, et on dérive ce polynôme. Avec
trois points au lieu de deux, l'erreur passe de l'ordre $h$ (formule classique)
à l'ordre $h^2$.

\Etapes
\Etape{1} Points équidistants $x_1 = x - h$, $x_2 = x$, $x_3 = x + h$. On part de
$f(t) \approx f(x_1)f_1(t) + f(x_2)f_2(t) + f(x_3)f_3(t)$ et on dérive les
trois polynômes de base. Par exemple
$f_1(t) = (t - x_2)(t - x_3)/\bigl((x_1 - x_2)(x_1 - x_3)\bigr)$ donne
$f_1'(t) = (2t - x_2 - x_3)/\bigl((x_1 - x_2)(x_1 - x_3)\bigr)$.

\Etape{2} Au point $t = x_2$ :
\[
\begin{aligned}
  f_1'(x_2) &= \frac{x_2 - x_3}{(x_1 - x_2)(x_1 - x_3)} = \frac{-h}{(-h)(-2h)} = -\frac{1}{2h},\\
  f_2'(x_2) &= \frac{2x_2 - x_1 - x_3}{(x_2 - x_1)(x_2 - x_3)} = 0,\\
  f_3'(x_2) &= \frac{x_2 - x_1}{(x_3 - x_1)(x_3 - x_2)} = \frac{h}{(2h)(h)} = \frac{1}{2h}.
\end{aligned}
\]
D'où la formule centrée
\[
  f'(x) \approx \frac{f(x+h) - f(x-h)}{2h} .
\]

\Etape{3} Erreur. Les développements de Taylor à l'ordre 3,
$f(x \pm h) = f(x) \pm hf'(x) + \tfrac{h^2}{2}f''(x) \pm \tfrac{h^3}{6}f'''(\xi_\pm)$,
donnent par différence
\[
  \frac{f(x+h) - f(x-h)}{2h} = f'(x) + \frac{h^2}{12}\bigl(f'''(\xi_+) + f'''(\xi_-)\bigr)
  = f'(x) + \frac{h^2}{6}f'''(\xi),
\]
la dernière égalité venant du théorème des valeurs intermédiaires ($f'''$
continue). On retrouve l'erreur du cours $E = -\tfrac{1}{6}h^2f'''(\xi)$ :
$f'(x) = \frac{f(x+h) - f(x-h)}{2h} - \frac{h^2}{6}f'''(\xi)$.

\Etape{4} Au point $t = x_1$ (dérivée en bout d'intervalle, utile au bord d'un
tableau), le même calcul donne $f_1'(x_1) = -3/(2h)$, $f_2'(x_1) = 2/h$,
$f_3'(x_1) = -1/(2h)$, c'est-à-dire la seconde formule du cours
$f'(x) \approx \bigl[-3f(x) + 4f(x+h) - f(x+2h)\bigr]/(2h)$, d'erreur
$+\tfrac13h^2f'''(\xi)$.

\Refs \BF{§~4.1, formules à trois points (4.4) et (4.5), p.~174--175}.

% ------------------------------------------------------------------
\section{Exercice IV.2 : vitesses à partir d'un tableau de positions}

\Enonce{$t = 1{,}0$ ; $1{,}1$ ; \dots ; $1{,}5$ et
$x = 1{,}6487$ ; $1{,}7333$ ; $1{,}8221$ ; $1{,}9155$ ; $2{,}0138$ ; $2{,}1170$.
Calculer les vitesses à $t = 1{,}05$ ; $1{,}15$ ; $1{,}25$ ; $1{,}35$ et $1{,}45$.}

\Pourquoi Les instants demandés sont les \emph{milieux} des intervalles du
tableau. La formule centrée de l'exercice IV.1, appliquée avec le pas $h/2$
autour du milieu $t + h/2$, n'utilise justement que $x(t)$ et $x(t+h)$ :
\[
  v\Bigl(t + \frac h2\Bigr) \approx \frac{x(t+h) - x(t)}{h},
  \qquad\text{erreur } -\frac{(h/2)^2}{6}x'''(\xi) = -\frac{h^2}{24}x'''(\xi).
\]
C'est la deuxième formule classique du cours ($f(x + h/2) - f(x - h/2)$ divisé
par $h$). On obtient donc une précision d'ordre $h^2$ avec les seules données.

\Etapes
\Etape{1} Différences successives des positions :
$0{,}0846$ ; $0{,}0888$ ; $0{,}0934$ ; $0{,}0983$ ; $0{,}1032$.

\Etape{2} Division par $h = 0{,}1$ :
\[
  v(1{,}05) \approx 0{,}846, \quad v(1{,}15) \approx 0{,}888, \quad
  v(1{,}25) \approx 0{,}934, \quad v(1{,}35) \approx 0{,}983, \quad
  v(1{,}45) \approx 1{,}032 .
\]

\Etape{3} Précision. Les positions sont données à $5\times 10^{-5}$ près ; une
différence divisée par $h = 0{,}1$ peut donc être fausse d'environ
$2 \times 5\times 10^{-5}/0{,}1 = 10^{-3}$. Les vitesses ne sont fiables qu'à
$10^{-3}$ près environ : c'est l'erreur d'arrondi, qui croît quand $h$ diminue,
alors que l'erreur de troncature décroît \BF{p.~179}.

\Verif On constate que les données coïncident avec $e^{t/2}$ à
$4{,}7\times 10^{-5}$ près (calcul fait sur les six points). Si c'est bien la
loi du mouvement, la vitesse exacte est $\tfrac12e^{t/2}$, soit
$0{,}8452$ ; $0{,}8886$ ; $0{,}9341$ ; $0{,}9820$ ; $1{,}0324$ : les écarts
avec nos valeurs restent inférieurs à $10^{-3}$, comme prévu à l'étape 3.

\Refs \BF{§~4.1, formule (4.5) p.~175 ; erreurs d'arrondi p.~178--179}.

% ------------------------------------------------------------------
\section{\texorpdfstring{Exercice IV.3 : trapèzes à la main pour $\int_1^2 \dd x/x$}{Exercice IV.3 : trapèzes à la main pour l'intégrale de 1/x}}

\Enonce{Calculer à la main $I = \int_1^2 \dd x/x$ pour $h = 0{,}2$, puis pour
$h = 0{,}1$. Comparer à la valeur exacte.}

\Pourquoi La valeur exacte est connue, $I = \ln 2 = 0{,}693147$ : l'exercice
permet de vérifier la formule composite des trapèzes et la loi d'erreur en
$h^2$ établie dans le cours, $E = -\tfrac{h^2}{12}(b-a)f''(\eta)$.

\Etapes
\Etape{1} $h = 0{,}2$ ($n = 5$). Valeurs de $f(x) = 1/x$ :
$1$ ; $0{,}833333$ ; $0{,}714286$ ; $0{,}625$ ; $0{,}555556$ ; $0{,}5$. Somme
des valeurs intérieures : $2{,}728175$.
\[
  T_{0,2} = \frac{0{,}2}{2}\bigl[1 + 2 \times 2{,}728175 + 0{,}5\bigr]
          = 0{,}1 \times 6{,}956349 = 0{,}695635 .
\]

\Etape{2} $h = 0{,}1$ ($n = 10$). Valeurs intérieures : $0{,}909091$ ;
$0{,}833333$ ; $0{,}769231$ ; $0{,}714286$ ; $0{,}666667$ ; $0{,}625$ ;
$0{,}588235$ ; $0{,}555556$ ; $0{,}526316$, de somme $6{,}187714$.
\[
  T_{0,1} = 0{,}05\bigl[1 + 2 \times 6{,}187714 + 0{,}5\bigr] = 0{,}693771 .
\]

\Etape{3} Comparaison. Erreurs : $T_{0,2} - \ln 2 = 0{,}002488$ et
$T_{0,1} - \ln 2 = 0{,}000624$. Leur rapport vaut $3{,}99 \approx 4$ : diviser
$h$ par 2 divise l'erreur par 4, c'est la loi en $h^2$. Les erreurs sont
positives, ce qu'annonce la formule car $f''(x) = 2/x^3 > 0$. Majoration : sur
$[1,2]$, $\abs{f''} \le 2$, d'où $\abs{E} \le h^2/6$, soit $0{,}0067$ et
$0{,}0017$ : elle est bien respectée.

\Etape{4} Prolongement (Richardson--Romberg). Comme l'erreur est en $h^2$, la
combinaison $(4T_{0,1} - T_{0,2})/3 = 0{,}693150$ élimine le terme principal :
l'erreur tombe à $3\times 10^{-6}$.

\Refs \BF{th.~4.5, p.~205 (trapèzes composites) ; intégration de Romberg,
§~4.5, p.~211--216}.

% ------------------------------------------------------------------
\section{Exercice IV.4 : trapèzes sans tableaux}

\Enonce{Écrire un autre algorithme qui n'utilise pas les tableaux.}

\Pourquoi Dans l'algorithme du cours, chaque valeur $f_k$ n'est utilisée
qu'une fois : il est inutile de la stocker. On l'ajoute à la somme dès qu'elle
est calculée. La mémoire utilisée ne dépend plus de $n$. On calcule les nœuds
par $x = a + kh$ plutôt que par additions répétées $x \leftarrow x + h$, pour
ne pas accumuler les erreurs d'arrondi.

\Etapes
\Etape{1} Initialiser la somme avec les deux extrémités (poids $1/2$).

\Etape{2} Ajouter les $n-1$ valeurs intérieures (poids $1$).

\Etape{3} Multiplier par $h$.

\begin{algo}[ : trapèzes composites sans tableaux]
  \Require la fonction $f$, $a$, $b$, $n$
  \State $h \leftarrow (b-a)/n$
  \State $S \leftarrow \bigl(f(a) + f(b)\bigr)/2$
  \For{$k$ allant de $1$ à $n-1$}
    \State $S \leftarrow S + f(a + kh)$
  \EndFor
  \State $S \leftarrow h\,S$
  \Print $S$
\end{algo}

\Verif Cet algorithme est programmé dans \fichier{ex11_trapeze.pas} et
\fichier{ex11_trapeze.f90} (exercice 11 de la fin du document) ; appliqué à
$1/x$ sur $[1,2]$, il redonne les valeurs de l'exercice IV.3.

\Refs \BF{p.~205} (même organisation dans l'algorithme 4.1).

% ------------------------------------------------------------------
\section{Exercice IV.5 : algorithme de Simpson}

\Enonce{Écrire un algorithme pour la méthode de Simpson.}

\Pourquoi La formule de Simpson s'applique sur des paires d'intervalles (un
polynôme de degré 2 sur trois points) : le nombre $n$ de sous-intervalles doit
donc être \emph{pair}. Les poids sont $1, 4, 2, 4, \dots, 2, 4, 1$ : $4$ aux
nœuds d'indice impair (milieux des paires), $2$ aux nœuds pairs intérieurs
(extrémités communes à deux paires). L'erreur composite vaut
$-\tfrac{b-a}{180}h^4f^{(4)}(\mu)$ \BF{th.~4.4, p.~204} : la méthode est
exacte pour les polynômes de degré 3.

\Etapes
\Etape{1} Vérifier que $n$ est pair.

\Etape{2} Accumuler séparément la somme $S_1$ des valeurs aux nœuds impairs et
la somme $S_2$ des valeurs aux nœuds pairs intérieurs (c'est ce que fait
l'algorithme 4.1 de Burden et Faires avec ses variables $XI_1$ et $XI_2$).

\Etape{3} $I \approx \tfrac h3\bigl[f(a) + 4S_1 + 2S_2 + f(b)\bigr]$.

\begin{algo}[ : Simpson composite]
  \Require $f$, $a$, $b$, $n$ (pair)
  \If{$n$ impair} \Print \og $n$ doit être pair \fg{} \Stop \EndIf
  \State $h \leftarrow (b-a)/n$ ; $S_1 \leftarrow 0$ ; $S_2 \leftarrow 0$
  \For{$k$ allant de $1$ à $n-1$}
    \State $x \leftarrow a + kh$
    \If{$k$ impair} \State $S_1 \leftarrow S_1 + f(x)$
    \Else \State $S_2 \leftarrow S_2 + f(x)$
    \EndIf
  \EndFor
  \State $I \leftarrow \frac{h}{3}\bigl(f(a) + 4S_1 + 2S_2 + f(b)\bigr)$
  \Print $I$
\end{algo}

\Verif Voir l'exercice IV.7 (calcul de $\pi$) et le programme
\fichier{ex12_mc.pas} (exercice 12 de la fin du document).

\Refs \BF{alg.~4.1, p.~205 ; th.~4.4, p.~204}.

% ------------------------------------------------------------------
\section{\texorpdfstring{Exercice IV.6 : formule d'intégration de degré 3 sur $a$, $c$, $d$, $b$}{Exercice IV.6 : formule d'intégration de degré 3 sur a, c, d, b}}

\Enonce{Établir une formule d'intégration numérique qui utilise une
interpolation de degré 3 et qui interpole les points $a$, $b$,
$c = (a+b)/4$ et $d = 3(a+b)/4$.}

\Attention Les points $(a+b)/4$ et $3(a+b)/4$ ne sont dans $[a,b]$, aux
quarts de l'intervalle, que si $a = 0$. On interprète donc l'énoncé comme
$c = a + (b-a)/4$ et $d = a + 3(b-a)/4$, ce qui coïncide avec lui lorsque
$a = 0$.

\Pourquoi Intégrer le polynôme d'interpolation de degré 3 revient à chercher
des poids $w_k$ tels que $\int_a^b f \approx \sum_k w_kf(x_k)$ soit exact pour
$1, x, x^2, x^3$. Les quatre nœuds étant symétriques par rapport au milieu de
$[a,b]$, les poids le sont aussi : il ne reste que deux inconnues.

\Etapes
\Etape{1} Changement de variable $x = a + (b-a)t$, $t \in [0,1]$ :
$\int_a^b f(x)\,\dd x = (b-a)\int_0^1 f\bigl(a + (b-a)t\bigr)\dd t$. Les nœuds
deviennent $t = 0$, $\tfrac14$, $\tfrac34$, $1$.

\Etape{2} Par symétrie, poids $\alpha$ en $0$ et $1$, et $\beta$ en $\tfrac14$ et
$\tfrac34$. Exactitude pour $1$ : $2\alpha + 2\beta = 1$. Exactitude pour
$t^2$ : $\alpha(0 + 1) + \beta\bigl(\tfrac1{16} + \tfrac9{16}\bigr) = \tfrac13$,
soit $\alpha + \tfrac58\beta = \tfrac13$.

\Etape{3} De la première, $\alpha = \tfrac12 - \beta$ ; en reportant :
$\tfrac12 - \tfrac38\beta = \tfrac13$, d'où $\beta = \tfrac49$ et
$\alpha = \tfrac1{18}$. L'exactitude pour $t$ et $t^3$ est automatique (les
fonctions $t - \tfrac12$ et $(t - \tfrac12)^3$ sont impaires par rapport au
milieu, et les poids sont symétriques).

\Etape{4} Formule obtenue :
\[
  \int_a^b f(x)\,\dd x \approx \frac{b-a}{18}\Bigl[f(a) + 8f(c) + 8f(d) + f(b)\Bigr] .
\]

\Etape{5} Erreur. Pour $t^4$, la formule donne
$\tfrac49\bigl(\tfrac1{256} + \tfrac{81}{256}\bigr) + \tfrac1{18} = \tfrac{19}{96}$
au lieu de $\tfrac15$ : l'écart (exact moins formule) vaut $\tfrac1{480}$. Comme
$(t^4)^{(4)} = 24$, l'erreur est, pour $f^{(4)}$ à peu près constante,
\[
  E \approx \frac{(b-a)^5}{11\,520}\,f^{(4)} .
\]
Le noyau de Peano de cette formule change légèrement de signe (calcul
numérique) : on ne peut pas écrire rigoureusement $E = K f^{(4)}(\xi)$, mais
on a la majoration $\abs{E} \le 9{,}13\times 10^{-5}(b-a)^5\max\abs{f^{(4)}}$.

\Verif Sur $\int_0^1 e^{x}\dd x = e - 1$ : l'erreur vaut $1{,}44\times 10^{-4}$,
proche de la prévision $e^{0{,}5}/11\,520 = 1{,}43\times 10^{-4}$. À titre de
comparaison, la formule de Newton-Cotes à quatre points équidistants
(\og 3/8 de Simpson \fg, erreur $-\tfrac{3}{80}h^5f^{(4)}$ avec
$h = (b-a)/3$, soit $(b-a)^5/6480$) donne $-2{,}58\times 10^{-4}$ : la
formule trouvée est environ deux fois plus précise.

\Refs \BF{§~4.3 : degré de précision, p.~196 ; formules de Newton-Cotes
fermées et règle des 3/8, th.~4.2, p.~196--197}.

% ------------------------------------------------------------------
\section{\texorpdfstring{Exercice IV.7 : Simpson à la main pour $\int_0^1 4/(1+x^2)\,\dd x$}{Exercice IV.7 : Simpson à la main pour l'intégrale de 4/(1+x²)}}

\Enonce{Utiliser la méthode de Simpson pour calculer manuellement
$I = \int_0^1 \frac{4}{1+x^2}\dd x$ avec $h = 0{,}1$.}

\Pourquoi Une primitive est $4\arctan x$, donc $I = 4\arctan 1 = \pi$. On peut
ainsi mesurer l'erreur exacte de la méthode.

\Etapes
\Etape{1} $n = 10$ (pair). Valeurs $f_k = 4/(1 + x_k^2)$ :
\begin{center}
\begin{tabular}{c|cccccc}
  $k$ & 0 & 1 & 2 & 3 & 4 & 5\\ \hline
  $f_k$ & 4 & 3,960396 & 3,846154 & 3,669725 & 3,448276 & 3,2\\[4pt]
  $k$ & 6 & 7 & 8 & 9 & 10 & \\ \hline
  $f_k$ & 2,941176 & 2,684564 & 2,439024 & 2,209945 & 2 &
\end{tabular}
\end{center}

\Etape{2} Sommes : indices impairs
$S_1 = 3{,}960396 + 3{,}669725 + 3{,}2 + 2{,}684564 + 2{,}209945 = 15{,}724629$ ;
indices pairs intérieurs
$S_2 = 3{,}846154 + 3{,}448276 + 2{,}941176 + 2{,}439024 = 12{,}674631$.

\Etape{3}
\[
  I \approx \frac{0{,}1}{3}\bigl[4 + 4 \times 15{,}724629 + 2 \times 12{,}674631 + 2\bigr]
    = \frac{0{,}1}{3} \times 94{,}247778 = 3{,}1415926 .
\]

\Etape{4} Comparaison : $\pi = 3{,}14159265$ ; l'erreur vaut
$-4\times 10^{-8}$. La majoration $\frac{b-a}{180}h^4\max\abs{f^{(4)}}$ avec
$\max_{[0,1]}\abs{f^{(4)}} = 96$ (atteint en $x = 0$) vaut $5{,}3\times
10^{-5}$ : elle est respectée, mais très pessimiste ici.

\Verif \fichier{calculs.py} : $S = 3{,}1415926139$, erreur
$-3{,}965\times 10^{-8}$.

\Refs \BF{th.~4.4, p.~204 ; alg.~4.1, p.~205}.

\chapter{Chapitre V : équations différentielles et intégro-différentielles}

Dans tout ce chapitre, les algorithmes ont été testés sur
$y' = -2xy$, $y(0) = 1$, de solution exacte $y = e^{-x^2}$, en calculant
$y(1)$ avec $n$ pas ($h = 1/n$).

% ------------------------------------------------------------------
\section{Exercice V.1 : algorithme de Runge-Kutta d'ordre 2}

\Enonce{Écrire l'algorithme de la méthode de Runge-Kutta d'ordre 2.}

\Pourquoi On a $y_{k+1} = y_k + \int_{x_k}^{x_{k+1}} f(x,y)\,\dd x$. La
méthode des trapèzes donnerait
$y_{k+1} = y_k + \tfrac h2\bigl[f(x_k,y_k) + f(x_{k+1},y_{k+1})\bigr]$, mais
$y_{k+1}$ est inconnu dans le second membre. On le remplace par la prévision
d'Euler $y_k + hf(x_k,y_k)$ : la formule devient explicite et reste d'ordre 2
(erreur locale $O(h^3)$, erreur globale $O(h^2)$). Il faut deux évaluations
de $f$ par pas ; on stocke la première dans une variable pour ne pas la
recalculer.

\Etapes
\Etape{1} $L_1 = hf(x_k, y_k)$ (pente au début du pas).

\Etape{2} $L_2 = hf(x_k + h, y_k + L_1)$ (pente au point prévu par Euler).

\Etape{3} $y_{k+1} = y_k + (L_1 + L_2)/2$, puis $x_{k+1} = x_k + h$.

\begin{algo}[ : Runge-Kutta d'ordre 2]
  \Require $f$, $x_0$, $y_0$, $h$, $n$
  \State $x \leftarrow x_0$ ; $y \leftarrow y_0$
  \For{$k$ allant de $1$ à $n$}
    \State $L_1 \leftarrow h\,f(x, y)$
    \State $L_2 \leftarrow h\,f(x + h, y + L_1)$
    \State $y \leftarrow y + (L_1 + L_2)/2$
    \State $x \leftarrow x + h$
    \Print $x$, $y$
  \EndFor
\end{algo}

\Verif Erreurs sur $y(1)$ : $1{,}17\times 10^{-3}$ ($n = 10$),
$3{,}01\times 10^{-4}$ ($n = 20$), $7{,}60\times 10^{-5}$ ($n = 40$). Chaque
division de $h$ par 2 divise l'erreur par $3{,}9$ puis $4{,}0$ : ordre 2
global.

\Refs \BF{méthode d'Euler modifiée, p.~286}.

% ------------------------------------------------------------------
\section{Exercice V.2 : algorithme de Runge d'ordre 4}

\Enonce{Établir l'algorithme de la méthode de Runge d'ordre 4.}

\Pourquoi La méthode applique la formule de Simpson à l'intégrale
$L = \int_{x_k}^{x_k+h} f\,\dd x$ ; les valeurs inconnues de $y$ au milieu et
à la fin du pas sont remplacées par des prévisions. Le cours donne les quatre
quantités $L_1, \dots, L_4$ à calculer dans un ordre précis, car chacune
utilise les précédentes.

\Etapes
\Etape{1} $L_1 = hf(x_k, y_k)$.

\Etape{2} $L_2 = hf(x_k + h, y_k + L_1)$ (prévision d'Euler en fin de pas).

\Etape{3} $L_3 = hf(x_k + h, y_k + L_2)$ (prévision améliorée en fin de pas).

\Etape{4} $L_4 = hf(x_k + h/2, y_k + L_1/2)$ (prévision au milieu du pas).

\Etape{5} $y_{k+1} = y_k + (L_1 + 4L_4 + L_3)/6$ : ce sont les poids
$1, 4, 1$ de Simpson.

\begin{algo}[ : Runge d'ordre 4 (formules du cours)]
  \Require $f$, $x_0$, $y_0$, $h$, $n$
  \State $x \leftarrow x_0$ ; $y \leftarrow y_0$
  \For{$k$ allant de $1$ à $n$}
    \State $L_1 \leftarrow h\,f(x, y)$
    \State $L_2 \leftarrow h\,f(x + h, y + L_1)$
    \State $L_3 \leftarrow h\,f(x + h, y + L_2)$
    \State $L_4 \leftarrow h\,f(x + h/2, y + L_1/2)$
    \State $y \leftarrow y + (L_1 + 4L_4 + L_3)/6$
    \State $x \leftarrow x + h$
    \Print $x$, $y$
  \EndFor
\end{algo}

\Verif et \textbf{remarque sur l'ordre.} Le cours annonce une erreur
$O(h^5)$. Les essais numériques montrent un ordre plus faible pour ces
formules précises : en divisant $h$ par 2, l'erreur globale est divisée par
environ 8 (et non 16), soit une erreur globale $O(h^3)$ et une erreur locale
$O(h^4)$.
\begin{center}
\begin{tabular}{lccc}
  \toprule
  & $n = 10$ & $n = 20$ & $n = 40$\\ \midrule
  $y' = y$ : Runge (cours) & $-1{,}05\times 10^{-4}$ & $-1{,}36\times 10^{-5}$ & $-1{,}74\times 10^{-6}$\\
  $y' = y$ : Runge-Kutta classique & $-2{,}08\times 10^{-6}$ & $-1{,}36\times 10^{-7}$ & $-8{,}67\times 10^{-9}$\\
  $y' = -2xy$ : Runge (cours) & $-7{,}48\times 10^{-5}$ & $-9{,}06\times 10^{-6}$ & $-1{,}11\times 10^{-6}$\\
  $y' = -2xy$ : Runge-Kutta classique & $1{,}63\times 10^{-6}$ & $1{,}03\times 10^{-7}$ & $6{,}41\times 10^{-9}$\\
  \bottomrule
\end{tabular}
\end{center}
La méthode de Runge-Kutta classique (dite de Kutta-Simpson dans le cours,
$y_{k+1} = y_k + (L_1 + 2L_2 + 2L_3 + L_4)/6$) a bien une erreur divisée par 16 :
c'est elle qu'il faut utiliser quand on veut l'ordre 4 \BF{alg.~5.2, p.~288}.

\Refs \BF{alg.~5.2, p.~288} ; calculs du tableau :
\fichier{tests_algorithmes.py} (fonctions \fichier{runge_cours} et
\texttt{rk4}).

% ------------------------------------------------------------------
\section{Exercice V.3 : algorithme d'Adams-Bashforth d'ordre 2}

\Enonce{Écrire l'algorithme de la méthode d'Adams-Bashforth du second ordre.}

\Pourquoi La formule $y_{k+1} = y_k + \tfrac h2\bigl[3f(x_k,y_k) -
f(x_{k-1},y_{k-1})\bigr]$ réutilise la pente du pas précédent : une seule
nouvelle évaluation de $f$ par pas (au lieu de deux pour RK2). C'est une
méthode \emph{à deux pas} : elle ne peut pas démarrer seule, il faut d'abord
$y_1$, calculé par une méthode à un pas du même ordre (RK2, comme le propose le
cours).

\Etapes
\Etape{1} Démarrage : $y_1$ par RK2 ; garder $f_0 = f(x_0, y_0)$ dans la variable
$f_{\text{prec}}$.

\Etape{2} Pour $k = 1, \dots, n-1$ : $f_k = f(x_k, y_k)$, puis
$y_{k+1} = y_k + \tfrac h2(3f_k - f_{\text{prec}})$.

\Etape{3} Décaler : $f_{\text{prec}} \leftarrow f_k$.

\begin{algo}[ : Adams-Bashforth d'ordre 2]
  \Require $f$, $x_0$, $y_0$, $h$, $n$
  \State $x \leftarrow x_0$ ; $y \leftarrow y_0$
  \State $f_{\text{prec}} \leftarrow f(x, y)$
  \State $y \leftarrow y + \frac h2\bigl(f_{\text{prec}} + f(x + h, y + h f_{\text{prec}})\bigr)$ \Comment{démarrage RK2}
  \State $x \leftarrow x + h$
  \For{$k$ allant de $1$ à $n-1$}
    \State $f_k \leftarrow f(x, y)$
    \State $y \leftarrow y + \frac h2(3f_k - f_{\text{prec}})$
    \State $f_{\text{prec}} \leftarrow f_k$
    \State $x \leftarrow x + h$
    \Print $x$, $y$
  \EndFor
\end{algo}

\Verif Erreurs sur $y(1)$ : $-6{,}15\times 10^{-3}$, $-1{,}54\times 10^{-3}$,
$-3{,}84\times 10^{-4}$ pour $n = 10, 20, 40$ : rapport 4, ordre 2.

\Refs \BF{méthode d'Adams-Bashforth à deux pas, p.~307}.

% ------------------------------------------------------------------
\section{Exercice V.4 : prédicteur-correcteur AB2 -- RK2 et programme}

\Enonce{Établir l'algorithme de la méthode de prédicteur-correcteur, utilisant
comme prédicteur la formule d'Adams-Bashforth du second ordre et comme
correcteur la formule de RK2, et le traduire en Pascal, Fortran ou C.}

\Pourquoi La formule des trapèzes
$y_{k+1} = y_k + \tfrac h2\bigl[f(x_k,y_k) + f(x_{k+1},y_{k+1})\bigr]$ est plus
précise qu'Adams-Bashforth, mais implicite. On \emph{prédit} donc $y_{k+1}$ par
Adams-Bashforth (explicite), puis on \emph{corrige} en mettant cette prévision
dans le second membre de la formule des trapèzes. On gagne en précision pour
deux évaluations de $f$ par pas.

\Etapes
\Etape{1} Démarrage : $y_1$ par RK2.

\Etape{2} Prédiction : $p = y_k + \tfrac h2(3f_k - f_{k-1})$.

\Etape{3} Correction : $y_{k+1} = y_k + \tfrac h2\bigl(f_k + f(x_{k+1}, p)\bigr)$.

\Etape{4} On peut répéter la correction (en remplaçant $p$ par la valeur
corrigée) pour se rapprocher de la solution de la formule implicite ; une seule
correction suffit à garder l'ordre 2.

\begin{algo}[ : prédicteur AB2 -- correcteur RK2]
  \Require $f$, $x_0$, $y_0$, $h$, $n$
  \State $x \leftarrow x_0$ ; $y \leftarrow y_0$ ; $f_{\text{prec}} \leftarrow f(x, y)$
  \State $y \leftarrow y + \frac h2\bigl(f_{\text{prec}} + f(x + h, y + h f_{\text{prec}})\bigr)$ ; $x \leftarrow x + h$
  \For{$k$ allant de $1$ à $n-1$}
    \State $f_k \leftarrow f(x, y)$
    \State $p \leftarrow y + \frac h2(3f_k - f_{\text{prec}})$ \Comment{prédicteur}
    \State $y \leftarrow y + \frac h2\bigl(f_k + f(x + h, p)\bigr)$ \Comment{correcteur}
    \State $f_{\text{prec}} \leftarrow f_k$ ; $x \leftarrow x + h$
    \Print $x$, $y$
  \EndFor
\end{algo}

\programme{C}{en C}{pc_ab2.c}

\Verif Le programme donne $y(1) = 0{,}36961867$ pour $n = 10$, au lieu de
$e^{-1} = 0{,}36787944$ : erreur $1{,}74\times 10^{-3}$, à comparer à
$6{,}15\times 10^{-3}$ pour Adams-Bashforth seul. Avec $n = 20$ et $40$ :
$3{,}66\times 10^{-4}$ et $8{,}38\times 10^{-5}$ (ordre 2).

\Refs \BF{principe des méthodes prédicteur-correcteur, p.~310--311 ;
alg.~5.4, p.~311} (version d'ordre 4).

% ------------------------------------------------------------------
\section{Exercice V.5 : algorithme de la méthode de tir}

\Enonce{Établir l'algorithme de la méthode de tir.}

\Pourquoi Les méthodes des sections précédentes avancent à partir de
conditions \emph{initiales} ; elles ne savent pas imposer la valeur
$y(b) = \eta_2$. On remplace la condition manquante $y'(a)$ par un paramètre
$\gamma$, on résout le problème à valeurs initiales, et on ajuste $\gamma$
pour que $F(\gamma) = y(b, \gamma) - \eta_2$ s'annule. La dérivée
$F'(\gamma)$ n'étant pas connue, le cours la remplace par un quotient de
différences entre deux tirs : c'est la méthode de la sécante. Il faut donc
deux valeurs de départ, $\gamma_0$ et $\gamma_1$.

\Etapes
\Etape{1} Écrire l'équation du second ordre comme un système : $y' = z$,
$z' = f(x, y, z)$, avec $y(a) = \eta_1$, $z(a) = \gamma$.

\Etape{2} Tirer avec $\gamma_0$ et $\gamma_1$ (intégration par Runge-Kutta
d'ordre 4 pour systèmes) et noter $y(b, \gamma_0)$, $y(b, \gamma_1)$.

\Etape{3} Nouvelle valeur par la formule du cours :
$\gamma_{n+1} = \gamma_n - (\gamma_n - \gamma_{n-1})\bigl(y(b,\gamma_n) - \eta_2\bigr)
\big/\bigl(y(b,\gamma_n) - y(b,\gamma_{n-1})\bigr)$.

\Etape{4} Arrêter quand $\abs{y(b,\gamma_n) - \eta_2} < \varepsilon$ ; garder la
dernière trajectoire calculée, qui est la solution approchée.

\begin{algo}[ : tir avec la sécante]
  \Require $f$, $a$, $b$, $\eta_1$, $\eta_2$, $\gamma_0$, $\gamma_1$, $N$ (pas), $\varepsilon$, $k_{\max}$
  \State $Y_0 \leftarrow y(b)$ obtenu par RK4 depuis $y(a) = \eta_1$, $z(a) = \gamma_0$
  \For{$k$ allant de $1$ à $k_{\max}$}
    \State $Y_1 \leftarrow y(b)$ obtenu par RK4 depuis $y(a) = \eta_1$, $z(a) = \gamma_1$
    \If{$\abs{Y_1 - \eta_2} < \varepsilon$}
      \Print $\gamma_1$ et la trajectoire $(x_i, y_i)$ \Stop
    \EndIf
    \State $\gamma \leftarrow \gamma_1 - (\gamma_1 - \gamma_0)(Y_1 - \eta_2)/(Y_1 - Y_0)$
    \State $\gamma_0 \leftarrow \gamma_1$ ; $Y_0 \leftarrow Y_1$ ; $\gamma_1 \leftarrow \gamma$
  \EndFor
  \Print \og pas de convergence \fg{}
\end{algo}
La procédure \og RK4 pour un système \fg{} est détaillée à l'exercice 5 de la fin
du document.

\Verif Pour $y'' = -y$, $y(0) = 0$, $y(\pi/2) = 1$ (solution $\sin x$, donc
$\gamma = 1$), avec $\gamma_0 = 0{,}5$, $\gamma_1 = 2$ : $\gamma = 1{,}0000000003$
après 2 itérations. Pour un problème linéaire, $y(b,\gamma)$ est une fonction
affine de $\gamma$ et la sécante trouve la valeur exacte (aux erreurs de RK4
près) dès la première correction.

\Refs \BF{§~11.2, sécante pour le tir p.~694 ; alg.~11.2, p.~696}
(version avec Newton).

% ------------------------------------------------------------------
\section{Exercice V.6 : discrétisation de l'équation intégro-différentielle}

\Enonce{Discrétiser l'équation intégro-différentielle
$y'(x) = 2y(x) + f(x) + \int_0^x k(x,t)\,y(t)\,\dd t$.}

\Attention Il faut une condition initiale $y(0) = y_0$, que l'énoncé ne donne
pas ; on la suppose connue.

\Pourquoi L'équation mélange une dérivée et une intégrale dont la borne
supérieure est $x$ (type Volterra). Sur une grille $x_i = ih$, l'intégrale au
point $x_i$ ne fait intervenir que $y_0, \dots, y_i$, c'est-à-dire des valeurs
déjà calculées quand on avance de $x_i$ à $x_{i+1}$. On peut donc combiner une
formule d'intégration (trapèzes) et un schéma pour la dérivée (Euler, puis
trapèzes pour monter en ordre).

\Etapes
\Etape{1} Intégrale au nœud $x_i$ par les trapèzes composites :
\[
  I_i = h\Bigl[\tfrac12k(x_i,x_0)y_0 + \sum_{j=1}^{i-1} k(x_i,x_j)y_j
        + \tfrac12k(x_i,x_i)y_i\Bigr], \qquad I_0 = 0 .
\]

\Etape{2} Second membre au nœud $x_i$ : $F_i = 2y_i + f(x_i) + I_i$.

\Etape{3} Schéma d'Euler : $y_{i+1} = y_i + hF_i$, pour $i = 0, 1, \dots, N-1$.
Il est explicite, d'ordre 1.

\Etape{4} Variante d'ordre 2 (trapèzes en $x$) :
$y_{i+1} = y_i + \tfrac h2(F_i + F_{i+1})$. Comme $F_{i+1}$ dépend
\emph{linéairement} de $y_{i+1}$ (par $2y_{i+1}$ et par le dernier terme
$\tfrac h2k(x_{i+1},x_{i+1})y_{i+1}$ de $I_{i+1}$), on isole $y_{i+1}$ :
\[
  y_{i+1}\Bigl[1 - \frac h2\Bigl(2 + \frac h2k(x_{i+1},x_{i+1})\Bigr)\Bigr]
  = y_i + \frac h2\Bigl[F_i + f(x_{i+1})
    + h\Bigl(\tfrac12k(x_{i+1},x_0)y_0 + \sum_{j=1}^{i}k(x_{i+1},x_j)y_j\Bigr)\Bigr].
\]

\begin{algo}[ : schéma d'Euler pour l'équation intégro-différentielle]
  \Require $f$, $k$, $y_0$, $h$, $N$
  \For{$i$ allant de $0$ à $N-1$}
    \State $I \leftarrow 0$
    \If{$i > 0$}
      \State $I \leftarrow \frac12k(x_i,x_0)y_0 + \frac12k(x_i,x_i)y_i$
      \For{$j$ allant de $1$ à $i-1$} \State $I \leftarrow I + k(x_i,x_j)\,y_j$ \EndFor
      \State $I \leftarrow h\,I$
    \EndIf
    \State $y_{i+1} \leftarrow y_i + h\bigl(2y_i + f(x_i) + I\bigr)$
  \EndFor
\end{algo}
(Il faut garder toutes les valeurs $y_j$ en mémoire, car l'intégrale les
utilise toutes.)

\Verif Test construit pour avoir une solution connue : $k = 1$,
$f(x) = 1 - 2e^{x}$, $y(0) = 1$ ; alors $y = e^{x}$ convient (car
$2e^{x} + 1 - 2e^{x} + (e^{x} - 1) = e^{x}$). Erreurs sur $y(1)$ :
\begin{center}
\begin{tabular}{lcccc}
  \toprule
  $N$ & 10 & 20 & 40 & 80\\ \midrule
  Euler & $-2{,}08\times 10^{-1}$ & $-1{,}17\times 10^{-1}$ & $-6{,}22\times 10^{-2}$ & $-3{,}21\times 10^{-2}$\\
  Trapèzes & $5{,}85\times 10^{-3}$ & $1{,}45\times 10^{-3}$ & $3{,}61\times 10^{-4}$ & $9{,}02\times 10^{-5}$\\
  \bottomrule
\end{tabular}
\end{center}
Les rapports successifs (environ 2 pour Euler, 4 pour les trapèzes)
confirment les ordres 1 et 2.

\Refs Formule des trapèzes : \BF{th.~4.5, p.~205} ; principe de l'intégration
pas à pas : chapitre V du cours.

\chapter{Chapitre VI : équations aux dérivées partielles}

\paragraph{Outil commun : l'analyse de stabilité de von Neumann.}
Plusieurs exercices demandent une condition de stabilité. On suit la démarche
du cours. \emph{Pourquoi} : une erreur commise à un instant (arrondi, donnée)
se propage par le schéma lui-même ; si elle est amplifiée à chaque pas, elle
finit par détruire la solution. Pour un schéma linéaire à coefficients
constants, on décompose la perturbation en modes de Fourier
$\varepsilon_{i,j} = \xi^{\,j}e^{Ipih}$ ($I^2 = -1$) ; chaque mode est
multiplié par le \emph{facteur d'amplification} $\xi$ à chaque pas de temps.
\emph{Comment} : on remplace $\varepsilon_{i,j}$ dans le schéma, on simplifie
par $\xi^{\,j}e^{Ipih}$, on obtient $\xi$ en fonction de $\theta = ph$, et le
schéma est stable si $\abs{\xi} \le 1$ pour tous les $\theta$. Les identités
utiles sont
\[
  e^{I\theta} + e^{-I\theta} - 2 = 2\cos\theta - 2 = -4\sin^2\frac{\theta}{2},
  \qquad e^{I\theta} - e^{-I\theta} = 2I\sin\theta .
\]
Le cours utilise $\sin(iph)\,e^{-rjk}$, ce qui revient au même ($e^{-rk} = \xi$).
On note $s = \sin^2(\theta/2) \in [0,1]$.

% ------------------------------------------------------------------
\section{Exercice VI.1 : Gauss-Seidel et Jacobi pour l'équation de Poisson}

\Enonce{Écrire le schéma d'itération de Gauss-Seidel et de Jacobi pour
résoudre l'équation de Poisson.}

\Pourquoi Le système (P2)--(P3) a $(n-1)(m-1)$ inconnues mais chaque équation
n'en contient que cinq. L'élimination de Gauss détruirait cette structure
creuse (remplissage) et coûterait cher ; les méthodes itératives n'utilisent
que les cinq voisins et ne stockent que la grille. C'est la méthode de
Gauss-Seidel que Burden et Faires utilisent pour ce système dans leur
algorithme 12.1 \BF{p.~738--739}.

\Etapes
\Etape{1} On isole $U_{i,j}$ dans (P2). Avec $r = h/k$ :
\[
  U_{i,j} = \frac{U_{i+1,j} + U_{i-1,j} + r^2\bigl(U_{i,j+1} + U_{i,j-1}\bigr) - h^2f_{i,j}}{2(1 + r^2)} .
\]

\Etape{2} Jacobi : tout le second membre est pris à l'itération $q$ ; il faut
une deuxième grille $V$ pour les nouvelles valeurs.

\Etape{3} Gauss-Seidel : on parcourt la grille ($i$ croissant, puis $j$
croissant) et on écrase les valeurs ; $U_{i-1,j}$ et $U_{i,j-1}$ sont alors
déjà à l'itération $q+1$ :
\[
  U^{(q+1)}_{i,j} = \frac{U^{(q)}_{i+1,j} + U^{(q+1)}_{i-1,j}
     + r^2\bigl(U^{(q)}_{i,j+1} + U^{(q+1)}_{i,j-1}\bigr) - h^2f_{i,j}}{2(1 + r^2)} .
\]

\Etape{4} Les valeurs au bord, fixées par (P3), ne sont jamais modifiées. Test
d'arrêt du chapitre I sur la somme des variations.

\begin{algo}[ : Gauss-Seidel pour (P2)--(P3)]
  \Require $a, b, c, d$, $n$, $m$, $f$, $g$, $\varepsilon$, $q_{\max}$
  \State $h \leftarrow (b-a)/n$ ; $k \leftarrow (d-c)/m$ ; $r^2 \leftarrow (h/k)^2$
  \State $U_{i,j} \leftarrow g(x_i, y_j)$ sur le bord ; $U_{i,j} \leftarrow 0$ à l'intérieur
  \For{$q$ allant de $1$ à $q_{\max}$}
    \State $\mathit{num} \leftarrow 0$ ; $\mathit{den} \leftarrow 0$
    \For{$i$ allant de $1$ à $n-1$}
      \For{$j$ allant de $1$ à $m-1$}
        \State $t \leftarrow \bigl[U_{i+1,j} + U_{i-1,j} + r^2(U_{i,j+1} + U_{i,j-1}) - h^2f(x_i,y_j)\bigr]/\bigl(2(1+r^2)\bigr)$
        \State $\mathit{num} \leftarrow \mathit{num} + \abs{t - U_{i,j}}$ ; $\mathit{den} \leftarrow \mathit{den} + \abs{t}$
        \State $U_{i,j} \leftarrow t$
      \EndFor
    \EndFor
    \If{$\mathit{num}/\mathit{den} < \varepsilon$} \Print $q$ et la grille $U$ \Stop \EndIf
  \EndFor
\end{algo}
Pour Jacobi, on écrit $t$ dans une grille $V$ ($V_{i,j} \leftarrow t$), et on
recopie $V$ dans $U$ à la fin de chaque itération.

\Verif Le schéma à cinq points est exact quand les dérivées quatrièmes de la
solution sont nulles \BF{p.~738}. Les algorithmes ont donc été testés sur des
solutions polynomiales : $u = x^2 - y^2$ (Laplace, exercice 17) et
$u = x^2 + y^2$ avec $f = 4$ (exercice 18) ; la solution discrète obtenue
coïncide avec la solution exacte à $10^{-11}$ près.

\Refs \BF{§~12.1, alg.~12.1, p.~738--739}.

% ------------------------------------------------------------------
\section{Exercice VI.2 : algorithme du schéma explicite (chaleur)}

\Enonce{Établir l'algorithme pour le schéma explicite.}

\Pourquoi Le schéma (C2) donne directement chaque $U_{i,j+1}$ à partir de la
ligne de temps $j$ : on avance ligne par ligne. Il faut deux tableaux (ancien et
nouveau temps), sinon $U_{i-1,j}$ serait écrasé par $U_{i-1,j+1}$ avant
d'être utilisé. Le schéma n'est stable que si $\lambda = a^2k/h^2 \le 1/2$
\BF{p.~746--747} : l'algorithme doit le vérifier avant de calculer.

\Etapes
\Etape{1} Lire $l$, $a$, $m$, $k$, le nombre de pas $N$ ; calculer $h = l/m$ et
$\lambda$ ; refuser si $\lambda > 1/2$.

\Etape{2} Condition initiale (C3) : $U_i = f(x_i)$, $U_0 = U_m = 0$.

\Etape{3} Pour chaque pas de temps : appliquer (C2) aux points intérieurs dans
le tableau $V$, imposer les conditions aux limites, recopier $V$ dans $U$.

\begin{algo}[ : schéma explicite]
  \Require $l$, $a$, $f$, $m$, $k$, $N$
  \State $h \leftarrow l/m$ ; $\lambda \leftarrow a^2k/h^2$
  \If{$\lambda > 1/2$} \Print \og schéma instable : diminuer $k$ \fg{} \Stop \EndIf
  \For{$i$ allant de $0$ à $m$} \State $U_i \leftarrow f(ih)$ \EndFor
  \State $U_0 \leftarrow 0$ ; $U_m \leftarrow 0$
  \For{$j$ allant de $1$ à $N$}
    \For{$i$ allant de $1$ à $m-1$}
      \State $V_i \leftarrow (1 - 2\lambda)U_i + \lambda(U_{i+1} + U_{i-1})$
    \EndFor
    \State $V_0 \leftarrow 0$ ; $V_m \leftarrow 0$
    \For{$i$ allant de $0$ à $m$} \State $U_i \leftarrow V_i$ \EndFor
    \Print $jk$, $U_0, \dots, U_m$
  \EndFor
\end{algo}

\Verif Programmé (avec d'autres conditions aux limites) dans
\fichier{ex14_explicite.pas} : exercice 14 de la fin du document.

\Refs \BF{§~12.2, méthode des différences progressives, p.~743--747}.

% ------------------------------------------------------------------
\section{Exercice VI.3 : stabilité de Crank-Nicolson et de Dufort-Frankel}

\Enonce{Établir les conditions de stabilité des schémas de Crank-Nicolson et de
Dufort-Frankel.}

\Question{Crank-Nicolson.}
\Pourquoi Le schéma moyenne les schémas progressif et régressif ; on s'attend à
une stabilité aussi bonne que celle du schéma régressif.

\Etape{1} Avec $\lambda = a^2k/h^2$, le schéma du cours s'écrit
$\varepsilon_{i,j+1} - \varepsilon_{i,j} = \frac\lambda2\bigl[\delta^2\varepsilon_{i,j}
+ \delta^2\varepsilon_{i,j+1}\bigr]$, où
$\delta^2\varepsilon_i = \varepsilon_{i+1} - 2\varepsilon_i + \varepsilon_{i-1}$.

\Etape{2} Pour un mode de Fourier, $\delta^2$ devient la multiplication par
$-4s$. D'où $\xi - 1 = \frac\lambda2(-4s)(1 + \xi)$, soit
\[
  \xi = \frac{1 - 2\lambda s}{1 + 2\lambda s} .
\]

\Etape{3} Pour $\lambda > 0$ et $s \in [0,1]$, le numérateur est toujours de
valeur absolue inférieure ou égale au dénominateur : $\abs{\xi} \le 1$. Le
schéma de Crank-Nicolson est \textbf{inconditionnellement stable}, sans
condition sur $\lambda$.

\Question{Dufort-Frankel.}
\Pourquoi Le schéma de Richardson est instable ; Dufort et Frankel remplacent
$U_{i,j}$ par la moyenne $(U_{i,j+1} + U_{i,j-1})/2$ dans le terme de
diffusion pour le stabiliser. Il y a trois niveaux de temps : on obtient une
équation du second degré en $\xi$.

\Etape{1} Le schéma s'écrit
$\varepsilon_{i,j+1} - \varepsilon_{i,j-1} = 2\lambda\bigl(\varepsilon_{i+1,j}
+ \varepsilon_{i-1,j} - \varepsilon_{i,j+1} - \varepsilon_{i,j-1}\bigr)$.

\Etape{2} Avec $\varepsilon_{i,j} = \xi^{\,j}e^{Ii\theta}$ et division par
$\xi^{\,j-1}e^{Ii\theta}$ :
$\xi^2 - 1 = 2\lambda(2\xi\cos\theta - \xi^2 - 1)$, c'est-à-dire
\[
  (1 + 2\lambda)\,\xi^2 - 4\lambda\cos\theta\,\xi - (1 - 2\lambda) = 0 .
\]

\Etape{3} Racines :
$\xi = \bigl[2\lambda\cos\theta \pm \sqrt{1 - 4\lambda^2\sin^2\theta}\,\bigr]/(1 + 2\lambda)$.
\begin{itemize}
  \item Si $4\lambda^2\sin^2\theta \le 1$, les racines sont réelles et
        $\abs{\xi} \le (2\lambda\abs{\cos\theta} + 1)/(1 + 2\lambda) \le 1$.
  \item Sinon, elles sont complexes conjuguées et
        $\abs{\xi}^2 = \abs{\text{produit des racines}} =
        (2\lambda - 1)/(2\lambda + 1) < 1$.
\end{itemize}
Dans tous les cas $\abs{\xi} \le 1$ : le schéma de Dufort-Frankel est
\textbf{inconditionnellement stable}, comme l'affirme le cours.

\Verif Le calcul numérique de $\max_\theta\abs{\xi}$ pour
$\lambda = 0{,}1$ ; $0{,}5$ ; $1$ ; $5$ ; $50$ donne $1$ dans tous les cas, pour
les deux schémas (le maximum $1$ est atteint pour $\theta = 0$).

\paragraph{Remarque importante (consistance de Dufort-Frankel).} La stabilité ne
suffit pas : l'erreur de troncature du schéma de Dufort-Frankel est
$O(k^2 + h^2 + k^2/h^2)$. Le schéma ne converge vers la solution de
l'équation de la chaleur que si $k/h \to 0$ [DF2]. Le test numérique le
confirme sur $u_t = u_{xx}$, $u(x,0) = \sin\pi x$ (solution
$e^{-\pi^2t}\sin\pi x$), erreur maximale à $t = 0{,}5$ :
\begin{center}
\begin{tabular}{lccc}
  \toprule
  & $h = 0{,}1$ & $h = 0{,}05$ & $h = 0{,}025$\\ \midrule
  $k/h = 0{,}1$ fixé & $2{,}96\times 10^{-3}$ & $3{,}17\times 10^{-3}$ & $3{,}23\times 10^{-3}$\\
  $\lambda = k/h^2 = 1$ fixé & $2{,}96\times 10^{-3}$ & $7{,}87\times 10^{-4}$ & $2{,}00\times 10^{-4}$\\
  \bottomrule
\end{tabular}
\end{center}
Avec $k/h$ constant, l'erreur ne diminue pas ; avec $k = h^2$, elle est
divisée par 4 à chaque fois.

\Refs Crank-Nicolson : \BF{p.~751--752} (inconditionnellement stable, ordre
$O(k^2 + h^2)$). Dufort-Frankel : [VN] et [DF2] (bibliographie).

% ------------------------------------------------------------------
\section{Exercice VI.4 : algorithmes régressif, Crank-Nicolson, Dufort-Frankel}

\Enonce{Établir les algorithmes pour les méthodes des différences régressives,
de Crank-Nicolson et de Dufort-Frankel.}

\Pourquoi Les deux premiers schémas sont implicites : à chaque pas, il faut
résoudre un système linéaire. Ce système est \emph{tridiagonal} (chaque
équation ne lie que $U_{i-1}$, $U_i$, $U_{i+1}$) et à diagonale strictement
dominante ; on le résout par la factorisation de Crout pour matrices
tridiagonales, en $O(m)$ opérations au lieu de $O(m^3)$. Comme la matrice ne
change pas d'un pas à l'autre, on pourrait même la factoriser une seule fois.
Le troisième schéma est explicite mais à trois niveaux : il lui faut
$U^{(0)}$ et $U^{(1)}$ pour démarrer.

\Question{Résolution d'un système tridiagonal (Crout).}
Système $\ell_iU_{i-1} + d_iU_i + u_iU_{i+1} = r_i$, $i = 1, \dots, M$
($\ell_1 = u_M = 0$). On factorise en une matrice bidiagonale inférieure (de
diagonale $L_i$) fois une bidiagonale supérieure à diagonale unité (de
sur-diagonale $S_i$), puis descente et remontée.

\begin{algo}[ : Crout tridiagonal]
  \State $L_1 \leftarrow d_1$ ; $S_1 \leftarrow u_1/L_1$ ; $z_1 \leftarrow r_1/L_1$
  \For{$i$ allant de $2$ à $M$}
    \State $L_i \leftarrow d_i - \ell_i\,S_{i-1}$
    \If{$i < M$} \State $S_i \leftarrow u_i/L_i$ \EndIf
    \State $z_i \leftarrow (r_i - \ell_i\,z_{i-1})/L_i$
  \EndFor
  \State $U_M \leftarrow z_M$
  \For{$i$ allant de $M-1$ à $1$} \State $U_i \leftarrow z_i - S_i\,U_{i+1}$ \EndFor
\end{algo}

\Question{Différences régressives.}
\Etape{1} (C5) aux points $i = 1, \dots, m-1$ : $\ell_i = u_i = -\lambda$,
$d_i = 1 + 2\lambda$, $r_i = U_i^{(j-1)}$ (les termes de bord sont nuls ici
car $U_0 = U_m = 0$).
\Etape{2} À chaque pas : résoudre le système tridiagonal, la solution devient
la nouvelle ligne.

\Question{Crank-Nicolson.}
\Etape{1} En multipliant le schéma du cours par $k$ :
$(1 + \lambda)U_i^{(j+1)} - \frac\lambda2\bigl(U_{i+1}^{(j+1)} + U_{i-1}^{(j+1)}\bigr)
 = (1 - \lambda)U_i^{(j)} + \frac\lambda2\bigl(U_{i+1}^{(j)} + U_{i-1}^{(j)}\bigr)$.
\Etape{2} Donc $\ell_i = u_i = -\lambda/2$, $d_i = 1 + \lambda$, et $r_i$ égal au
second membre calculé avec l'ancienne ligne.

\Question{Dufort-Frankel.}
\Etape{1} En multipliant le schéma du cours par $2k$ et en regroupant les
termes en $U_i^{(j+1)}$ :
$U_i^{(j+1)} = \bigl[(1 - 2\lambda)U_i^{(j-1)} + 2\lambda\bigl(U_{i+1}^{(j)} + U_{i-1}^{(j)}\bigr)\bigr]
\big/(1 + 2\lambda)$.
\Etape{2} Démarrage : $U^{(0)} = f$, puis $U^{(1)}$ par un pas de Crank-Nicolson
(un schéma à un pas d'ordre 2).
\Etape{3} Trois tableaux : ancien ($j-1$), courant ($j$) et nouveau ($j+1$).

\begin{algo}[ : différences régressives et Crank-Nicolson]
  \Require $l$, $a$, $f$, $m$, $k$, $N$, choix du schéma
  \State $h \leftarrow l/m$ ; $\lambda \leftarrow a^2k/h^2$
  \For{$i$ allant de $0$ à $m$} \State $U_i \leftarrow f(ih)$ \EndFor
  \State $U_0 \leftarrow 0$ ; $U_m \leftarrow 0$
  \For{$j$ allant de $1$ à $N$}
    \For{$i$ allant de $1$ à $m-1$}
      \If{régressif}
        \State $\ell_i \leftarrow -\lambda$ ; $d_i \leftarrow 1 + 2\lambda$ ; $u_i \leftarrow -\lambda$ ; $r_i \leftarrow U_i$
      \Else \Comment{Crank-Nicolson}
        \State $\ell_i \leftarrow -\lambda/2$ ; $d_i \leftarrow 1 + \lambda$ ; $u_i \leftarrow -\lambda/2$
        \State $r_i \leftarrow (1-\lambda)U_i + \frac\lambda2(U_{i+1} + U_{i-1})$
      \EndIf
    \EndFor
    \State résoudre le système tridiagonal (Crout) ; ranger la solution dans $U_1, \dots, U_{m-1}$
    \Print $jk$, $U$
  \EndFor
\end{algo}

\begin{algo}[ : Dufort-Frankel]
  \State $U^{\text{anc}}_i \leftarrow f(ih)$ ; $U^{\text{cour}} \leftarrow$ un pas de Crank-Nicolson depuis $U^{\text{anc}}$
  \For{$j$ allant de $2$ à $N$}
    \For{$i$ allant de $1$ à $m-1$}
      \State $U^{\text{nouv}}_i \leftarrow \bigl[(1 - 2\lambda)U^{\text{anc}}_i
         + 2\lambda(U^{\text{cour}}_{i+1} + U^{\text{cour}}_{i-1})\bigr]/(1 + 2\lambda)$
    \EndFor
    \State $U^{\text{nouv}}_0 \leftarrow 0$ ; $U^{\text{nouv}}_m \leftarrow 0$
    \State $U^{\text{anc}} \leftarrow U^{\text{cour}}$ ; $U^{\text{cour}} \leftarrow U^{\text{nouv}}$
  \EndFor
\end{algo}

\Verif Sur $u_t = u_{xx}$, $u(x,0) = \sin\pi x$, erreur maximale à $t = 0{,}5$
pour $(h, k) = (0{,}1 ; 0{,}01)$, $(0{,}05 ; 0{,}005)$, $(0{,}025 ; 0{,}0025)$
(donc $\lambda = 1, 2, 4$, valeurs pour lesquelles le schéma explicite serait
instable) : régressif $2{,}19\times 10^{-3}$ ; $9{,}78\times 10^{-4}$ ;
$4{,}63\times 10^{-4}$ (ordre 1 en $k$) ; Crank-Nicolson $2{,}68\times 10^{-4}$ ;
$6{,}61\times 10^{-5}$ ; $1{,}65\times 10^{-5}$ (ordre 2). Aucun des deux ne
diverge. Pour Dufort-Frankel, voir la remarque de l'exercice VI.3.

\Refs \BF{alg.~6.7, p.~427 (Crout tridiagonal) ; alg.~12.2, p.~749 (régressif) ;
alg.~12.3, p.~753 (Crank-Nicolson)}.

% ------------------------------------------------------------------
\section{Exercice VI.5 : équation d'ondes, stabilité et algorithme}

\Enonce{a) Établir la condition de stabilité du schéma explicite (O2).
b) Établir l'algorithme de ce schéma.}

\Question{a) Stabilité.}
\Pourquoi Le schéma (O2) a trois niveaux de temps : comme pour Dufort-Frankel,
le facteur d'amplification est racine d'une équation du second degré.

\Etape{1} Avec $\varepsilon_{i,j} = \xi^{\,j}e^{Ii\theta}$ et $\lambda = ak/h$,
(O2) donne
$\xi = 2(1 - \lambda^2) + \lambda^2(2\cos\theta) - \xi^{-1}$, soit
\[
  \xi^2 - 2\bigl(1 - 2\lambda^2s\bigr)\xi + 1 = 0, \qquad s = \sin^2\frac\theta2 .
\]

\Etape{2} Le produit des racines vaut $1$. Si les racines sont réelles et
distinctes, l'une est de module $> 1$ : instabilité. Il faut donc des racines
complexes conjuguées (ou une racine double), c'est-à-dire un discriminant
négatif ou nul : $\bigl(1 - 2\lambda^2s\bigr)^2 \le 1$, soit
$0 \le \lambda^2s \le 1$. Alors $\abs{\xi} = 1$.

\Etape{3} Cette condition doit être vraie pour tous les modes, donc pour
$s = 1$ : $\lambda^2 \le 1$, c'est-à-dire
\[
  \lambda = \frac{ak}{h} \le 1 .
\]
C'est la condition du cours (condition de Courant-Friedrichs-Lewy) ; Burden et
Faires donnent $\lambda \le 1$ \BF{p.~762}. Le cas limite $\lambda = 1$
est admis par le livre ; le cours écrit $\lambda < 1$.

\Verif $\max_\theta\abs{\xi} = 1$ pour $\lambda = 0{,}9$ et $\lambda = 1$, mais
$1{,}33$ pour $\lambda = 1{,}01$.

\Question{b) Algorithme.}
\Pourquoi Pour calculer la ligne $j+1$, il faut les lignes $j$ et $j-1$ : on
initialise avec (O4) et (O6), puis on applique (O2). La formule (O6) est
utilisée plutôt que (O5) pour que le démarrage soit du même ordre ($k^2$) que le
schéma.

\begin{algo}[ : équation d'ondes, schéma explicite]
  \Require $l$, $a$, $f$, $g$, $m$, $k$, $N$
  \State $h \leftarrow l/m$ ; $\lambda \leftarrow ak/h$
  \If{$\lambda > 1$} \Print \og schéma instable \fg{} \Stop \EndIf
  \For{$i$ allant de $0$ à $m$} \State $U^{\text{anc}}_i \leftarrow f(ih)$ \EndFor
  \State $U^{\text{anc}}_0 \leftarrow 0$ ; $U^{\text{anc}}_m \leftarrow 0$ \Comment{(O3), (O4)}
  \For{$i$ allant de $1$ à $m-1$}
    \State $U^{\text{cour}}_i \leftarrow (1 - \lambda^2)U^{\text{anc}}_i
      + \frac{\lambda^2}{2}(U^{\text{anc}}_{i+1} + U^{\text{anc}}_{i-1}) + k\,g(ih)$ \Comment{(O6)}
  \EndFor
  \State $U^{\text{cour}}_0 \leftarrow 0$ ; $U^{\text{cour}}_m \leftarrow 0$
  \For{$j$ allant de $1$ à $N-1$}
    \For{$i$ allant de $1$ à $m-1$}
      \State $U^{\text{nouv}}_i \leftarrow 2(1 - \lambda^2)U^{\text{cour}}_i
        + \lambda^2(U^{\text{cour}}_{i+1} + U^{\text{cour}}_{i-1}) - U^{\text{anc}}_i$ \Comment{(O2)}
    \EndFor
    \State $U^{\text{nouv}}_0 \leftarrow 0$ ; $U^{\text{nouv}}_m \leftarrow 0$
    \State $U^{\text{anc}} \leftarrow U^{\text{cour}}$ ; $U^{\text{cour}} \leftarrow U^{\text{nouv}}$
    \Print $(j+1)k$, $U^{\text{cour}}$
  \EndFor
\end{algo}

\Verif Ce schéma est le cas $a = d = 0$ du programme \fichier{ex16_onde.c}
(exercice 16 de la fin du document). Pour $f = \sin\pi x$, $g = 0$, $l = 1$,
vitesse $1$ (solution exacte $\sin\pi x\cos\pi t$), $h = 0{,}02$, $k = 0{,}01$
($\lambda = 0{,}5$) : en $x = 0{,}5$, $t = 1{,}5$, on obtient $-5{,}8\times
10^{-4}$ pour une valeur exacte nulle.

\Refs \BF{(12.19)--(12.25), p.~758--760 ; alg.~12.4, p.~761 ; stabilité
p.~762}.

\chapter{Exercices de fin de document (1 à 7)}

% ------------------------------------------------------------------
\section{Exercice 1 : maximum de la loi de Planck}

\Question{1- Équation vérifiée par $x_m$.}
\Pourquoi Avec $x = hc/(\lambda kT)$, on a $\lambda = hc/(xkT)$, et l'intensité
devient une fonction de $x$ seul (à un facteur constant près). Comme
$x \mapsto \lambda$ est strictement monotone ($\dd x/\dd\lambda =
-hc/(\lambda^2kT) \ne 0$), chercher le maximum en $\lambda$ revient à chercher
le maximum en $x$.

\Etape{1} $I = 8\pi hc\,\lambda^{-5}/(e^{x} - 1) = 8\pi hc\,(kT/hc)^5\;x^5/(e^{x} - 1)$.

\Etape{2} On dérive $\varphi(x) = x^5/(e^{x} - 1)$ :
\[
  \varphi'(x) = \frac{5x^4(e^{x} - 1) - x^5e^{x}}{(e^{x} - 1)^2} = 0
  \iff 5(e^{x} - 1) = xe^{x} .
\]

\Etape{3} L'équation de $x_m$ est donc
\[
  f(x) = (x - 5)e^{x} + 5 = 0, \qquad\text{ou encore}\qquad x = 5\bigl(1 - e^{-x}\bigr) .
\]
Attention : $x = 0$ est solution évidente ($f(0) = -5 + 5 = 0$) mais
correspond à $\lambda \to \infty$, pas au maximum ; on cherche la racine
strictement positive.

\Question{2- Algorithme de Newton-Raphson.}
\Pourquoi $f$ est dérivable, de dérivée simple $f'(x) = (x - 4)e^{x}$ : Newton
converge vite si l'on part près de la racine. Comme $e^{-x}$ est petit pour $x$
grand, l'écriture $x = 5(1 - e^{-x})$ suggère $x_0 = 5$. Sur $[4{,}9 ; 5]$,
$f(4{,}9) < 0 < f(5)$, $f' > 0$ et $f'' = (x - 3)e^{x} > 0$, donc
$f(x_0)f''(x_0) > 0$ en $x_0 = 5$ : d'après la remarque du cours, la suite est
monotone et converge.

\Etapes
\Etape{1} $x_0 = 5$, $m = \min\abs{f'} = f'(4{,}9) = 120{,}86$ sur $[4{,}9 ; 5]$.

\Etape{2} $x_{n+1} = x_n - \dfrac{(x_n - 5)e^{x_n} + 5}{(x_n - 4)e^{x_n}}$.

\Etape{3} Arrêt quand $\abs{f(x_n)}/m < E$ (algorithme de l'exercice II.1).

\begin{algo}[ : Newton pour $x_m$]
  \Require $E$, $k_{\max}$
  \State $x \leftarrow 5$ ; $m \leftarrow (4{,}9 - 4)e^{4{,}9}$ ; $k \leftarrow 0$
  \Repeat
    \State $x \leftarrow x - \bigl((x-5)e^{x} + 5\bigr)/\bigl((x-4)e^{x}\bigr)$ ; $k \leftarrow k + 1$
  \Until{$\abs{(x-5)e^{x} + 5}/m < E$ \textbf{ou} $k \ge k_{\max}$}
  \Print $x_m = x$ et $\lambda_mT = hc/(k\,x_m)$
\end{algo}

\Verif Itérés : $x_1 = 4{,}966310265$, $x_2 = 4{,}965115686$,
$x_3 = 4{,}965114231746$, $x_4 = 4{,}965114231744$ (convergence quadratique :
le nombre de chiffres exacts double environ à chaque pas). Cette valeur est
celle de la suite A094090 de l'OEIS ($4{,}96511423174427\ldots$) [OEIS]. Avec
les valeurs exactes des constantes ($h = 6{,}62607015\times 10^{-34}$~J\,s,
$c = 299\,792\,458$~m/s, $k = 1{,}380649\times 10^{-23}$~J/K), on obtient la loi
de Wien $\lambda_mT = hc/(kx_m) = 2{,}8978\times 10^{-3}$~m\,K.

\Refs \BF{alg.~2.3, p.~67} ; [OEIS].

% ------------------------------------------------------------------
\section{Exercice 2 : volume d'un gaz de Van der Waals}

\Pourquoi Pour $T$ et $P$ donnés, $V$ est racine de
$g(V) = (P + a/V^2)(V - b) - RMT$ ; on ne sait pas l'isoler, il faut une
méthode numérique : Newton. En multipliant par $V^2 > 0$, on obtient une
équation polynomiale de degré 3, sans division par $V$ :
\[
  f(V) = PV^3 - (Pb + K)V^2 + aV - ab = 0, \qquad K = RMT .
\]
Elle se dérive facilement, $f'(V) = 3PV^2 - 2(Pb + K)V + a$, et s'évalue par
Horner (chapitre II). Le point de départ naturel est le volume du gaz parfait,
$V_0 = K/P$ (on retrouve Van der Waals quand $a = b = 0$). Seules les racines
$V > b$ ont un sens physique. Sous la température critique, le polynôme peut
avoir trois racines réelles (liquide et gaz) ; partir de $V_0$ oriente vers la
racine \og gaz \fg.

\Etapes
\Etape{1} Lire $a$, $b$, $R$, $M$ puis, pour chaque couple $(T, P)$ :
$K = RMT$, $V = K/P$.

\Etape{2} Itérer $V \leftarrow V - f(V)/f'(V)$ jusqu'à
$\abs{\Delta V} < \varepsilon\abs{V}$ (test relatif, $V$ pouvant être très
petit ou très grand selon les unités).

\Etape{3} Vérifier $V > b$ et afficher.

\begin{algo}[ : Van der Waals]
  \Require $a$, $b$, $R$, $M$, $\varepsilon$, $k_{\max}$
  \While{il reste des couples $(T, P)$ à traiter}
    \Read $T$, $P$
    \State $K \leftarrow RMT$ ; $V \leftarrow K/P$ ; $k \leftarrow 0$
    \Repeat
      \State $f \leftarrow ((PV - (Pb + K))V + a)V - ab$ \Comment{Horner}
      \State $f' \leftarrow (3PV - 2(Pb + K))V + a$
      \State $\Delta \leftarrow f/f'$ ; $V \leftarrow V - \Delta$ ; $k \leftarrow k + 1$
    \Until{$\abs{\Delta} < \varepsilon\abs{V}$ \textbf{ou} $k \ge k_{\max}$}
    \If{$V \le b$ \textbf{ou} $k \ge k_{\max}$} \Print \og échec \fg{} \Else \ \Print $T$, $P$, $V$ \EndIf
  \EndWhile
\end{algo}

\Verif Pour le dioxyde de carbone, avec $a = 3{,}592$~L$^2$\,atm/mol$^2$,
$b = 0{,}04267$~L/mol [VdW] et $R = 0{,}08206$~L\,atm/(mol\,K), pour une mole
($K = RT$) à $T = 300$~K : à $P = 10$~atm, $V = 2{,}354688$~L en 5 itérations
(gaz parfait : $2{,}4618$~L) ; à $P = 50$~atm, $V = 0{,}358484$~L en 7
itérations (gaz parfait : $0{,}4924$~L). Dans les deux cas les deux autres
racines du polynôme sont complexes, la racine trouvée est donc la seule racine
réelle.

\Refs \BF{alg.~2.3, p.~67 ; Horner, alg.~2.7, p.~94} ; constantes : [VdW].

% ------------------------------------------------------------------
\section{Exercice 3 : facteur de compressibilité (De Santis)}

\Attention L'expression $z = (1 + y + y^2 - y^3)/(1 - y)^3$, avec
$y = b/(4v)$, est celle de Carnahan et Starling pour un gaz de sphères dures,
reprise dans l'équation d'état de Carnahan-Starling-De Santis (1976) [CSD].
Dans le document scanné, le dénominateur est imprimé $(1 - y^3)$ ; c'est
$(1 - y)^3$ qui est correct (voir la transcription du cours). On connaît $z$
et $v$, on cherche $b$.

\Question{1- Principe de la méthode de Newton.}
\Pourquoi On remplace la courbe de $f$ par sa tangente au point $x_n$ ; le point
où la tangente coupe l'axe donne $x_{n+1} = x_n - f(x_n)/f'(x_n)$ (chapitre II).
La convergence est quadratique près d'une racine simple.

\Question{2- Algorithme.}
\Pourquoi On cherche d'abord $y$, puis $b = 4vy$. On multiplie l'équation par
$(1 - y)^3$ pour supprimer le pôle $y = 1$ :
\[
  f(y) = z(1 - y)^3 - (1 + y + y^2 - y^3) = 0, \qquad
  f'(y) = -3z(1 - y)^2 - (1 + 2y - 3y^2) .
\]
Pour le départ : quand $y$ est petit, $(1 - y)^{-3} = 1 + 3y + 6y^2 + \dots$,
d'où $z = 1 + 4y + 10y^2 + \dots$ ; on peut donc prendre $y_0 = (z - 1)/4$.
Physiquement $0 < y < 1$.

\Etapes
\Etape{1} Lire $z$, $v$, $\varepsilon$, $k_{\max}$ ; $y \leftarrow (z - 1)/4$
(ou une valeur donnée).

\Etape{2} Itérer $y \leftarrow y - f(y)/f'(y)$ jusqu'à
$\abs{\Delta y} < \varepsilon$.

\Etape{3} $b = 4vy$.

\begin{algo}[ : De Santis]
  \Read $z$, $v$, $y_0$, $\varepsilon$, $k_{\max}$
  \State $y \leftarrow y_0$
  \For{$k$ allant de $1$ à $k_{\max}$}
    \State $f \leftarrow z(1-y)^3 - (1 + y + y^2 - y^3)$
    \State $f' \leftarrow -3z(1-y)^2 - (1 + 2y - 3y^2)$
    \State $\Delta \leftarrow f/f'$ ; $y \leftarrow y - \Delta$
    \If{$\abs{\Delta} < \varepsilon$} \Print $y$, $b = 4vy$ \Stop \EndIf
  \EndFor
  \Print \og pas de convergence \fg{}
\end{algo}

\Question{3- Programme.}
\programme{desantis.py}

\Verif Pour $y = 0{,}3$, la formule donne $z = 1{,}363/0{,}343 = 3{,}9737609$.
Avec cette valeur de $z$, $v = 1$ et $y_0 = 0{,}5$, le programme retrouve en 6
itérations $y = 0{,}2999999984$ et $b = 1{,}1999999937$ (la valeur de $z$ n'ayant que 8
chiffres, on ne peut pas espérer mieux) ; avec le départ $y_0 = (z-1)/4 =
0{,}7434$, la convergence vers $0{,}3$ demande 8 itérations.

\Refs \BF{alg.~2.3, p.~67} ; [CSD].

% ------------------------------------------------------------------
\section{\texorpdfstring{Exercice 4 : $y' = -ay - by^4 + c$ par prédicteur-correcteur}{Exercice 4 : y' = -ay - by\textasciicircum 4 + c par prédicteur-correcteur}}

\Question{1)- Schéma itératif.}
\Pourquoi On applique l'exercice V.4 avec $f(x, y) = -ay - by^4 + c$ (qui ne
dépend pas de $x$). Le prédicteur AB2 est explicite mais il lui faut deux
valeurs ; le démarrage se fait par RK2.

\Etapes
\Etape{1} Démarrage : $y_1 = y_0 + \frac h2\bigl[f(y_0) + f\bigl(y_0 + hf(y_0)\bigr)\bigr]$.

\Etape{2} Pour $k \ge 1$, prédiction
$p_{k+1} = y_k + \frac h2\bigl[3f(y_k) - f(y_{k-1})\bigr]$.

\Etape{3} Correction $y_{k+1} = y_k + \frac h2\bigl[f(y_k) + f(p_{k+1})\bigr]$,
avec $f(y) = -ay - by^4 + c$.

\Question{2)- Algorithme.} C'est celui de l'exercice V.4, avec cette fonction
$f$ ; on garde $f(y_{k-1})$ dans une variable pour ne calculer que deux
valeurs de $f$ par pas.

\Question{3)- Programme.}
\programme{ex4_pc.py}

\Verif Pour $a = 1$, $b = 0{,}5$, $c = 2$, $y_0 = 0$, $x = 2$ et $n = 40$, les
programme donne $y(2) = 1{,}1425354104$ ; une intégration de référence
très précise (SciPy) donne $1{,}1425030979$ : erreur $3\times 10^{-5}$. La
solution tend vers l'état d'équilibre $y^*$, racine de $-y - 0{,}5y^4 + 2 = 0$,
soit $y^* = 1{,}1439$.

\Refs \BF{p.~307 (AB2) ; p.~310--311 (prédicteur-correcteur)}.

% ------------------------------------------------------------------
\section{\texorpdfstring{Exercice 5 : $y'' + \frac{a}{x}y' + be^{cy} = 0$ par la méthode de tir}{Exercice 5 : y'' + (a/x)y' + b exp(cy) = 0 par la méthode de tir}}

\Pourquoi La condition $y'(0) = 0$ est en $x = 0$ et la condition $y(1) = 1$
en $x = 1$ : il manque la valeur $y(0)$ pour pouvoir partir de $x = 0$. On la
remplace par le paramètre de tir $\gamma = y(0)$ et on cherche $\gamma$ tel
que $F(\gamma) = y(1, \gamma) - 1 = 0$. Newton demande
$F'(\gamma) = \partial y(1,\gamma)/\partial\gamma$ : on l'obtient en dérivant
l'équation par rapport à $\gamma$, ce qui donne une deuxième équation
différentielle (méthode de Burden et Faires, alg.~11.2).

\Question{1)- Schéma itératif de Newton pour $\gamma$.}
\Etape{1} Posons $z(x) = \partial y(x,\gamma)/\partial\gamma$. En dérivant
l'équation et les conditions initiales $y(0) = \gamma$, $y'(0) = 0$ par
rapport à $\gamma$ :
\[
  z'' + \frac ax z' + bc\,e^{cy}z = 0, \qquad z(0) = 1, \quad z'(0) = 0 .
\]

\Etape{2} Newton : $\displaystyle \gamma_{n+1} = \gamma_n - \frac{y(1, \gamma_n) - 1}{z(1, \gamma_n)}$.

\Etape{3} À chaque itération, on intègre ensemble les équations en $y$ et en
$z$ (même $\gamma_n$), puis on corrige $\gamma$.

\Question{Le point singulier $x = 0$.}
Le terme $\frac ax y'$ n'est pas défini en $x = 0$. Mais $y'(0) = 0$, donc
$y'(x)/x \to y''(0)$ quand $x \to 0$. L'équation en $x = 0$ devient
$y''(0) + a\,y''(0) + be^{c\gamma} = 0$, d'où
\[
  y''(0) = -\frac{b\,e^{c\gamma}}{1 + a}, \qquad\text{et de même}\qquad
  z''(0) = -\frac{bc\,e^{c\gamma}}{1 + a}.
\]
Ce sont les valeurs à utiliser au premier pas.

\Question{2)- Méthode de Runge-Kutta d'ordre 4 pour un système.}
On pose $Y = (Y_1, Y_2, Y_3, Y_4) = (y, y', z, z')$ ; le système s'écrit
$Y' = F(x, Y)$ avec
\[
  F(x, Y) = \Bigl(Y_2,\; -\frac ax Y_2 - be^{cY_1},\; Y_4,\;
  -\frac ax Y_4 - bc\,e^{cY_1}Y_3\Bigr)
\]
(et les limites ci-dessus pour $x = 0$). RK4 s'applique composante par
composante, avec les mêmes formules que pour une équation scalaire :
\[
  K_1 = hF(x_i, Y_i), \qquad K_2 = hF\bigl(x_i + \tfrac h2, Y_i + \tfrac12K_1\bigr),
\]
\[
  K_3 = hF\bigl(x_i + \tfrac h2, Y_i + \tfrac12K_2\bigr), \qquad
  K_4 = hF(x_i + h, Y_i + K_3),
\]
\[
  Y_{i+1} = Y_i + \tfrac16\bigl(K_1 + 2K_2 + 2K_3 + K_4\bigr).
\]
Il est essentiel de calculer \emph{toutes} les composantes de $K_1$ avant de
passer à $K_2$, etc.

\Question{3)- Algorithme complet.}
\begin{algo}[ : tir avec Newton]
  \Require $a$, $b$, $c$, $N$, $\gamma_0$, $\varepsilon$, $k_{\max}$
  \State $h \leftarrow 1/N$ ; $\gamma \leftarrow \gamma_0$
  \For{$k$ allant de $1$ à $k_{\max}$}
    \State $x \leftarrow 0$ ; $Y \leftarrow (\gamma, 0, 1, 0)$ \Comment{conditions initiales}
    \For{$i$ allant de $1$ à $N$}
      \State $K_1 \leftarrow hF(x, Y)$ ; $K_2 \leftarrow hF(x + h/2, Y + K_1/2)$
      \State $K_3 \leftarrow hF(x + h/2, Y + K_2/2)$ ; $K_4 \leftarrow hF(x + h, Y + K_3)$
      \State $Y \leftarrow Y + (K_1 + 2K_2 + 2K_3 + K_4)/6$ ; $x \leftarrow x + h$
    \EndFor
    \State $\Delta \leftarrow (Y_1 - 1)/Y_3$ \Comment{$Y_1 = y(1)$, $Y_3 = z(1)$}
    \State $\gamma \leftarrow \gamma - \Delta$
    \If{$\abs{\Delta} < \varepsilon$} \Print $\gamma$ (et refaire un tir pour afficher $y(x_i)$) \Stop \EndIf
  \EndFor
  \Print \og pas de convergence \fg{}
\end{algo}
où $F(0, Y)$ utilise les limites $y''(0)$ et $z''(0)$.

\Verif
\begin{itemize}
  \item Cas test exact $c = 0$ : l'équation devient linéaire,
        $y = \gamma - \frac{b}{2(1+a)}x^2$, d'où $\gamma = 1 + \frac{b}{2(1+a)}$.
        Pour $a = b = 1$ : $\gamma = 1{,}25$. L'algorithme ($N = 200$) trouve
        $1{,}2500000000$ en 2 itérations.
  \item $a = 1$, $b = 1$, $c = 0{,}5$ : $\gamma = 1{,}4961487192$ (5 itérations
        depuis $\gamma_0 = 1$), avec $y(1) = 1$ à $10^{-12}$ près. Ce problème a
        une seconde solution, $\gamma = 9{,}5942$, que Newton atteint si l'on part
        plus haut : le point de départ choisit la solution.
  \item $a = 2$, $b = 0{,}5$, $c = 1$ : $\gamma = 1{,}2739793373$.
  \item Pour $a = b = c = 1$, il n'y a \emph{pas} de solution : le calcul de
        $y(1, \gamma)$ pour $\gamma \in [-5, 5]$ montre que son maximum vaut
        $0{,}693$ (vers $\gamma \approx 2{,}1$), toujours inférieur à $1$.
        Une méthode de tir qui ne converge pas peut donc signaler un problème
        sans solution, et pas seulement un mauvais $\gamma_0$.
\end{itemize}

\Refs \BF{§~11.2, p.~693--698, alg.~11.2, p.~696 ; RK4 pour les systèmes, alg.~5.7,
p.~333}.

% ------------------------------------------------------------------
\section{\texorpdfstring{Exercice 6 : potentiel entre cathode et anode, $V'' = aV^{-1/2}$}{Exercice 6 : potentiel entre cathode et anode, V'' = a V\textasciicircum (-1/2)}}

\Question{1-a- Itérations RK2.}
\Pourquoi On intègre $y' = f(x,y)$ entre $x_k$ et $x_{k+1}$ par les trapèzes et on
remplace l'inconnue $y_{k+1}$ du second membre par la prévision d'Euler
(exercice V.1) :
\[
  y_{k+1} = y_k + \frac h2\Bigl[f(x_k, y_k) + f\bigl(x_k + h, y_k + hf(x_k, y_k)\bigr)\Bigr].
\]

\Question{1-b- Algorithme pour (1).}
\Pourquoi On pose $W = V'$ (c'est l'opposé du champ électrique) : (1) devient le
système $V' = W$, $W' = aV^{-1/2}$, avec $V(0) = W(0) = 0$. RK2 s'applique au
vecteur $(V, W)$. Mais en $x = 0$, $V = 0$ et le second membre $a/\sqrt{V}$
est infini : on ne peut pas démarrer en $0$. On cherche alors une solution de
la forme $V = Cx^{p}$ : $V'' = Cp(p-1)x^{p-2}$ et $aV^{-1/2} = aC^{-1/2}x^{-p/2}$
imposent $p - 2 = -p/2$, soit $p = 4/3$, puis $C\cdot\frac43\cdot\frac13 =
aC^{-1/2}$, soit $C = (9a/4)^{2/3}$. On retrouve la loi de Child-Langmuir
$V = (9a/4)^{2/3}x^{4/3}$ [CL]. On démarre donc en $x = \varepsilon$ petit, avec
$V(\varepsilon) = C\varepsilon^{4/3}$ et
$W(\varepsilon) = \frac43C\varepsilon^{1/3}$.

\begin{algo}[ : RK2 pour $V'' = aV^{-1/2}$]
  \Require $a$, $l$, $\varepsilon$, $n$
  \State $C \leftarrow (9a/4)^{2/3}$ ; $x \leftarrow \varepsilon$ ; $V \leftarrow C\varepsilon^{4/3}$ ;
         $W \leftarrow \frac43C\varepsilon^{1/3}$ ; $h \leftarrow (l - \varepsilon)/n$
  \For{$i$ allant de $1$ à $n$}
    \State $k_{1V} \leftarrow hW$ ; $k_{1W} \leftarrow ha/\sqrt{V}$
    \State $k_{2V} \leftarrow h(W + k_{1W})$ ; $k_{2W} \leftarrow ha/\sqrt{V + k_{1V}}$
    \State $V \leftarrow V + (k_{1V} + k_{2V})/2$ ; $W \leftarrow W + (k_{1W} + k_{2W})/2$ ; $x \leftarrow x + h$
    \Print $x$, $V$, $W$
  \EndFor
\end{algo}

\Question{1-c- Programme.}
Le même fichier contient aussi la fonction \texttt{tir} de la question 2.
\programme{ex6_rk2.py}

\Verif $a = 1$, $l = 1$, $\varepsilon = 10^{-2}$ ; valeur exacte
$V(1) = (9/4)^{2/3} = 1{,}71707136$. Résultats : $n = 100$ : $1{,}72325821$ ;
$n = 200$ : $1{,}71871098$ ; $n = 400$ : $1{,}71749474$.
Les erreurs ($6{,}2\times 10^{-3}$ ; $1{,}6\times 10^{-3}$ ;
$4{,}2\times 10^{-4}$) sont divisées par 4 : ordre 2. Partir trop près de la
singularité dégrade la précision : avec $\varepsilon = 10^{-3}$ et $n = 400$,
on obtient $1{,}7256$.

\Question{2-a- Principe de la méthode de tir.}
\Pourquoi Les conditions sont maintenant $V(0) = 0$ et $V(l) = V_0$ : une
condition à chaque bout. On choisit la pente inconnue $\gamma = V'(0)$
(c'est-à-dire le champ sur la cathode), on intègre le problème à valeurs
initiales, et on ajuste $\gamma$ pour que $V(l, \gamma) = V_0$ (exercice V.5).

\Question{2-b- Algorithme.}
\Etape{1} Démarrage hors de la singularité : pour $\gamma > 0$, près de $0$,
$V \approx \gamma x$, donc $V'' \approx a/\sqrt{\gamma x}$ ; en intégrant deux
fois : $W(\varepsilon) \approx \gamma + 2a\sqrt{\varepsilon/\gamma}$ et
$V(\varepsilon) \approx \gamma\varepsilon + \frac{4a}{3}\varepsilon^{3/2}/\sqrt\gamma$.

\Etape{2} Deux tirs $\gamma_0$, $\gamma_1$, puis la formule de la sécante du cours
jusqu'à $\abs{V(l, \gamma) - V_0} < \text{tolérance}$.

\Etape{3} Remarque physique : pour $\gamma = 0$ on retrouve Child-Langmuir,
$V(l) = Cl^{4/3}$. Le calcul montre que $V(l, \gamma)$ croît avec $\gamma$
($1{,}7173$ ; $1{,}7233$ ; $1{,}8429$ ; $2{,}1223$ ; $2{,}8746$ pour
$\gamma = 0{,}01$ ; $0{,}1$ ; $0{,}5$ ; $1$ ; $2$, avec $a = l = 1$) : une solution
avec $V \ge 0$ n'existe que si $V_0 \ge Cl^{4/3}$.

\begin{algo}[ : tir pour $V(0) = 0$, $V(l) = V_0$]
  \Require $a$, $l$, $V_0$, $\varepsilon$, $n$, $\gamma_0$, $\gamma_1$, tolérance, $k_{\max}$
  \State $R_0 \leftarrow V(l, \gamma_0) - V_0$ \Comment{RK2 depuis $x = \varepsilon$, étape 1}
  \For{$k$ allant de $1$ à $k_{\max}$}
    \State $R_1 \leftarrow V(l, \gamma_1) - V_0$
    \If{$\abs{R_1} < \text{tolérance}$} \Print $\gamma_1$ \Stop \EndIf
    \State $\gamma \leftarrow \gamma_1 - (\gamma_1 - \gamma_0)R_1/(R_1 - R_0)$
    \State $\gamma_0 \leftarrow \gamma_1$ ; $R_0 \leftarrow R_1$ ; $\gamma_1 \leftarrow \gamma$
  \EndFor
\end{algo}

\Verif $a = l = 1$, $V_0 = 2{,}5$, $\gamma_0 = 0{,}1$, $\gamma_1 = 1$ : la fonction
\texttt{tir} de \fichier{ex6_rk2.py} donne $\gamma = 1{,}526741$
($\varepsilon = 10^{-2}$, $n = 4000$), $1{,}529373$ ($\varepsilon = 10^{-3}$,
$n = 16\,000$) et $1{,}529628$ ($\varepsilon = 10^{-4}$, $n = 64\,000$), en 6 ou
7 itérations ; une intégration de référence très précise (SciPy) donne
$\gamma = 1{,}52966$. Près de la singularité $x = 0$, il faut à la fois un
départ $\varepsilon$ petit (le démarrage n'est qu'approché) et un pas fin
(la dérivée seconde y est très grande).

\Refs Child-Langmuir : [CL] ; \BF{RK pour les systèmes, alg.~5.7, p.~333 ;
tir, §~11.2, p.~694}.

% ------------------------------------------------------------------
\section{Exercice 7 : équation du troisième ordre par la méthode de tir}

\Attention Une équation du troisième ordre demande trois conditions ; l'énoncé
n'en donne que deux, $y(a) = y_1$ et $y(b) = y_2$. Le problème est donc
incomplet. On suppose ici une troisième condition donnée en $x = a$,
$y'(a) = y'_1$ (avec une condition en $b$, par exemple $y'(b)$, la démarche
serait la même avec deux paramètres de tir).

\Question{1)- Problème aux valeurs initiales.}
\Pourquoi En $x = a$, on connaît $y$ et $y'$ ; il manque $y''(a)$, que l'on
remplace par le paramètre $\gamma$. On écrit l'équation comme un système du
premier ordre pour pouvoir utiliser les méthodes du chapitre V.

\Etape{1} $Y_1 = y$, $Y_2 = y'$, $Y_3 = y''$ :
\[
  Y_1' = Y_2, \quad Y_2' = Y_3, \quad Y_3' = f(x, Y_1, Y_2, Y_3), \qquad
  Y_1(a) = y_1, \quad Y_2(a) = y'_1, \quad Y_3(a) = \gamma .
\]

\Etape{2} On cherche $\gamma$ tel que $Y_1(b, \gamma) = y_2$ (sécante).

\Question{2-a)- Principe du prédicteur-correcteur (premier ordre).}
Voir l'exercice V.4 : on prédit $y_{k+1}$ par une formule explicite
(Adams-Bashforth d'ordre 2, $p = y_k + \frac h2(3f_k - f_{k-1})$), puis on le
corrige par une formule implicite (trapèzes,
$y_{k+1} = y_k + \frac h2(f_k + f(x_{k+1}, p))$), dans laquelle on utilise la
valeur prédite à la place de l'inconnue.

\Question{2-b)- Algorithme pour le problème du troisième ordre.}
\Pourquoi Les formules sont les mêmes, appliquées au vecteur $Y$ (toutes les
composantes sont prédites, puis toutes sont corrigées). On les place dans la
boucle de tir.

\begin{algo}[ : tir + prédicteur-correcteur pour $y''' = f(x, y, y', y'')$]
  \Require $f$, $a$, $b$, $y_1$, $y'_1$, $y_2$, $N$, $\gamma_0$, $\gamma_1$, $\varepsilon$, $k_{\max}$
  \State $h \leftarrow (b-a)/N$
  \Statex \textbf{Procédure} $\mathrm{PC}(\gamma)$ :
  \State \quad $x \leftarrow a$ ; $Y \leftarrow (y_1, y'_1, \gamma)$ ; $F_p \leftarrow F(x, Y)$
  \State \quad $Y \leftarrow Y + \frac h2\bigl(F_p + F(x + h, Y + hF_p)\bigr)$ ; $x \leftarrow x + h$ \Comment{démarrage RK2}
  \State \quad \textbf{pour} $k = 1$ à $N-1$ : $F_k \leftarrow F(x, Y)$ ; $P \leftarrow Y + \frac h2(3F_k - F_p)$ ;
  \State \quad\quad $Y \leftarrow Y + \frac h2\bigl(F_k + F(x + h, P)\bigr)$ ; $F_p \leftarrow F_k$ ; $x \leftarrow x + h$
  \State \quad retourner $Y_1$ \Comment{$\approx y(b, \gamma)$}
  \Statex \textbf{Programme principal} :
  \State $R_0 \leftarrow \mathrm{PC}(\gamma_0) - y_2$
  \For{$k$ allant de $1$ à $k_{\max}$}
    \State $R_1 \leftarrow \mathrm{PC}(\gamma_1) - y_2$
    \If{$\abs{R_1} < \varepsilon$} \Print $\gamma_1$ \Stop \EndIf
    \State $\gamma \leftarrow \gamma_1 - (\gamma_1 - \gamma_0)R_1/(R_1 - R_0)$ ;
           $\gamma_0 \leftarrow \gamma_1$ ; $R_0 \leftarrow R_1$ ; $\gamma_1 \leftarrow \gamma$
  \EndFor
\end{algo}
avec $F(x, Y) = \bigl(Y_2, Y_3, f(x, Y_1, Y_2, Y_3)\bigr)$.

\Verif Test avec solution connue $y = \sin x + x^2$ : elle vérifie
$y''' = -\cos x = 2x - y'$, $y(0) = 0$, $y'(0) = 1$, $y(1) = \sin 1 + 1$, et
$y''(0) = 2$. L'algorithme retrouve $\gamma = 2{,}00025923$ ($N = 20$),
$2{,}00006339$ ($N = 40$), $2{,}00001560$ ($N = 80$) : erreur divisée par 4,
ordre 2 du prédicteur-correcteur.

\Refs \BF{p.~310--311 ; systèmes et équations d'ordre supérieur, §~5.9,
p.~331--335}.

\chapter{Exercices de fin de document (8 à 13)}

% ------------------------------------------------------------------
\section{Exercice 8 : différences finies et méthode de Jacobi}

\Question{1-a- Dérivée première centrée.}
\Pourquoi C'est le calcul de l'exercice IV.1 : on interpole $y$ par le polynôme
de degré 2 passant par $x - h$, $x$, $x + h$, et on le dérive en $x$.
\Etapes
\Etape{1} Les dérivées des polynômes de base en $x$ valent $-1/(2h)$, $0$ et
$1/(2h)$ (calcul détaillé à l'exercice IV.1).
\Etape{2} D'où $y'(x) \approx \bigl(y(x+h) - y(x-h)\bigr)/(2h)$, avec l'erreur
$-\frac{h^2}{6}y'''(\xi)$.

\Question{1-b- Dérivée seconde.}
\Pourquoi On applique la formule de a) deux fois, à la fonction $y'$, avec le pas
$h/2$ : $y''(x) \approx \bigl[y'(x + h/2) - y'(x - h/2)\bigr]/h$.
\Etapes
\Etape{1} Par a) avec le pas $h/2$ : $y'(x + h/2) \approx \bigl(y(x+h) - y(x)\bigr)/h$
et $y'(x - h/2) \approx \bigl(y(x) - y(x-h)\bigr)/h$.
\Etape{2} D'où
\[
  y''(x) \approx \frac{y(x+h) - 2y(x) + y(x-h)}{h^2} .
\]
\Etape{3} Erreur : en additionnant les développements de Taylor de $y(x \pm h)$
jusqu'à l'ordre 4, les termes impairs disparaissent et il reste
$-\frac{h^2}{12}y^{(4)}(\xi)$. C'est pourquoi l'énoncé suppose $y$ de classe
$C^4$.

\Question{2- Système algébrique.}
\Pourquoi On écrit l'équation aux nœuds $x_i = a + ih$, $i = 1, \dots, n-1$,
$h = (b-a)/n$, en remplaçant $y'$ et $y''$ par les formules centrées. Les
inconnues sont $y_1, \dots, y_{n-1}$ ; $y_0 = c$ et $y_n = d$ sont connues.
\Etapes
\Etape{1} $\dfrac{y_{i+1} - 2y_i + y_{i-1}}{h^2}
          = p_i\dfrac{y_{i+1} - y_{i-1}}{2h} + q_iy_i + g_i$.
\Etape{2} En multipliant par $h^2$ et en regroupant :
\[
  \Bigl(1 + \frac h2p_i\Bigr)y_{i-1} - \bigl(2 + h^2q_i\bigr)y_i
  + \Bigl(1 - \frac h2p_i\Bigr)y_{i+1} = h^2g_i, \qquad i = 1, \dots, n-1 .
\]
\Etape{3} Pour $i = 1$ et $i = n-1$, les termes connus $y_0 = c$ et $y_n = d$
passent au second membre. Le système est \emph{tridiagonal}, d'erreur de
troncature $O(h^2)$. Burden et Faires montrent qu'il a une solution unique si
$q \ge 0$ et $h < 2/L$, avec $L = \max\abs{p}$ \BF{th.~11.3, p.~702} ; la
condition $h\abs{p_i} < 2$ rend aussi $1 \pm \frac h2p_i$ positifs.

\Question{3-a- Principe de Jacobi.}
On isole $y_i$ dans l'équation $i$ et on calcule les nouvelles valeurs à partir
des anciennes uniquement (exercice I.7) :
\[
  y_i^{(k+1)} = \frac{h^2g_i - \bigl(1 + \frac h2p_i\bigr)y_{i-1}^{(k)}
                - \bigl(1 - \frac h2p_i\bigr)y_{i+1}^{(k)}}{-\bigl(2 + h^2q_i\bigr)} .
\]
Le coefficient diagonal $2 + h^2q_i$ est au moins égal à la somme
$\abs{1 + \frac h2p_i} + \abs{1 - \frac h2p_i} = 2$ des autres (si $q_i \ge 0$
et $h\abs{p_i} \le 2$) : la matrice est à diagonale dominante, ce qui favorise
la convergence (lente, car la dominance n'est pas stricte quand $q_i = 0$).

\Question{3-b- Algorithme.}
\begin{algo}[ : différences finies + Jacobi]
  \Require $a$, $b$, $c$, $d$, $n$, $p$, $q$, $g$, $\varepsilon$, $k_{\max}$
  \State $h \leftarrow (b-a)/n$
  \For{$i$ allant de $1$ à $n-1$}
    \State $x \leftarrow a + ih$ ; $B_i \leftarrow 1 + \frac h2p(x)$ ; $A_i \leftarrow -(2 + h^2q(x))$
    \State $C_i \leftarrow 1 - \frac h2p(x)$ ; $D_i \leftarrow h^2g(x)$ ; $y_i \leftarrow 0$
  \EndFor
  \State $y_0 \leftarrow c$ ; $y_n \leftarrow d$ ; $k \leftarrow 0$
  \Repeat
    \State $k \leftarrow k + 1$ ; $\mathit{num} \leftarrow 0$ ; $\mathit{den} \leftarrow 0$
    \For{$i$ allant de $1$ à $n-1$}
      \State $z_i \leftarrow (D_i - B_iy_{i-1} - C_iy_{i+1})/A_i$
      \State $\mathit{num} \leftarrow \mathit{num} + \abs{z_i - y_i}$ ; $\mathit{den} \leftarrow \mathit{den} + \abs{z_i}$
    \EndFor
    \For{$i$ allant de $1$ à $n-1$} \State $y_i \leftarrow z_i$ \EndFor
  \Until{$\mathit{num}/\mathit{den} < \varepsilon$ \textbf{ou} $k \ge k_{\max}$}
  \Print $k$, $(x_i, y_i)$
\end{algo}

\Question{3-c- Programme Pascal.}
\programme{Pascal}{en Pascal}{ex8_jacobi.pas}

\Verif Le programme a été testé sur l'exemple~1 de \BF{p.~703} :
$y'' = -\frac2xy' + \frac2{x^2}y + \frac{\sin(\ln x)}{x^2}$, $1 \le x \le 2$,
$y(1) = 1$, $y(2) = 2$, $n = 10$. Avec $\varepsilon = 10^{-12}$, Jacobi converge
en 450 itérations et donne $y(1{,}1) = 1{,}09260052$, $y(1{,}5) = 1{,}48112026$,
$y(1{,}9) = 1{,}89392110$, etc. : \emph{toutes} les valeurs coïncident avec la
table~11.3 de Burden et Faires (obtenue par résolution directe). Le nombre élevé
d'itérations illustre la lenteur de Jacobi (Gauss-Seidel ou la résolution
tridiagonale directe, exercice VI.4, seraient plus rapides).

\Refs \BF{(4.9), p.~178 ; §~11.3, alg.~11.3 et th.~11.3, p.~700--703}.

% ------------------------------------------------------------------
\section{Exercice 9 : oscillateur de Duffing forcé et amorti}

\Question{1)- Équation du mouvement et adimensionnement.}
\Pourquoi La force qui dérive du potentiel est $-\dd u/\dd y = -ky - k_1y^3$ ;
la viscosité linéaire donne $-\alpha\dot y$ ; la force extérieure est
$F\cos\omega t$. Adimensionner réduit le nombre de paramètres et donne des
nombres d'ordre 1, plus sûrs numériquement.

\Etapes
\Etape{1} Deuxième loi de Newton :
\[
  M\ddot y + \alpha\dot y + ky + k_1y^3 = F\cos\omega t .
\]
\Etape{2} Échelle de temps : la pulsation propre $\omega_0 = \sqrt{k/M}$,
$\tau = \omega_0t$. Échelle de longueur : le déplacement statique
$y_s = F/k$, $Y = y/y_s$. Alors $\dot y = y_s\omega_0Y'$ et
$\ddot y = y_s\omega_0^2Y'' = (F/M)Y''$ (le prime désigne $\dd/\dd\tau$).
\Etape{3} En reportant et en divisant par $F$ :
\[
  Y'' + \mu Y' + Y + \beta Y^3 = \cos(\Omega\tau),
  \qquad \mu = \frac{\alpha}{M\omega_0} = \frac{\alpha}{\sqrt{kM}},
  \quad \beta = \frac{k_1F^2}{k^3},
  \quad \Omega = \frac{\omega}{\omega_0} .
\]
Il ne reste que trois paramètres sans dimension. (D'autres choix d'échelle de
longueur sont possibles, par exemple $\sqrt{k/k_1}$, qui donne $\beta = 1$ et
une amplitude de forçage sans dimension.) Conditions initiales :
$Y(0) = Y'(0) = 0$.

\Question{2)- Algorithme RK4.}
\Pourquoi RK4 est une méthode à un pas : il suffit d'écrire l'équation du second
ordre comme un système du premier ordre et d'appliquer RK4 au vecteur
(voir l'exercice 5).

\Etape{1} $X_1 = Y$, $X_2 = Y'$ : $X_1' = X_2$,
$X_2' = \cos(\Omega\tau) - \mu X_2 - X_1 - \beta X_1^3$.
\Etape{2} À chaque pas, calcul des quatre pentes vectorielles, puis combinaison
$1, 2, 2, 1$.

\begin{algo}[ : RK4 pour le Duffing adimensionné]
  \Require $\mu$, $\beta$, $\Omega$, $h$, $N$
  \State $\tau \leftarrow 0$ ; $X_1 \leftarrow 0$ ; $X_2 \leftarrow 0$
  \For{$i$ allant de $1$ à $N$}
    \State $k_1 \leftarrow hX_2$ ; $l_1 \leftarrow hG(\tau, X_1, X_2)$
    \State $k_2 \leftarrow h(X_2 + l_1/2)$ ; $l_2 \leftarrow hG(\tau + h/2, X_1 + k_1/2, X_2 + l_1/2)$
    \State $k_3 \leftarrow h(X_2 + l_2/2)$ ; $l_3 \leftarrow hG(\tau + h/2, X_1 + k_2/2, X_2 + l_2/2)$
    \State $k_4 \leftarrow h(X_2 + l_3)$ ; $l_4 \leftarrow hG(\tau + h, X_1 + k_3, X_2 + l_3)$
    \State $X_1 \leftarrow X_1 + (k_1 + 2k_2 + 2k_3 + k_4)/6$ ; $X_2 \leftarrow X_2 + (l_1 + 2l_2 + 2l_3 + l_4)/6$
    \State $\tau \leftarrow \tau + h$
    \Print $\tau$, $X_1$, $X_2$
  \EndFor
\end{algo}
avec $G(\tau, X_1, X_2) = \cos(\Omega\tau) - \mu X_2 - X_1 - \beta X_1^3$.

\Question{3)- Moindres carrés pour les $a_n$ et $b_n$.}
\Pourquoi On dispose de beaucoup plus de valeurs $y_i$ que d'inconnues
($2N$) : on ne peut pas interpoler, on minimise
\[
  R = \sum_{i=1}^{m}\Bigl[y_i - \sum_{n=1}^{N}\bigl(a_n\cos n\omega t_i + b_n\sin n\omega t_i\bigr)\Bigr]^2 .
\]
Le modèle est linéaire en les $a_n$, $b_n$ : c'est le cadre du chapitre III avec
les $2N$ fonctions $g = \cos n\omega t$ et $\sin n\omega t$.

\Etapes
\Etape{1} Équations normales : $\partial R/\partial a_n = 0$ et
$\partial R/\partial b_n = 0$ donnent, pour $n = 1, \dots, N$,
\[
  \sum_{p=1}^{N}\Bigl[a_p\sum_i\cos p\omega t_i\cos n\omega t_i
  + b_p\sum_i\sin p\omega t_i\cos n\omega t_i\Bigr] = \sum_i y_i\cos n\omega t_i ,
\]
et la même équation avec $\sin n\omega t_i$ à la place de $\cos n\omega t_i$
en facteur. C'est un système $2N \times 2N$, symétrique ; on calcule ses
coefficients avec l'algorithme de l'exercice III.2 et on le résout par Gauss.

\Etape{2} Cas particulier très utile : si les $t_i$ sont équidistants et
couvrent exactement une période $T = 2\pi/\omega$ ($t_i = t_0 + iT/m$,
$i = 0, \dots, m-1$, avec $m > 2N$), les sommes croisées sont nulles
(orthogonalité discrète) et $\sum_i\cos^2 = \sum_i\sin^2 = m/2$. Le système
devient diagonal :
\[
  a_n = \frac2m\sum_i y_i\cos n\omega t_i, \qquad
  b_n = \frac2m\sum_i y_i\sin n\omega t_i .
\]

\Etape{3} Précautions. Le modèle suppose un régime périodique de période
$2\pi/\omega$ et sans terme constant : il faut écarter le régime transitoire
(qui dépend des conditions initiales) avant de prendre les $y_i$. En variables
sans dimension, la pulsation à utiliser est $\Omega$ et le temps $\tau$.

\begin{algo}[ : coefficients de Fourier par moindres carrés]
  \Require $m$ valeurs $(t_i, y_i)$ en régime permanent, $\omega$, $N$
  \For{$i$ allant de $1$ à $m$}
    \For{$n$ allant de $1$ à $N$}
      \State $G_{i,n} \leftarrow \cos n\omega t_i$ ; $G_{i,N+n} \leftarrow \sin n\omega t_i$
    \EndFor
  \EndFor
  \State former $A_{pq} = \sum_i G_{ip}G_{iq}$ et $B_p = \sum_i G_{ip}y_i$ ($p, q = 1, \dots, 2N$) \Comment{exercice III.2}
  \State résoudre $A\,c = B$ par Gauss \Comment{$c = (a_1, \dots, a_N, b_1, \dots, b_N)$}
  \Statex (si échantillonnage uniforme sur une période : $a_n \leftarrow \frac2m\sum_iy_i\cos n\omega t_i$, $b_n \leftarrow \frac2m\sum_iy_i\sin n\omega t_i$)
  \Print $a_n$, $b_n$
\end{algo}

\Verif Avec $\mu = 0{,}2$, $\beta = 0{,}1$, $\Omega = 1{,}2$, RK4 avec $64$ pas par
période, $100$ périodes écartées puis $m = 64$ points sur une période,
$N = 5$ : $a_1 = 2{,}07476$, $b_1 = 2{,}12381$, $a_3 = -0{,}04201$,
$b_3 = 0{,}04485$, $a_5 = -0{,}00091$, $b_5 = -0{,}00083$, et $a_2 = b_2 = a_4 =
b_4 = 0$ (le terme $Y^3$ étant impair, seules les harmoniques impaires
apparaissent). Les équations normales et les formules d'orthogonalité donnent
les mêmes valeurs (écart $3\times 10^{-12}$), et le reste vaut
$R = 2\times 10^{-8}$.

\Refs \BF{RK4 pour les systèmes, alg.~5.7, p.~333 ; moindres carrés, §~8.1 ;
approximation trigonométrique discrète, §~8.5, p.~545--550}.

% ------------------------------------------------------------------
\section{Exercice 10 : méthode de Gauss}

\Question{1- Principe.} On élimine les inconnues une à une : à l'étape $k$, on
choisit l'équation pivotale (plus grand $\abs{a_{ik}}$, $i \ge k$), on soustrait
de chaque ligne $l > k$ le multiple $a_{lk}/a_{kk}$ de la ligne $k$ pour
annuler le coefficient de $x_k$. Après $n-1$ étapes, le système est
triangulaire ; on le résout par remontée. \emph{Pourquoi le pivot} : éviter la
division par zéro et l'amplification des erreurs d'arrondi (exercice I.2).

\Question{2- Algorithme de triangularisation.} C'est l'algorithme de l'exercice
I.3 (boucle sur $k$ avec recherche du pivot, échange, puis élimination).

\Question{3- Résolution du système triangulaire.} C'est l'algorithme de
remontée de l'exercice I.1.

\Question{4- Programme Pascal.}
Le programme contient les deux procédures : \texttt{triangularise} (questions 1
et 2) et \texttt{remontee} (question 3).
\programme{Pascal}{en Pascal}{gauss.pas}

\Verif Sur le système $2x_2 + x_3 = 5$, $x_1 + x_2 + x_3 = 6$,
$2x_1 + x_2 + 3x_3 = 13$ (où $a_{11} = 0$), le programme donne
$x = (7/3, 4/3, 7/3)$ (détail du calcul à l'exercice I.3) ; sur un système
$6\times 6$ aléatoire avec $a_{11} = 0$, l'erreur est de $4{,}4\times 10^{-11}$.

\Refs \BF{alg.~6.1, p.~368 ; alg.~6.2, p.~378}.

% ------------------------------------------------------------------
\section{Exercice 11 : le pendule simple}

\Question{1-a- Algorithme RK2.}
\Pourquoi On écrit $Y'' + \omega_0^2\sin Y = 0$ comme un système :
$Y' = Z$, $Z' = -\omega_0^2\sin Y$, avec $Y(0) = Y_0$, $Z(0) = Z_0$ (pour un
pendule lâché sans vitesse, $Z_0 = 0$). RK2 (exercice V.1) s'applique au vecteur
$(Y, Z)$ : on calcule les deux pentes de Euler pour les deux composantes,
puis on fait la moyenne.

\begin{algo}[ : RK2 pour le pendule]
  \Require $\omega_0$, $Y_0$, $Z_0$, $t_{\max}$, $n$
  \State $h \leftarrow t_{\max}/n$ ; $t \leftarrow 0$ ; $Y \leftarrow Y_0$ ; $Z \leftarrow Z_0$
  \For{$i$ allant de $1$ à $n$}
    \State $k_{1Y} \leftarrow hZ$ ; $k_{1Z} \leftarrow -h\omega_0^2\sin Y$
    \State $k_{2Y} \leftarrow h(Z + k_{1Z})$ ; $k_{2Z} \leftarrow -h\omega_0^2\sin(Y + k_{1Y})$
    \State $Y \leftarrow Y + (k_{1Y} + k_{2Y})/2$ ; $Z \leftarrow Z + (k_{1Z} + k_{2Z})/2$ ; $t \leftarrow t + h$
    \Print $t$, $Y$, $Z$
  \EndFor
\end{algo}

\Question{1-b- Programme Fortran.}
\programme{[90]Fortran}{en Fortran 90}{ex11_pendule.f90}

\Verif $\omega_0 = 1$, $Y_0 = 2\pi/9$, $Z_0 = 0$, $h = 10^{-3}$, comparaison
avec une intégration de référence (SciPy, tolérance $10^{-12}$) :
\begin{center}
\begin{tabular}{lcccc}
  \toprule
  $t$ & 1 & 2 & 3 & 4\\ \midrule
  RK2 & $0{,}39768602$ & $-0{,}25372859$ & $-0{,}67967973$ & $-0{,}51945770$\\
  référence & $0{,}39768609$ & $-0{,}25372840$ & $-0{,}67967966$ & $-0{,}51945795$\\
  \bottomrule
\end{tabular}
\end{center}
Les écarts sont inférieurs à $3\times 10^{-7}$.

\Question{2-a- Formule composite des trapèzes et algorithme pour (2).}
\Pourquoi On découpe $[0, \pi/2]$ en $n$ intervalles de longueur
$h = \pi/(2n)$ ; sur chacun, l'aire sous la courbe est remplacée par celle du
trapèze : $\frac h2\bigl(f(x_{k-1}) + f(x_k)\bigr)$. En sommant, chaque nœud
intérieur est compté deux fois :
\[
  \int_0^{\pi/2} f \approx h\Bigl[\tfrac12f(0) + \sum_{k=1}^{n-1}f(kh) + \tfrac12f(\pi/2)\Bigr],
  \qquad f(x) = \bigl(1 - A\sin^2x\bigr)^{-1/2},
\]
d'erreur $-\frac{h^2}{12}\cdot\frac\pi2f''(\eta)$ (chapitre IV). Puis
$T/T_0 = \frac2\pi\times$ intégrale. L'intégrande est défini et régulier car
$A = \sin^2(Y_0/2) < 1$ pour $\abs{Y_0} < \pi$.

\begin{algo}[ : période par les trapèzes]
  \Require $Y_0$, $n$
  \State $A \leftarrow \sin^2(Y_0/2)$ ; $h \leftarrow \pi/(2n)$
  \State $S \leftarrow \bigl(f(0) + f(\pi/2)\bigr)/2$
  \For{$k$ allant de $1$ à $n-1$} \State $S \leftarrow S + f(kh)$ \EndFor
  \Print $T/T_0 = \frac2\pi\,h\,S$
\end{algo}

\Question{2-b- Programmes.}
\programme{Pascal}{en Pascal}{ex11_trapeze.pas}
\programme{[90]Fortran}{en Fortran 90}{ex11_trapeze.f90}

\Verif Avec seulement $n = 4$ :
$T/T_0 = 1{,}00190719$ ; $1{,}00766903$ ; $1{,}01740880$ ; $1{,}03134052$ pour
$Y_0 = \pi/18$, $\pi/9$, $\pi/6$, $2\pi/9$, et $n = 8$ donne les mêmes huit
chiffres. Cette précision, bien meilleure que ce qu'annonce l'erreur en $h^2$,
s'explique : l'intégrande est régulier, pair et de période $\pi$, et les
trapèzes sur une période d'une fonction périodique régulière convergent
exponentiellement vite [TW]. Les valeurs coïncident avec la formule exacte
$T = \frac{2T_0}\pi K\bigl(\sin(Y_0/2)\bigr)$ ($K$ : intégrale elliptique
complète de première espèce) [LA].

\Question{3- Moindres carrés $T/T_0 = a + bY_0^2$.}
\Pourquoi Le modèle est linéaire en $a$ et $b$ avec la variable $X = Y_0^2$ :
c'est la droite des moindres carrés (chapitre III) en $X$.
\Etapes
\Etape{1} $X_i = Y_{0,i}^2 = 0{,}030462$ ; $0{,}121847$ ; $0{,}274156$ ;
$0{,}487388$. Sommes ($m = 4$) : $\sum X = 0{,}913852$,
$\sum X^2 = 0{,}328483$, $\sum T/T_0 = 4{,}06$,
$\sum X\,T/T_0 = 0{,}934993$.
\Etape{2} Système :
$4a + 0{,}913852\,b = 4{,}06$ ; $0{,}913852\,a + 0{,}328483\,b = 0{,}934993$.
Déterminant $0{,}478806$.
\Etape{3} $a = (0{,}328483 \times 4{,}06 - 0{,}913852 \times 0{,}934993)/0{,}478806
      = 1{,}000814$ ; $b = (4 \times 0{,}934993 - 0{,}913852 \times 4{,}06)/0{,}478806
      = 0{,}062093$.

\Attention Les valeurs exactes de $T/T_0$ (question 2) sont $1{,}0019$ ;
$1{,}0077$ ; $1{,}0174$ ; $1{,}0313$. Trois valeurs du tableau sont correctes à
l'arrondi près, mais la troisième ($1{,}020$ pour $\pi/6$) ne l'est pas (la
valeur exacte est $1{,}0174$) ; c'est aussi le point du plus grand écart à la
droite ($0{,}0022$). Avec les valeurs exactes, les moindres carrés donnent
$a = 0{,}999856$ et $b = 0{,}064452$. Pour comparaison, le développement
théorique pour les petites amplitudes est $T/T_0 = 1 + Y_0^2/16 + \dots$, soit
$a = 1$ et $b = 0{,}0625$ [LA].

\Refs \BF{alg.~5.7, p.~333 ; th.~4.5, p.~205 ; §~8.1} ; [TW], [LA].

% ------------------------------------------------------------------
\section{Exercice 12 : capacité calorifique et chaleur absorbée}

\Question{1- Chaleur $Q$ par Simpson.}
\Pourquoi À volume constant, $\dd Q = C_v\,\dd T$, donc
$Q = \int_{0{,}1}^{0{,}5} C_v(T)\,\dd T$. On a trois points équidistants
($h = 0{,}2$) : c'est exactement le cas de la formule de Simpson simple.
\Etapes
\Etape{1} $Q \approx \frac h3\bigl[C_v(0{,}1) + 4C_v(0{,}3) + C_v(0{,}5)\bigr]$.
\Etape{2} $Q \approx \frac{0{,}2}{3}\bigl[0{,}211 + 4 \times 0{,}694 + 1{,}361\bigr]\times 10^{-4}
         = \frac{0{,}2}{3} \times 4{,}348 \times 10^{-4}$.
\Etape{3} $Q \approx 0{,}289867 \times 10^{-4}$ cal/mol.

\Question{2-a- Système vérifié par $a$ et $b$.}
Ce sont les équations normales de la droite (complément du chapitre III) :
\[
  m\,a + \Bigl(\sum T_i\Bigr)b = \sum C_i, \qquad
  \Bigl(\sum T_i\Bigr)a + \Bigl(\sum T_i^2\Bigr)b = \sum T_iC_i .
\]

\Question{2-b- Algorithmes.}
\Pourquoi Les quatre sommes se calculent en une seule boucle sur les mesures ;
le système $2\times 2$ se résout par la règle de Cramer (pas besoin de Gauss
pour deux inconnues), à condition que le déterminant
$m\sum T_i^2 - (\sum T_i)^2$ soit non nul.
\begin{algo}[ : droite des moindres carrés]
  \State $s_1 \leftarrow 0$ ; $s_2 \leftarrow 0$ ; $s_y \leftarrow 0$ ; $s_{ty} \leftarrow 0$
  \For{$i$ allant de $1$ à $m$}
    \State $s_1 \leftarrow s_1 + T_i$ ; $s_2 \leftarrow s_2 + T_i^2$ ; $s_y \leftarrow s_y + C_i$ ; $s_{ty} \leftarrow s_{ty} + T_iC_i$
  \EndFor
  \State $D \leftarrow m\,s_2 - s_1^2$
  \State $a \leftarrow (s_2s_y - s_1s_{ty})/D$ ; $b \leftarrow (m\,s_{ty} - s_1s_y)/D$
  \Print $a$, $b$
\end{algo}

\Question{2-c- Programme Pascal.} Il contient aussi la question 2-d.
\programme{Pascal}{en Pascal}{ex12_mc.pas}

\Question{2-d- $Q$ par Simpson composite sur les mesures.}
\Pourquoi Simpson composite exige des abscisses équidistantes et un nombre
\emph{pair} d'intervalles, donc un nombre \emph{impair} de mesures. Les poids
sont $1, 4, 2, 4, \dots, 4, 1$ (exercice IV.5). Si l'on préfère utiliser la
droite ajustée, on intègre exactement $a + bT$.
\begin{algo}[ : Simpson composite sur un tableau]
  \If{$m$ pair} \Print \og nombre de points impair requis \fg{} \Stop \EndIf
  \State $h \leftarrow T_2 - T_1$ ; $Q \leftarrow C_1 + C_m$
  \For{$i$ allant de $2$ à $m-1$}
    \If{$i$ pair} \State $Q \leftarrow Q + 4C_i$ \Else \State $Q \leftarrow Q + 2C_i$ \EndIf
  \EndFor
  \State $Q \leftarrow hQ/3$
  \Print $Q$
\end{algo}
(Avec la numérotation $i = 1, \dots, m$, les points d'indice pair sont les
milieux des paires d'intervalles.)

\Verif Sur les trois mesures données, le programme affiche le système
$3a + 0{,}9b = 2{,}266$ ; $0{,}9a + 0{,}35b = 0{,}9098$, puis
$a = -0{,}107167$ et $b = 2{,}875$ (en $10^{-4}$ cal/(K\,mol)), et
$Q = 0{,}289867$ ($\times 10^{-4}$ cal/mol) par Simpson. L'intégrale exacte de la
droite, $a(0{,}5 - 0{,}1) + b(0{,}5^2 - 0{,}1^2)/2 = 0{,}302133$, diffère de 4~\% :
avec seulement trois points, une droite décrit mal des données qui ne sont pas
alignées.

\Refs \BF{th.~4.4, p.~204 ; alg.~4.1, p.~205 ; §~8.1, p.~507--508}.

% ------------------------------------------------------------------
\section{Exercice 13 : équation intégrale de Fredholm}

\Attention L'énoncé écrit $t_i = ih$ pour $n$ points. Pour que la formule des
trapèzes utilise les bornes $a$ et $b$, on prend, comme le cours (chapitre V),
$t_i = a + (i-1)h$, $i = 1, \dots, n$, avec $h = (b-a)/(n-1)$ ; ainsi
$t_1 = a$ et $t_n = b$.

\Question{1)- Coefficients $c_i$.}
\Pourquoi La formule composite des trapèzes donne le poids $h/2$ aux deux
extrémités et $h$ aux points intérieurs. On l'écrit sous la forme
$(b-a)\sum c_i\,(\dots)$ du cours en divisant les poids par $b - a$.
\Etapes
\Etape{1} $\int_a^b k(x,t)\phi(t)\,\dd t \approx h\bigl[\tfrac12k(x,t_1)\phi_1
  + k(x,t_2)\phi_2 + \dots + k(x,t_{n-1})\phi_{n-1} + \tfrac12k(x,t_n)\phi_n\bigr]$.
\Etape{2} Comme $h = (b-a)/(n-1)$ :
\[
  c_1 = c_n = \frac{1}{2(n-1)}, \qquad c_i = \frac{1}{n-1} \quad (i = 2, \dots, n-1),
\]
et l'on vérifie que $\sum c_i = 1$ (la formule est exacte pour une constante).

\Question{2)- Le système algébrique.}
\Pourquoi L'équation (2) est vraie pour tout $x$ ; en particulier aux nœuds.
En prenant $x = t_j$, le membre de gauche devient $\phi(t_j) = \phi_j$, et le
membre de droite ne contient que les inconnues $\phi_1, \dots, \phi_n$ : on
obtient $n$ équations linéaires à $n$ inconnues, c'est-à-dire (3)
\[
  \phi_j = f_j + (b-a)\sum_{i=1}^{n}c_i\,k(t_j, t_i)\,\phi_i, \qquad j = 1, \dots, n .
\]
(C'est la méthode de Nyström.)

\Question{3)- Les matrices $A$, $X$, $b$.}
On passe les inconnues à gauche :
$\phi_j - (b-a)\sum_i c_ik(t_j,t_i)\phi_i = f_j$. Donc
\[
  A_{ji} = \delta_{ji} - (b-a)\,c_i\,k(t_j, t_i), \qquad
  X = (\phi_1, \dots, \phi_n)^{t}, \qquad b = (f_1, \dots, f_n)^{t},
\]
où $\delta_{ji} = 1$ si $j = i$ et $0$ sinon : $A = I - (b-a)\,K\,C$, avec
$K_{ji} = k(t_j, t_i)$ et $C = \operatorname{diag}(c_1, \dots, c_n)$.

\Question{4)- Résolution par Gauss.}
a) Triangularisation : algorithme de l'exercice I.3 (avec pivot maximal).
b) Système triangulaire : remontée, exercice I.1.
c) Programmes : procédures \texttt{triangularise} et \texttt{remontee} de
\fichier{gauss.pas} (exercice 10, ci-dessus) et leur traduction en Fortran :
\programme{[90]Fortran}{en Fortran 90}{gauss.f90}

\Verif Sur l'exemple du cours, $\phi(x) = \frac{5x}6 - \frac19 +
\frac13\int_0^1(t + x)\phi(t)\,\dd t$ (solution exacte $\phi = x$, donc
$k(x,t) = (x + t)/3$), l'erreur maximale de la méthode des trapèzes vaut
$2{,}48\times 10^{-2}$ ; $6{,}17\times 10^{-3}$ ; $1{,}54\times 10^{-3}$ ;
$3{,}85\times 10^{-4}$ pour $n = 3$, $5$, $9$, $17$ : divisée par 4 quand $h$ est
divisé par 2 (ordre 2 des trapèzes). Le cours, avec Simpson et 3 points, obtient
la solution exacte, car Simpson intègre exactement le produit
$(t + x)\,t$, polynôme de degré 2 en $t$.

\Question{5)- Zéro de $\phi$ par bissection.}
\Pourquoi La résolution donne $\phi$ aux nœuds seulement ; mais la formule (2)
permet ensuite de calculer $\phi(x)$ en \emph{n'importe quel} $x$ (interpolation
de Nyström), ce dont la bissection a besoin. Si $\phi(a)\phi(b) < 0$, la
bissection (exercice II.2) converge vers un zéro.
\begin{algo}[ : zéro de $\phi$]
  \State résoudre (4) : $\phi_1, \dots, \phi_n$
  \State définir $\phi(x) = f(x) + (b-a)\sum_{i=1}^{n}c_i\,k(x,t_i)\,\phi_i$
  \State $\alpha \leftarrow a$ ; $\beta \leftarrow b$
  \If{$\phi(\alpha)\phi(\beta) > 0$} \Print \og pas de changement de signe \fg{} \Stop \EndIf
  \While{$\beta - \alpha > \varepsilon$}
    \State $\gamma \leftarrow (\alpha + \beta)/2$
    \If{$\phi(\alpha)\phi(\gamma) \le 0$} \State $\beta \leftarrow \gamma$ \Else \State $\alpha \leftarrow \gamma$ \EndIf
  \EndWhile
  \Print $(\alpha + \beta)/2$
\end{algo}

\Verif Avec le noyau $(x + t)/3$ sur $[0,1]$ et $f(x) = x - \frac12 - \frac1{36}$,
la solution exacte est $\phi = x - \frac12$ (zéro en $0{,}5$). La bissection sur
la solution de Nyström donne $0{,}494718$ ($n = 5$), $0{,}498680$ ($n = 9$),
$0{,}499670$ ($n = 17$) : l'erreur suit celle de la discrétisation (ordre 2).

\Refs \BF{alg.~6.2, p.~378 ; alg.~2.1, p.~49 ; th.~4.5, p.~205}.

\chapter{Exercices de fin de document (14 à 19)}

% ------------------------------------------------------------------
\section{Exercice 14 : température dans un mur}

\Attention À $t = 0$, la face $x = L$ est à $T_0$ (condition initiale) mais doit
être à $T_L$ (condition aux limites) : les deux conditions ne se raccordent
pas si $T_L \ne T_0$. On impose $T_L$ dès le premier nœud de temps
($T_{m,0} = T_L$) ; ce choix n'influe que sur les tout premiers pas. En $x = 0$,
$T_0 + T_1\sin 0 = T_0$ : pas de conflit. Le temps $t$ est mesuré dans l'unité
implicite de la loi $\sin(\pi t/2)$ (période 4).

\Question{1-a- Forme discrète par différences progressives.}
\Pourquoi C'est le schéma explicite du cours (C2) appliqué à
$\partial T/\partial t = \alpha\,\partial^2T/\partial x^2$ ; seules les conditions
aux limites changent.
\Etapes
\Etape{1} Grille $x_i = i\Delta x$, $\Delta x = L/m$ ; $t_j = j\Delta t$ ;
$T_{i,j} \approx T(x_i, t_j)$ ; $\lambda = \alpha\Delta t/\Delta x^2$.
\Etape{2} Dérivée en temps progressive, dérivée en espace centrée :
\[
  \frac{T_{i,j+1} - T_{i,j}}{\Delta t} = \alpha\,\frac{T_{i+1,j} - 2T_{i,j} + T_{i-1,j}}{\Delta x^2}
  \;\Longrightarrow\;
  T_{i,j+1} = \lambda T_{i+1,j} + (1 - 2\lambda)T_{i,j} + \lambda T_{i-1,j},
\]
pour $i = 1, \dots, m-1$ et $j \ge 0$.
\Etape{3} Conditions : $T_{i,0} = T_0$ ($0 \le i < m$), $T_{m,j} = T_L$,
$T_{0,j} = T_0 + T_1\sin(\pi j\Delta t/2)$.
\Etape{4} Erreur de discrétisation : $O(\Delta t)$ pour la dérivée en temps et
$O(\Delta x^2)$ pour la dérivée en espace, soit $O(\Delta t + \Delta x^2)$
\BF{p.~747}. Le schéma est stable si $\lambda \le \frac12$ (exercice 15).

\Question{1-b- Algorithme.}
\begin{algo}[ : mur, schéma explicite]
  \Require $L$, $\alpha$, $T_0$, $T_1$, $T_L$, $m$, $\Delta t$, $t_{\max}$
  \State $\Delta x \leftarrow L/m$ ; $\lambda \leftarrow \alpha\Delta t/\Delta x^2$
  \If{$\lambda > 1/2$} \Print \og instable : réduire $\Delta t$ \fg{} \EndIf
  \For{$i$ allant de $0$ à $m-1$} \State $U_i \leftarrow T_0$ \EndFor
  \State $U_m \leftarrow T_L$
  \For{$j$ allant de $1$ à $t_{\max}/\Delta t$}
    \For{$i$ allant de $1$ à $m-1$} \State $V_i \leftarrow (1 - 2\lambda)U_i + \lambda(U_{i+1} + U_{i-1})$ \EndFor
    \State $V_0 \leftarrow T_0 + T_1\sin(\pi j\Delta t/2)$ ; $V_m \leftarrow T_L$
    \For{$i$ allant de $0$ à $m$} \State $U_i \leftarrow V_i$ \EndFor
  \EndFor
  \Print $(i\Delta x, U_i)$
\end{algo}

\Question{1-c- Programme Pascal.}
\programme{Pascal}{en Pascal}{ex14_explicite.pas}

\Verif
\begin{itemize}
  \item Sans oscillation ($T_1 = 0$, $T_0 = 20$, $T_L = 100$, $\alpha = 1$,
        $L = 1$, $m = 10$, $\Delta t = 0{,}004$, $\lambda = 0{,}4$), à $t = 2$ le
        programme donne le profil linéaire $20, 28, 36, \dots, 100$ : c'est le
        régime permanent exact $T = T_0 + (T_L - T_0)x/L$ (le schéma reproduit
        exactement les fonctions affines en $x$).
  \item $\alpha = 0{,}1$, $T_0 = 20$, $T_1 = 10$, $T_L = 30$, $\Delta t = 0{,}04$
        ($\lambda = 0{,}4$) : à $t = 3$, $T(0) = 10$ (car $\sin(3\pi/2) = -1$),
        $T(0{,}5) = 24{,}36$, $T(L) = 30$.
  \item Avec $\Delta t = 0{,}06$ ($\lambda = 0{,}6 > \frac12$), le calcul explose
        (valeurs de l'ordre de $10^{5}$ à $t = 3$) : l'instabilité prévue.
\end{itemize}

\Question{2-a- Schéma de Richardson.}
\Pourquoi On remplace la dérivée en temps par la différence centrée
$(T_{i,j+1} - T_{i,j-1})/(2\Delta t)$, d'erreur $O(\Delta t^2)$, pour gagner un
ordre en temps.
\Etapes
\Etape{1} $T_{i,j+1} = T_{i,j-1} + 2\lambda\bigl(T_{i+1,j} - 2T_{i,j} + T_{i-1,j}\bigr)$,
erreur de discrétisation $O(\Delta t^2 + \Delta x^2)$ \BF{(12.14), p.~751}.
\Etape{2} Stabilité (méthode de l'introduction du chapitre VI) :
$\xi - \xi^{-1} = 2\lambda(-4s)$, soit $\xi^2 + 8\lambda s\,\xi - 1 = 0$. Le
produit des racines vaut $-1$ et le discriminant $64\lambda^2s^2 + 4$ est
positif : les racines sont réelles, distinctes, de produit de module 1, donc
l'une a un module $> 1$ dès que $s > 0$. Le schéma est
\textbf{instable quel que soit $\lambda$} (le cours et \BF{p.~751} le
signalent).

\Question{2-b- Algorithme.}
\Pourquoi Le schéma a trois niveaux : il faut $T_{i,0}$ (condition initiale) et
$T_{i,1}$, donné par un pas du schéma explicite, comme le suggère l'énoncé.
\begin{algo}[ : Richardson]
  \State $U^{\text{anc}} \leftarrow$ condition initiale ; $U^{\text{cour}} \leftarrow$ un pas explicite (question 1)
  \For{$j$ allant de $1$ à $N-1$}
    \For{$i$ allant de $1$ à $m-1$}
      \State $U^{\text{nouv}}_i \leftarrow U^{\text{anc}}_i + 2\lambda(U^{\text{cour}}_{i+1} - 2U^{\text{cour}}_i + U^{\text{cour}}_{i-1})$
    \EndFor
    \State $U^{\text{nouv}}_0 \leftarrow T_0 + T_1\sin\bigl(\pi(j+1)\Delta t/2\bigr)$ ; $U^{\text{nouv}}_m \leftarrow T_L$
    \State $U^{\text{anc}} \leftarrow U^{\text{cour}}$ ; $U^{\text{cour}} \leftarrow U^{\text{nouv}}$
  \EndFor
\end{algo}

\Question{2-c- Programme Fortran.}
\programme{[90]Fortran}{en Fortran 90}{ex14_richardson.f90}

\Verif Même avec $\lambda = 0{,}1$ (très en dessous de la limite du schéma
explicite), le programme donne $T(L/2) = -3{,}3\times 10^{2}$ à $t = 0{,}3$,
$-3{,}7\times 10^{17}$ à $t = 1{,}2$ et $-2{,}2\times 10^{47}$ à $t = 3$ : la
divergence confirme l'analyse. En pratique, pour l'ordre 2 en temps, on utilise
Crank-Nicolson (exercice VI.4).

\Refs \BF{p.~743--747 (schéma explicite) ; (12.14), p.~751 (Richardson)}.

% ------------------------------------------------------------------
\section{Exercice 15 : conditions de stabilité}

\Question{1) Schéma explicite à une dimension.}
\Pourquoi On applique l'analyse de von Neumann (introduction du chapitre VI) :
une perturbation vérifie le même schéma que la solution.
\Etapes
\Etape{1} $\varepsilon_{i,j+1} = (1 - 2\lambda)\varepsilon_{i,j} + \lambda(\varepsilon_{i+1,j} + \varepsilon_{i-1,j})$,
$\lambda = a^2k/h^2$.
\Etape{2} Avec $\varepsilon_{i,j} = \xi^{\,j}e^{Ii\theta}$ :
$\xi = 1 - 2\lambda + 2\lambda\cos\theta = 1 - 4\lambda s$, $s = \sin^2(\theta/2)$.
\Etape{3} $\abs{\xi} \le 1 \iff -1 \le 1 - 4\lambda s \le 1$. L'inégalité de
droite est toujours vraie ; celle de gauche s'écrit $\lambda s \le \frac12$ pour
tout $s \in [0,1]$, donc (cas le plus défavorable $s = 1$, mode le plus
oscillant)
\[
  \frac{a^2k}{h^2} \le \frac12 .
\]

\Question{2) Problème plan.}
\Etape{1} Schéma :
$u_{i,l}^{j+1} = u_{i,l}^{j} + \lambda_1\delta_x^2u^{j} + \lambda_2\delta_y^2u^{j}$, avec
$\lambda_1 = a^2k/h_1^2$, $\lambda_2 = a^2k/h_2^2$.
\Etape{2} Mode $\xi^{\,j}e^{I(i\theta_1 + l\theta_2)}$ :
$\xi = 1 - 4\lambda_1s_1 - 4\lambda_2s_2$.
\Etape{3} Le cas le plus défavorable est $s_1 = s_2 = 1$ :
$\xi = 1 - 4(\lambda_1 + \lambda_2) \ge -1$, d'où
\[
  a^2k\Bigl(\frac{1}{h_1^2} + \frac{1}{h_2^2}\Bigr) \le \frac12 ,
  \qquad\text{soit}\qquad \frac{a^2k}{h^2} \le \frac14 \ \text{ si } h_1 = h_2 = h .
\]
Le pas de temps admissible est deux fois plus petit qu'en dimension 1.

\Verif $\max\abs{\xi} = 1$ pour $(\lambda_1, \lambda_2) = (0{,}25 ; 0{,}25)$ et
$(0{,}3 ; 0{,}2)$, mais $1{,}2$ pour $(0{,}3 ; 0{,}25)$.

\Question{3) Advection-diffusion dans le plan.}
\Etapes
\Etape{1} Schéma (temps progressif, espace centré, $h$ dans les deux
directions) avec $d = \nu k/h^2$, $c_1 = u_1k/h$, $c_2 = u_2k/h$ :
\[
  \Omega^{j+1}_{i,l} = \Omega^{j}_{i,l}
  - \frac{c_1}{2}\bigl(\Omega^{j}_{i+1,l} - \Omega^{j}_{i-1,l}\bigr)
  - \frac{c_2}{2}\bigl(\Omega^{j}_{i,l+1} - \Omega^{j}_{i,l-1}\bigr)
  + d\bigl(\delta_x^2 + \delta_y^2\bigr)\Omega^{j}_{i,l} .
\]
\Etape{2} Facteur d'amplification :
\[
  \xi = 1 - 4d(s_1 + s_2) - I\bigl(c_1\sin\theta_1 + c_2\sin\theta_2\bigr).
\]
\Etape{3} \emph{Condition établie par l'énoncé.} Pour $\theta_1 = \theta_2 = \pi$
($s_1 = s_2 = 1$, sinus nuls), $\xi = 1 - 8d$, et $\abs{\xi} \le 1$ impose
\[
  \frac{\nu k}{h^2} \le \frac14 .
\]
C'est une condition \emph{nécessaire} ; elle est aussi suffisante si
$u_1 = u_2 = 0$ (on retrouve la question 2 avec $a^2 = \nu$).
\Etape{4} \emph{Condition supplémentaire due à l'advection.} Pour $\theta$ petit,
$\abs{\xi}^2 \approx 1 - 2d\abs{\theta}^2 + (c_1\theta_1 + c_2\theta_2)^2$ ; en
choisissant $\theta$ parallèle à $(c_1, c_2)$, il faut
$c_1^2 + c_2^2 \le 2d$. Le calcul numérique de $\max\abs{\xi}$ sur tous les
$(\theta_1, \theta_2)$ confirme que le schéma est stable si et seulement si
$d \le \frac14$ \emph{et} $c_1^2 + c_2^2 \le 2d$ (par exemple, avec $d = 0{,}1$ :
stable pour $c_1 = c_2 = 0{,}316$, instable pour $0{,}33$). Ce résultat est en
accord avec l'étude de Hindmarsh, Gresho et Griffiths [HGG].

\Question{Cas unidimensionnel.}
Avec $s = \sin^2(\theta/2)$ et $\sin^2\theta = 4s(1-s)$ :
\[
  \abs{\xi}^2 = (1 - 4ds)^2 + 4c^2s(1-s) = 1 + s\bigl[(4c^2 - 8d) + s(16d^2 - 4c^2)\bigr].
\]
Le crochet est une fonction affine de $s$ ; il est négatif ou nul sur $[0,1]$ si
et seulement s'il l'est en $s = 0$ et en $s = 1$, c'est-à-dire
\[
  c^2 \le 2d \qquad\text{et}\qquad d = \frac{\nu k}{h^2} \le \frac12 .
\]
La condition de diffusion devient $\nu k/h^2 \le \frac12$ (au lieu de
$\frac14$), et l'advection ajoute $(uk/h)^2 \le 2\nu k/h^2$.

\Verif $d = 0{,}5$ : stable pour $c = 1$, instable pour $c = 1{,}01$ ;
$d = 0{,}2$ : stable pour $c = 0{,}6$ ($c^2 = 0{,}36 \le 0{,}4$), instable pour
$c = 0{,}65$.

\Refs \BF{p.~746--747} ; [VN] ; [HGG].

% ------------------------------------------------------------------
\section{Exercice 16 : équation d'ondes amortie et non linéaire}

\Question{1)- Signification des termes.} Par analogie avec l'équation de la
corde vibrante du cours,
$\rho U_{tt} + kU_t = TU_{xx} + f$ :
\begin{itemize}
  \item $\partial^2u/\partial t^2$ : accélération (terme d'inertie) ;
  \item $a\,\partial u/\partial t$ : frottement visqueux, amortissement ($a > 0$) ;
  \item $c^2\,\partial^2u/\partial x^2$ : force de rappel élastique ; $c$ est la
        célérité des ondes ($c^2 = T/\rho$ pour une corde) ;
  \item $du^3$ : force non linéaire cubique (selon le signe de $d$, elle
        s'oppose au déplacement, $d < 0$, ou l'amplifie, $d > 0$).
\end{itemize}
Conditions : extrémités fixes, déplacement initial $f$, vitesse initiale $g$.

\Question{2)- Forme discrète (différences centrées).}
\Pourquoi Les différences centrées sont d'ordre 2 en temps et en espace, comme
pour (O2). Le terme d'amortissement centré fait apparaître $u_{i,j+1}$ des deux
côtés ; on le regroupe : le schéma reste explicite.
\Etapes
\Etape{1} Avec $x_i = ih$ ($h = l/m$), $t_j = jk$ :
\[
  \frac{u_{i,j+1} - 2u_{i,j} + u_{i,j-1}}{k^2} + a\frac{u_{i,j+1} - u_{i,j-1}}{2k}
  = c^2\frac{u_{i+1,j} - 2u_{i,j} + u_{i-1,j}}{h^2} + d\,u_{i,j}^3 .
\]
\Etape{2} On multiplie par $k^2$, on pose $\lambda = ck/h$ et $\beta = ak/2$ :
\[
  (1 + \beta)\,u_{i,j+1} = 2u_{i,j} - (1 - \beta)\,u_{i,j-1}
  + \lambda^2\bigl(u_{i+1,j} - 2u_{i,j} + u_{i-1,j}\bigr) + k^2d\,u_{i,j}^3 ,
\]
pour $i = 1, \dots, m-1$, $j \ge 1$.
\Etape{3} Conditions : $u_{0,j} = u_{m,j} = 0$ ; $u_{i,0} = f_i = f(ih)$.
\Etape{4} Stabilité de la partie linéaire ($d = 0$) : $\xi$ vérifie
$(1+\beta)\xi^2 - 2(1 - 2\lambda^2s)\xi + (1 - \beta) = 0$. Si
$\lambda \le 1$ (et $0 \le \beta < 1$), on vérifie que les racines sont de module
$\le 1$ (le produit des racines vaut $(1-\beta)/(1+\beta) < 1$ et, si elles sont
réelles, le polynôme est positif en $\xi = \pm1$). La condition est donc encore
$ck/h \le 1$ ; numériquement, $\max\abs{\xi} = 1$ pour $\lambda = 1$, $\beta = 0{,}1$,
mais $1{,}71$ pour $\lambda = 1{,}05$. Le terme $du^3$ n'est pas couvert par cette
analyse linéaire.

\Question{3)- Expression de $u_{i,2}$.}
\Pourquoi Le schéma demande deux niveaux pour démarrer. Le niveau $t = k$
s'obtient par Taylor, comme (O6) du cours, mais en utilisant l'équation
elle-même pour $u_{tt}(x,0)$ ; le niveau $t = 2k$ s'obtient ensuite par le
schéma avec $j = 1$.
\Etapes
\Etape{1} $u(x,k) = u(x,0) + k\,u_t(x,0) + \frac{k^2}2u_{tt}(x,0) + O(k^3)$, et
d'après l'équation, $u_{tt}(x,0) = c^2f''(x) + d\,f(x)^3 - a\,g(x)$. D'où
\[
  u_{l,1} = f_l + k\,g_l + \frac{k^2}{2}\Bigl[c^2\frac{f_{l+1} - 2f_l + f_{l-1}}{h^2} + d\,f_l^3 - a\,g_l\Bigr],
\]
qui utilise $f_{l-1}$, $f_l$, $f_{l+1}$ et $g_l$.
\Etape{2} Le schéma avec $j = 1$ :
\[
  u_{i,2} = \frac{2u_{i,1} - (1 - \beta)f_i + \lambda^2\bigl(u_{i+1,1} - 2u_{i,1} + u_{i-1,1}\bigr)
           + k^2d\,u_{i,1}^3}{1 + \beta} .
\]
Il fait intervenir $u_{i-1,1}$, $u_{i,1}$, $u_{i+1,1}$, donc (étape 1)
$f_{i-2}, \dots, f_{i+2}$ et $g_{i-1}, g_i, g_{i+1}$ : c'est l'expression
demandée.
\Etape{3} Forme entièrement développée dans le cas $d = 0$ (vérifiée
numériquement) :
\[
\begin{aligned}
  (1 + \beta)\,u_{i,2} ={}& \bigl(1 - 4\lambda^2 + 3\lambda^4 + \beta\bigr)f_i
    + 2\lambda^2(1 - \lambda^2)\bigl(f_{i+1} + f_{i-1}\bigr)
    + \frac{\lambda^4}{2}\bigl(f_{i+2} + f_{i-2}\bigr)\\
  &+ k(1 - \beta)\Bigl[2(1 - \lambda^2)g_i + \lambda^2\bigl(g_{i+1} + g_{i-1}\bigr)\Bigr].
\end{aligned}
\]
Contrôle : si $f \equiv F$ constante et $g = 0$, la somme des coefficients de $f$
vaut $1 + \beta$, d'où $u_{i,2} = F$, ce qui est correct loin des bords.

\Question{4)- Algorithme.}
\begin{algo}[ : onde amortie non linéaire]
  \Require $a$, $c$, $d$, $l$, $m$, $k$, $t_{\max}$, $f$, $g$
  \State $h \leftarrow l/m$ ; $\lambda \leftarrow ck/h$ ; $\beta \leftarrow ak/2$
  \If{$\lambda > 1$} \Print \og instable \fg{} \EndIf
  \For{$i$ allant de $0$ à $m$} \State $U^{\text{anc}}_i \leftarrow f(ih)$ \EndFor
  \State $U^{\text{anc}}_0 \leftarrow 0$ ; $U^{\text{anc}}_m \leftarrow 0$
  \For{$i$ allant de $1$ à $m-1$}
    \State $U^{\text{cour}}_i \leftarrow U^{\text{anc}}_i + kg_i + \frac{k^2}2\bigl[c^2(U^{\text{anc}}_{i+1} - 2U^{\text{anc}}_i + U^{\text{anc}}_{i-1})/h^2 + d(U^{\text{anc}}_i)^3 - ag_i\bigr]$
  \EndFor
  \State $U^{\text{cour}}_0 \leftarrow 0$ ; $U^{\text{cour}}_m \leftarrow 0$
  \For{$j$ allant de $1$ à $t_{\max}/k - 1$}
    \For{$i$ allant de $1$ à $m-1$}
      \State $U^{\text{nouv}}_i \leftarrow \bigl[2U^{\text{cour}}_i - (1-\beta)U^{\text{anc}}_i + \lambda^2(U^{\text{cour}}_{i+1} - 2U^{\text{cour}}_i + U^{\text{cour}}_{i-1}) + k^2d(U^{\text{cour}}_i)^3\bigr]/(1+\beta)$
    \EndFor
    \State $U^{\text{nouv}}_0 \leftarrow 0$ ; $U^{\text{nouv}}_m \leftarrow 0$
    \State $U^{\text{anc}} \leftarrow U^{\text{cour}}$ ; $U^{\text{cour}} \leftarrow U^{\text{nouv}}$
  \EndFor
  \Print $(ih, U^{\text{cour}}_i)$
\end{algo}

\Question{5)- Programme en C.}
\programme{C}{en C}{ex16_onde.c}

\Verif $l = 1$, $c = 1$, $f = \sin\pi x$, $g = 0$, $h = 0{,}02$, $k = 0{,}01$,
valeur en $x = 0{,}5$ et $t = 1{,}5$ :
\begin{itemize}
  \item $a = d = 0$ : $-5{,}8\times 10^{-4}$ (exact : $\sin(\pi/2)\cos(1{,}5\pi) = 0$) ;
  \item $a = 0{,}4$, $d = 0$ : $-0{,}054756$ ; la solution exacte
        $e^{-at/2}\sin\pi x\bigl[\cos\omega t + \frac{a}{2\omega}\sin\omega t\bigr]$,
        $\omega = \sqrt{\pi^2 - a^2/4}$, vaut $-0{,}054337$ ;
  \item $a = 0{,}4$, $d = -2$ (pas de solution exacte connue), avec $k = h/2$ :
        $0{,}093020$ ($m = 50$), $0{,}093271$ ($m = 100$), $0{,}093334$
        ($m = 200$) ; les écarts successifs ($2{,}5\times 10^{-4}$ puis
        $6{,}3\times 10^{-5}$) sont divisés par 4 : le schéma converge à
        l'ordre 2.
\end{itemize}

\Refs \BF{(12.19)--(12.25), p.~758--760 ; alg.~12.4, p.~761 ; p.~762}.

% ------------------------------------------------------------------
\section{Exercice 17 : pression dans la paroi d'un tube}

\Attention La position du rectangle intérieur n'est pas cotée ; on le suppose
centré, comme sur la figure (marges de 250~mm en $x$ et de 200~mm en $y$).

\Question{1)- Adimensionnement.}
\Pourquoi On veut des grandeurs sans dimension d'ordre 1 et une équation sans
paramètre. En prenant la \emph{même} échelle $L$ pour $x$ et $y$, l'équation de
Laplace garde sa forme ; avec des échelles différentes, un coefficient
$(L/l)^2$ apparaîtrait devant $\partial^2/\partial Y^2$.
\Etapes
\Etape{1} $X = x/L$, $Y = y/L$, $u = p/p_s$. Comme
$\partial^2p/\partial x^2 = (p_s/L^2)\,\partial^2u/\partial X^2$ :
\[
  \frac{\partial^2u}{\partial X^2} + \frac{\partial^2u}{\partial Y^2} = 0 .
\]
\Etape{2} Domaine : $0 \le X \le 1$, $0 \le Y \le 0{,}8$, privé du rectangle
$0{,}25 < X < 0{,}75$, $0{,}2 < Y < 0{,}6$.
\Etape{3} Conditions : $u = 0$ sur le bord extérieur, $u = 1$ sur le bord
intérieur.

\Question{2-a)- Équation discrète.}
\Pourquoi On remplace chaque dérivée seconde par la différence centrée
(exercice 8) ; c'est le cas $f = 0$ de (P2).
\Etape{1} Avec $X_i = ih$, $Y_j = jk$ :
$\frac{u_{i+1,j} - 2u_{i,j} + u_{i-1,j}}{h^2} + \frac{u_{i,j+1} - 2u_{i,j} + u_{i,j-1}}{k^2} = 0$.
\Etape{2} Pour $h = k$ : $u_{i,j} = \frac14\bigl(u_{i+1,j} + u_{i-1,j} + u_{i,j+1} + u_{i,j-1}\bigr)$
(la valeur au centre est la moyenne des quatre voisins).

\Question{2-b)- Conditions aux limites discrètes.}
\Pourquoi Pour que les deux bords tombent sur des nœuds de la grille, il faut
que $h$ divise $1$, $0{,}8$, $0{,}25$ et $0{,}2$ : on prend $h = k = 0{,}05$
(50~mm), soit $i = 0, \dots, 20$ et $j = 0, \dots, 16$.
\Etape{1} Bord extérieur : $u_{0,j} = u_{20,j} = 0$, $u_{i,0} = u_{i,16} = 0$.
\Etape{2} Bord intérieur : $u_{i,j} = 1$ pour $(i = 5$ ou $15$, $4 \le j \le 12)$
et pour $(j = 4$ ou $12$, $5 \le i \le 15)$.
\Etape{3} Inconnues : les nœuds avec $0 < i < 20$, $0 < j < 16$, hors du
rectangle $5 \le i \le 15$, $4 \le j \le 12$.

\Question{3-a)- Méthode de Gauss-Seidel.} On parcourt les nœuds inconnus et on
remplace chaque $u_{i,j}$ par la moyenne de ses voisins, en utilisant tout de
suite les valeurs déjà mises à jour ; on recommence jusqu'à ce que la variation
relative soit inférieure à $\varepsilon$ (exercices I.7 et VI.1).

\Question{3-b)- Algorithme.}
\begin{algo}[ : Gauss-Seidel sur le cadre]
  \State fixer $u$ sur les deux bords ; $u \leftarrow 0$ aux nœuds inconnus
  \Repeat
    \State $\mathit{num} \leftarrow 0$ ; $\mathit{den} \leftarrow 0$
    \For{$i$ allant de $1$ à $19$}
      \For{$j$ allant de $1$ à $15$}
        \If{$(i,j)$ n'est pas dans $[5,15]\times[4,12]$}
          \State $t \leftarrow (u_{i+1,j} + u_{i-1,j} + u_{i,j+1} + u_{i,j-1})/4$
          \State $\mathit{num} \leftarrow \mathit{num} + \abs{t - u_{i,j}}$ ; $\mathit{den} \leftarrow \mathit{den} + \abs{t}$ ; $u_{i,j} \leftarrow t$
        \EndIf
      \EndFor
    \EndFor
  \Until{$\mathit{num}/\mathit{den} < \varepsilon$}
  \Print $p_{i,j} = p_s\,u_{i,j}$
\end{algo}
Remarque : le domaine et les conditions sont symétriques par rapport aux droites
$X = 0{,}5$ et $Y = 0{,}4$ ; on pourrait ne calculer qu'un quart du domaine.

\Verif Avec $u = X^2 - Y^2$ imposé sur les bords (fonction harmonique, pour
laquelle le schéma est exact), l'algorithme retrouve $u$ à $4\times 10^{-12}$
près (124 itérations). Pour le problème posé ($\varepsilon = 10^{-10}$), il
converge en 102 itérations et toutes les valeurs restent entre 0 et 1, comme le
veut le principe du maximum discret du cours.

\Question{4-a)- Tube cylindrique.}
\Pourquoi Avec des bords circulaires concentriques, les coordonnées polaires
$(r, \varphi)$ épousent la géométrie, et le problème ne dépend pas de $\varphi$
(symétrie de révolution).
\Etape{1} Laplacien en coordonnées polaires :
$\frac{\partial^2p}{\partial r^2} + \frac1r\frac{\partial p}{\partial r}
 + \frac1{r^2}\frac{\partial^2p}{\partial\varphi^2} = 0$.
\Etape{2} Sans dépendance en $\varphi$ :
$\frac{\dd^2u}{\dd r^2} + \frac1r\frac{\dd u}{\dd r} = 0$,
$u(R_1) = 1$, $u(R_2) = 0$ (avec $\rho = r/R_2$, l'équation garde la même forme
sur $[\frac12, 1]$).
\Etape{3} Solution exacte (utile pour vérifier) :
$\frac{\dd}{\dd r}\bigl(r\frac{\dd u}{\dd r}\bigr) = 0$ donne
$u = A\ln r + B$, d'où $u = \dfrac{\ln(r/R_2)}{\ln(R_1/R_2)}$.

\Question{4-b)- Discrétisation.}
\Etape{1} $r_i = R_1 + i\Delta r$, $\Delta r = (R_2 - R_1)/N$ ; différences
centrées :
\[
  \Bigl(1 - \frac{\Delta r}{2r_i}\Bigr)u_{i-1} - 2u_i + \Bigl(1 + \frac{\Delta r}{2r_i}\Bigr)u_{i+1} = 0,
  \qquad u_0 = 1, \quad u_N = 0 .
\]
C'est l'exercice 8 avec $p = -1/r$, $q = g = 0$ : système tridiagonal,
résolu par Crout ou par Jacobi/Gauss-Seidel.
\Etape{2} Si le problème dépendait de $\varphi$, on utiliserait une grille
$(r_i, \varphi_j)$ avec $\varphi_j = j\Delta\varphi$ et la condition de périodicité
$u_{i,j+M} = u_{i,j}$ :
$\frac{u_{i+1,j} - 2u_{i,j} + u_{i-1,j}}{\Delta r^2}
 + \frac{u_{i+1,j} - u_{i-1,j}}{2r_i\Delta r}
 + \frac{u_{i,j+1} - 2u_{i,j} + u_{i,j-1}}{r_i^2\Delta\varphi^2} = 0$.

\Verif Avec $N = 40$, l'erreur maximale par rapport à la solution exacte est
$4{,}8\times 10^{-6}$.

\Refs \BF{§~12.1, alg.~12.1, p.~738--739 ; §~11.3}.

% ------------------------------------------------------------------
\section{Exercice 18 : équation de Poisson sur une plaque entaillée}

\Question{1)- Équation discrète.}
\Pourquoi Le domaine n'est pas un rectangle, mais il est formé de morceaux de
rectangles : si la grille est choisie pour que tous les bords (y compris ceux
de l'entaille) passent par des nœuds, on peut utiliser le schéma à cinq points
(P2) partout à l'intérieur, sans formule spéciale près des bords.
\Etapes
\Etape{1} Grille : $h = L/(3p)$ et $k = l/(2q)$ avec $p, q$ entiers ; alors les
abscisses $L/3$ et $2L/3$ correspondent à $i = p$ et $i = 2p$, et la hauteur
$l/2$ de l'entaille à $j = q$.
\Etape{2} En tout nœud intérieur, avec $r = h/k$ :
\[
  2(1 + r^2)u_{i,j} - (u_{i+1,j} + u_{i-1,j}) - r^2(u_{i,j+1} + u_{i,j-1}) = -h^2f_{i,j} .
\]
\Etape{3} Nœuds intérieurs : $0 < i < 3p$, $0 < j < 2q$, sauf ceux de
l'entaille et de ses bords ($p \le i \le 2p$ et $j \le q$). Tous les autres
nœuds sont sur la frontière et reçoivent $u_{i,j} = g(x_i, y_j)$.

\Question{2-a)- Méthode de Jacobi.} On calcule de nouvelles valeurs en tous les
nœuds intérieurs à partir des seules anciennes valeurs (deux grilles), puis on
remplace l'ancienne grille ; on recommence jusqu'au test d'arrêt.

\Question{2-b)- Algorithme.}
\begin{algo}[ : Jacobi sur la plaque entaillée]
  \Require $L$, $l$, $p$, $q$, $f$, $g$, $\varepsilon$, $q_{\max}$
  \State $h \leftarrow L/(3p)$ ; $k \leftarrow l/(2q)$ ; $r^2 \leftarrow (h/k)^2$
  \State \textbf{intérieur}$(i,j)$ : vrai si $0 < i < 3p$, $0 < j < 2q$ et non ($p \le i \le 2p$ et $j \le q$)
  \State $u_{i,j} \leftarrow g(ih, jk)$ si non intérieur$(i,j)$ ; $u_{i,j} \leftarrow 0$ sinon
  \For{$q'$ allant de $1$ à $q_{\max}$}
    \State $\mathit{num} \leftarrow 0$ ; $\mathit{den} \leftarrow 0$
    \For{tous les $(i,j)$ intérieurs}
      \State $v_{i,j} \leftarrow \bigl[u_{i+1,j} + u_{i-1,j} + r^2(u_{i,j+1} + u_{i,j-1}) - h^2f(ih, jk)\bigr]/\bigl(2(1+r^2)\bigr)$
      \State $\mathit{num} \leftarrow \mathit{num} + \abs{v_{i,j} - u_{i,j}}$ ; $\mathit{den} \leftarrow \mathit{den} + \abs{v_{i,j}}$
    \EndFor
    \State $u_{i,j} \leftarrow v_{i,j}$ pour tous les $(i,j)$ intérieurs
    \If{$\mathit{num}/\mathit{den} < \varepsilon$} \Print $u$ \Stop \EndIf
  \EndFor
\end{algo}

\Verif $L = 1$, $l = 0{,}8$, $p = 4$, $q = 3$, solution test $u = x^2 + y^2$
(donc $f = 4$ et $g = x^2 + y^2$) : Jacobi converge en 154 itérations et
l'erreur maximale est $6\times 10^{-12}$ (le schéma est exact pour ce polynôme).

\Refs \BF{§~12.1, p.~734--739}.

% ------------------------------------------------------------------
\section{Exercice 19 : équation de Laplace sur une plaque triangulaire}

\Attention On prend pour $n$ la normale \emph{extérieure} au domaine (convention
usuelle). Si $n$ est la normale intérieure, il suffit de changer $a$ en $-a$.

\Question{1)- Discrétisation.}
\Pourquoi Deux difficultés : l'hypoténuse est oblique, et les deux autres côtés
portent une condition de Neumann. Pour la première, on choisit des pas tels que
l'hypoténuse passe par des nœuds. Pour la seconde, on écrit l'équation aussi aux
nœuds des côtés de Neumann, en introduisant des points \og fictifs \fg{} hors
du domaine, éliminés grâce à la condition sur la dérivée.
\Etapes
\Etape{1} Sommets $(0,0)$, $(l,0)$, $(0,h)$ ; hypoténuse $x/l + y/h = 1$. Pas
$\Delta x = l/N$, $\Delta y = h/N$ : le nœud $(i,j)$ est sur l'hypoténuse si
$i + j = N$. Inconnues : $i, j \ge 0$, $i + j < N$.
\Etape{2} Hypoténuse : $u_{i,j} = 0$ si $i + j = N$.
\Etape{3} Nœuds intérieurs et nœuds des côtés, schéma à cinq points
($\rho = \Delta x/\Delta y$) :
\[
  u_{i,j} = \frac{u_{i+1,j} + u_{i-1,j} + \rho^2(u_{i,j+1} + u_{i,j-1})}{2(1 + \rho^2)} .
\]
\Etape{4} Côté $y = 0$ : la normale extérieure est $-\vec e_y$, donc
$-\partial u/\partial y = a$. Par différence centrée,
$(u_{i,1} - u_{i,-1})/(2\Delta y) = -a$, d'où le point fictif
$u_{i,-1} = u_{i,1} + 2a\Delta y$. De même sur $x = 0$ :
$u_{-1,j} = u_{1,j} + 2a\Delta x$. Au coin $(0,0)$, on utilise les deux.
\Etape{5} Lien avec le cours. En reportant $u_{i,-1}$ dans l'équation à cinq
points écrite en $(i,0)$, on obtient
$u_{i,1} - \bigl(1 + \frac{\Delta y^2}{\Delta x^2}\bigr)u_{i,0}
 + \frac{\Delta y^2}{2\Delta x^2}\bigl(u_{i+1,0} + u_{i-1,0}\bigr) = \Delta y\,\frac{\partial u}{\partial y}(x_i, 0)$,
c'est-à-dire exactement la formule (P4) du cours avec
$\partial u/\partial y = -a$ : les deux approches coïncident.

\begin{center}
\begin{tikzpicture}[scale=0.9]
  \foreach \i in {0,...,5} \foreach \j in {0,...,5} {
    \pgfmathtruncatemacro{\s}{\i+\j}
    \ifnum\s<5 \fill (\i,\j) circle (2pt); \fi
    \ifnum\s=5 \draw (\i,\j) circle (2.5pt); \fi
  }
  \foreach \j in {0,...,4} \draw (-1,\j) node {$\times$};
  \foreach \i in {0,...,4} \draw (\i,-1) node {$\times$};
  \draw[thick] (0,0) -- (5,0) -- (0,5) -- cycle;
  \node[right] at (5.4,4.2) {$\bullet$ inconnue};
  \node[right] at (5.4,3.5) {$\circ$ hypoténuse ($u = 0$)};
  \node[right] at (5.4,2.8) {$\times$ point fictif};
\end{tikzpicture}
\end{center}

\Question{2)- Algorithme.}
\begin{algo}[ : Gauss-Seidel sur le triangle]
  \Require $l$, $h$, $a$, $N$, $\varepsilon$
  \State $\Delta x \leftarrow l/N$ ; $\Delta y \leftarrow h/N$ ; $\rho^2 \leftarrow (\Delta x/\Delta y)^2$
  \State $u_{i,j} \leftarrow 0$ pour $i + j \le N$
  \Repeat
    \State $\mathit{num} \leftarrow 0$ ; $\mathit{den} \leftarrow 0$
    \For{$i$ allant de $0$ à $N-1$}
      \For{$j$ allant de $0$ à $N-1-i$}
        \State $u_O \leftarrow u_{i-1,j}$ si $i > 0$, sinon $u_{1,j} + 2a\Delta x$
        \State $u_S \leftarrow u_{i,j-1}$ si $j > 0$, sinon $u_{i,1} + 2a\Delta y$
        \State $t \leftarrow \bigl[u_{i+1,j} + u_O + \rho^2(u_{i,j+1} + u_S)\bigr]/\bigl(2(1+\rho^2)\bigr)$
        \State $\mathit{num} \leftarrow \mathit{num} + \abs{t - u_{i,j}}$ ; $\mathit{den} \leftarrow \mathit{den} + \abs{t}$ ; $u_{i,j} \leftarrow t$
      \EndFor
    \EndFor
  \Until{$\mathit{num}/\mathit{den} < \varepsilon$}
  \Print $u_{i,j}$
\end{algo}
(Les valeurs $u_{i+1,j}$ et $u_{i,j+1}$ utilisées au voisinage de l'hypoténuse
sont les zéros imposés.)

\Verif Pour $l = h$, la fonction $u = a(l - x - y)$ est harmonique, nulle sur
l'hypoténuse, et vérifie $-\partial u/\partial y = -\partial u/\partial x = a$ :
c'est la solution exacte. Avec $a = 2$, $l = h = 1$, $N = 10$, l'algorithme
la retrouve à $4\times 10^{-12}$ près (543 itérations), ce qui valide le
traitement des points fictifs. (Pour $l \ne h$, il n'existe pas de solution
affine, mais l'algorithme est le même.)

\Refs Cours, chapitre VI (conditions de Neumann, formule (P4)) ;
\BF{§~12.1, alg.~12.1, p.~738--739}.

\chapter*{Bibliographie}
\addcontentsline{toc}{chapter}{Bibliographie}

\begin{description}[style=nextline,leftmargin=2.2cm,labelwidth=2cm]
  \item[{[BF]}] R. L. Burden, J. D. Faires, A. M. Burden, \textit{Numerical
    Analysis}, 10\textsuperscript{e} édition, Cengage Learning, 2016. Les
    renvois donnent les numéros d'algorithmes, de théorèmes ou d'équations et
    les pages imprimées du livre.
  \item[{[Cours]}] P. Woafo, \textit{Méthodes numériques pour physiciens et
    ingénieurs}, Université de Yaoundé I (transcription LaTeX dans le dossier
    \fichier{latex/} de ce dépôt).
  \item[{[OEIS]}] The On-Line Encyclopedia of Integer Sequences, suite A094090
    (solution positive de $5(1 - e^{u}) + ue^{u} = 0$, constante de la loi de
    Wien), \url{https://oeis.org/A094090}.
  \item[{[VdW]}] Purdue University, \og Deviations from the Ideal Gas Law \fg{}
    (table des constantes de Van der Waals),
    \url{https://chemed.chem.purdue.edu/genchem/topicreview/bp/ch4/deviation.php}.
  \item[{[CSD]}] Équation de Carnahan-Starling pour les sphères dures :
    \og Playing with Marbles: Structural and Thermodynamic Properties of
    Hard-Sphere Systems \fg, \url{https://arxiv.org/abs/1310.5578} ;
    équation de Carnahan-Starling-De Santis :
    \url{https://iifiir.org/en/fridoc/a-modified-carnahan-starling-de-santis-equation-of-state-to-compute-106555}.
  \item[{[CL]}] Loi de Child-Langmuir : \og Child's Law \fg, SIMION,
    \url{https://simion.com/info/childs_law.html} ; \og A new approach to the
    Child-Langmuir law \fg, \url{https://arxiv.org/abs/1506.07417}.
  \item[{[LA]}] F. M. S. Lima, P. Arun, \og An accurate formula for the period of
    a simple pendulum oscillating beyond the small-angle regime \fg,
    \url{https://arxiv.org/abs/physics/0510206} (formule
    $T = \frac{2T_0}{\pi}K(\sin\frac{\theta_0}{2})$).
  \item[{[TW]}] L. N. Trefethen, J. A. C. Weideman, \og The exponentially
    convergent trapezoidal rule \fg, \textit{SIAM Review} 56(3), 2014,
    p.~385--458.
  \item[{[VN]}] \og Von Neumann stability analysis \fg, Wikipédia (anglais),
    \url{https://en.wikipedia.org/wiki/Von_Neumann_stability_analysis}
    (stabilité inconditionnelle du schéma de Dufort-Frankel).
  \item[{[DF2]}] \og New analysis of the Du Fort-Frankel methods \fg,
    \textit{Journal of Scientific Computing}, 2012,
    doi:10.1007/s10915-012-9627-2,
    \url{https://cris.iucc.ac.il/en/publications/new-analysis-of-the-du-fort-frankel-methods/}
    (erreur de troncature $O(k^2 + h^2 + k^2/h^2)$).
  \item[{[HGG]}] A. C. Hindmarsh, P. M. Gresho, D. F. Griffiths, \og The
    stability of explicit Euler time-integration for certain finite difference
    approximations of the multi-dimensional advection-diffusion equation \fg,
    \textit{International Journal for Numerical Methods in Fluids} 4(9), 1984,
    doi:10.1002/fld.1650040905.
\end{description}

\section*{Où trouver d'autres exercices corrigés}

Les exercices du cours reprennent des exercices classiques. Le livre [BF]
donne les réponses d'une sélection d'exercices à la fin de l'ouvrage
(\og Answers for Selected Exercises \fg, p.~787 et suivantes). En français,
l'ouvrage d'A. Quarteroni, F. Saleri et P. Gervasio, \textit{Calcul
scientifique : cours, exercices corrigés et illustrations en MATLAB et Octave}
(Springer), traite les mêmes méthodes avec des exercices corrigés.


\end{document}
